% A LaTeX template for MSc Thesis submissions to 
% Politecnico di Milano (PoliMi) - School of Industrial and Information Engineering
%
% S. Bonetti, A. Gruttadauria, G. Mescolini, A. Zingaro
% e-mail: template-tesi-ingind@polimi.it
%
% Last Revision: October 2021
%
% Copyright 2021 Politecnico di Milano, Italy. NC-BY

\documentclass{Configuration_Files/PoliMi3i_thesis}

%------------------------------------------------------------------------------
%	REQUIRED PACKAGES AND  CONFIGURATIONS
%------------------------------------------------------------------------------

% CONFIGURATIONS
\usepackage{parskip} % For paragraph layout
\usepackage{setspace} % For using single or double spacing
\usepackage{emptypage} % To insert empty pages
\usepackage{multicol} % To write in multiple columns (executive summary)
\setlength\columnsep{15pt} % Column separation in executive summary
\setlength\parindent{0pt} % Indentation
\raggedbottom  

% PACKAGES FOR TITLES
\usepackage{titlesec}
% \titlespacing{\section}{left spacing}{before spacing}{after spacing}
\titlespacing{\section}{0pt}{3.3ex}{2ex}
\titlespacing{\subsection}{0pt}{3.3ex}{1.65ex}
\titlespacing{\subsubsection}{0pt}{3.3ex}{1ex}
\usepackage{color}
\usepackage{datetime}

% PACKAGES FOR LANGUAGE AND FONT
\usepackage[english]{babel} % The document is in English  
\usepackage[utf8]{inputenc} % UTF8 encoding
\usepackage[T1]{fontenc} % Font encoding
\usepackage[11pt]{moresize} % Big fonts

% PACKAGES FOR IMAGES
\usepackage{graphicx}
\usepackage{transparent} % Enables transparent images
\usepackage{eso-pic} % For the background picture on the title page
\usepackage{subfig} % Numbered and caption subfigures using \subfloat.
\usepackage{tikz} % A package for high-quality hand-made figures.
\usetikzlibrary{}
\graphicspath{{./Images/}} % Directory of the images
\usepackage{caption} % Coloured captions
\usepackage{xcolor} % Coloured captions
\usepackage{amsthm,thmtools,xcolor} % Coloured "Theorem"
\usepackage{float}

% STANDARD MATH PACKAGES
\usepackage{amsmath}
\usepackage{amsthm}
\usepackage{amssymb}
\usepackage{amsfonts}
\usepackage{bm}
\usepackage[overload]{empheq} % For braced-style systems of equations.
\usepackage{fix-cm} % To override original LaTeX restrictions on sizes

% PACKAGES FOR TABLES
\usepackage{tabularx}
\usepackage{longtable} % Tables that can span several pages
\usepackage{colortbl}
\usepackage{blindtext}
\usepackage{booktabs}
\usepackage{makecell}
\usepackage{bbm}
\usepackage{multirow}

% PACKAGES FOR ALGORITHMS (PSEUDO-CODE)
\usepackage{algorithm}
\usepackage{algorithmic}

% PACKAGES FOR REFERENCES & BIBLIOGRAPHY
\usepackage[colorlinks=true,linkcolor=black,anchorcolor=black,citecolor=black,filecolor=black,menucolor=black,runcolor=black,urlcolor=black]{hyperref} % Adds clickable links at references
\usepackage{cleveref}


% Reference
\usepackage[square, numbers, sort&compress]{natbib} % Square brackets, citing references with numbers, citations sorted by appearance in the text and compressed
\bibliographystyle{unsrtnat} % You may use a different style adapted to your field


% PACKAGES FOR CODE
\usepackage{listingsutf8}


% OTHER PACKAGES
\usepackage{pdfpages} % To include a pdf file
\usepackage{afterpage}
\usepackage{lipsum} % DUMMY PACKAGE
\usepackage{fancyhdr} % For the headers
\usepackage{tcolorbox}
\usepackage[version=4]{mhchem} % Chemical Formula
\fancyhf{}

% Input of configuration file. Do not change config.tex file unless you really know what you are doing. 
\input{Configuration_Files/config}

%----------------------------------------------------------------------------
%	NEW COMMANDS DEFINED
%----------------------------------------------------------------------------

% EXAMPLES OF NEW COMMANDS
\newcommand{\bea}{\begin{eqnarray}} % Shortcut for equation arrays
\newcommand{\eea}{\end{eqnarray}}
\newcommand{\e}[1]{\times 10^{#1}}  % Powers of 10 notation


% Embedded Code
\lstset{inputencoding=utf8/latin1}
\definecolor{codegreen}{rgb}{0,0.6,0}
\definecolor{codegray}{rgb}{0.5,0.5,0.5}
\definecolor{codepurple}{rgb}{0.58,0,0.82}
\definecolor{backcolour}{rgb}{1,1,1}
\lstdefinestyle{mystyle}{
    backgroundcolor=\color{backcolour},   
    commentstyle=\color{codegreen},
    keywordstyle=\color{magenta},
    numberstyle=\tiny\color{codegray},
    stringstyle=\color{codepurple},
    basicstyle=\ttfamily\footnotesize,
    breakatwhitespace=false,         
    breaklines=true,                 
    captionpos=b,                    
    keepspaces=true,                 
    numbers=left,                    
    numbersep=5pt,                  
    showspaces=false,                
    showstringspaces=false,
    showtabs=false,                  
    tabsize=2
}
\lstset{style=mystyle}

%----------------------------------------------------------------------------
%	ADD YOUR PACKAGES (be careful of package interaction)
%----------------------------------------------------------------------------
\usepackage{siunitx}


%----------------------------------------------------------------------------
%	Watermark
%----------------------------------------------------------------------------
\usepackage{background}
\backgroundsetup{
    scale=8,
    angle = 45,
    color=black!20,
    position=current page.center,
    contents={}
}

%----------------------------------------------------------------------------
%	ADD YOUR DEFINITIONS AND COMMANDS (be careful of existing commands)
%----------------------------------------------------------------------------
\DeclarePairedDelimiter\ceil{\lceil}{\rceil}
%----------------------------------------------------------------------------
%	BEGIN OF YOUR DOCUMENT
%----------------------------------------------------------------------------

\begin{document}

\fancypagestyle{plain}{%
\fancyhf{} % Clear all header and footer fields
\fancyhead[RO,RE]{\thepage} %RO=right odd, RE=right even
\renewcommand{\headrulewidth}{0pt}
\renewcommand{\footrulewidth}{0pt}}

%----------------------------------------------------------------------------
%	TITLE PAGE
%----------------------------------------------------------------------------

\pagestyle{empty} % No page numbers
\frontmatter % Use roman page numbering style (i, ii, iii, iv...) for the preamble pages

\puttitle{
	title=Hot-Spot Detection \\[0.3cm] in Metal PBF Processes \\[0.5cm] \large A Systematic Literature Review of ML-Based Methods, % Title of the thesis
	name=Federico Cantarelli, % Author Name and Surname
	course=Management Engineering - Ingegneria Gestionale, % Study Programme (in Italian)
	ID  = 992964,  % Student ID number (numero di matricola)
	advisor={Prof. Bianca Maria Colosimo}, % Supervisor name
	coadvisor={Prof. Piercesare Secchi}, % Co-Supervisor name, remove this line if there is none
	academicyear={2022-23},  % Academic Year
} % These info will be put into your Title page 

\newdate{date}{05}{10}{2023}

\newpage
\thispagestyle{empty}
\begin{minipage}{0.6\textwidth}
    \vspace{18cm}
    \setstretch{1.5}
    \scriptsize Scritto con \LaTeX\ usando il templeta creato da S. Bonetti, \\ A. Gruttadauria, G. Mescolini e A. Zingaro. Disponibile su \href{https://www.overleaf.com/latex/templates/classical-format-thesis-scuola-di-ingegneria-industriale-e-dellinformazione-politecnico-di-milano/dkmvtndqkyxg}{Overleaf}. \\[4ex]
    Milano, \displaydate{date} \\ Prima e unica edizione. \\[4ex] 
    Quest’opera è stata rilasciata con licenza Creative Commons\\ \href{https://creativecommons.org/licenses/by-nc-nd/4.0/deed.it}{Attribuzione - Non commerciale - Non opere derivate 4.0 Internazionale}\\[3ex] 
    This work is licensed under a licence Creative Commons\\ \href{https://creativecommons.org/licenses/by-nc-nd/4.0/deed.en}{Attribution-NonCommercial-NoDerivatives 4.0 International}
\end{minipage}


\newpage
\thispagestyle{empty}
\begin{minipage}{\textwidth}
    \vspace{2cm}
    \begin{flushright}
    \textit{Ai miei genitori,\\[2ex]perché mi sono sempre stati d'esempio,\\[2ex]fin dal primo giorno.}
    \end{flushright}
\end{minipage}
\newpage
\blankpage


\thispagestyle{empty}
\begin{minipage}{\textwidth}
    \vspace{3cm}
    \begin{center}
    \textit{\large Executive Summary.}
    \end{center}
\end{minipage}
\newpage
\blankpage


\thispagestyle{empty}
\begin{minipage}{\textwidth}
    \vspace{3cm}
    \begin{center}
    \textit{\large Main Matter.}
    \end{center}
\end{minipage}
\cleardoublepage


%----------------------------------------------------------------------------
%	LIST OF CONTENTS/FIGURES/TABLES/SYMBOLS
%----------------------------------------------------------------------------

% TABLE OF CONTENTS
\startpreamble
\setcounter{page}{1} % Set page counter to 1
\thispagestyle{empty}
\tableofcontents % Table of contents 
\thispagestyle{empty}
\cleardoublepage

\thispagestyle{empty}
% LIST OF FIGURE
\listoffigures
\cleardoublepage

% LIST OF TABLE
\listoftables

\cleardoublepage


%----------------------------------------------------------------------------
%	PREAMBLE PAGES: ABSTRACT (inglese e italiano), EXECUTIVE SUMMARY
%----------------------------------------------------------------------------


% ACKNOWLEDGEMENTS
\chapter*{Ringraziamenti}
Voglio ringraziare la mia relatrice, la Prof.ssa Bianca Maria Colosimo, e il mio correlatore, il Prof. Piercesare Secchi, per avermi accettato e seguito durante il mio progetto di tesi, soprattutto durante le difficoltà degli ultimi mesi. Voglio anche ringraziare entrambi per le chiacchiere stimolanti durante i nostri meeting, e per avermi fatto appassionare all'analisi dei dati, alla statistica e al controllo dei processi. Voglio ringraziare Simone, con cui ho condiviso un pezzo importante di vita negli ultimi due anni, con cui ho riso, litigato, discusso e parlato durante le lunghi notti milanesi. Voglio ringraziare i miei amici Paolo e Valerio, per tutte le serate trascorse insieme a bere buon vino e a divertirci, in cui non pensavamo a nulla se non a vivere. Voglio ringraziare i miei compagni di viaggio, Paola, Giovanni, Edoardo, Matteo, con cui ho condiviso ansia, corse contro il tempo e successi. Voglio ringraziare Giorgia, che con un sorriso ha saputo tirarmi su il morale nelle giornate no. Voglio ringraziare Luca e Sasha, perché in questi anni, durante l'estate, mi hanno insegnato cosa vuol dire impegnarsi, appassionarsi al proprio lavoro, a risolvere i problemi con quello che si ha a disposizione e perché mi hanno fatto vedere quanto è bello sporcarsi le mani di grasso. Voglio ringraziare Cristiano, che con la sua esperienza e saggezza ha saputo consigliarmi sulle scelte da fare durante il progetto di tesi, tirandomi su nei momenti di sconforto (questo ragazzo farà strada). Voglio ringraziare i miei genitori, che con i loro sacrifici mi hanno dato la libertà più grande che ci possa essere: la libertà di poter scegliere. Voglio ringraziare i miei nonni, Silvana, Franca, Graziano e Adriano che mi hanno cresciuto da piccolo e che hanno contribuito a farmi diventare la persona che sono oggi. In particolare voglio ringraziare mia nonna Silvana, la mia madrina, perché so che sono sempre stato nelle sue preghiere. Grazie anche agli amici di sempre. Grazie ai momenti felici. Grazie a quello che sarà domani. Grazie alle emozioni, a tutte le emozioni.
\\[1.5cm]
\includegraphics[width=0.2 \textwidth]{Images/firma.eps}


% ABSTRACT IN ENGLISH
\chapter*{Abstract}
% ABSTRACT IN ENGLISH
\begin{flushright}
    \textbf{\textcolor{red}{Da completare.}} \hfill
\end{flushright}
\textbf{Keywords:} here, the keywords, of your thesis



% ABSTRACT IN ITALIAN
\chapter*{Sommario}
% ABSTRACT IN ITALIAN
La manifattura additiva (MA) di materiali metallici, è un approccio inovativo alla produzione che permette di ottenere componenti metallici attraverso la fusione selettiva di materiali metallici. Queste nuove tecnologie, ci permettono di produrre forme estremamente complesse che non sarebbe possibile ottenere con le tecniche manifatturiere tradizionali. La recente innovazione tecnologica, ha accelerato l'adozione di queste tecniche in ambito industriale, soprattutto per i processi di Electron Beam Melting e di Selective Laser Sintering. Questi processi utilizzano una sorgente esterna di energia concentrata per la fusione selettiva della polvere metallica, livello dopo livello. Questa tesi descrive in dettaglio i processi Powder Bed Fusion, e fornisce un quadro dei principi operazionali di base, le sfide, le criticità e i vantaggi rispetto ad altri processi manifatturieri, i difetti tipici che possono verificarsi, e i diversi approcci per individuarli. La tesi è incentrata sugli hot-spot, delle aree circoscritte caratterizzate da un'elevata temperatura sostenuta nel tempo e da un processo di raffreddamento anomalo, spesso relativi a parti circondate da polvere non fusa. Nella MA metallica, la temperatura assume un ruolo critico, in quanto la qualità dei componenti fabbricati è influenzata dal comportamento termico durante la fase di stampa. Infatti, gli hot-spot possono portare a delle modifiche della microstruttura del materiale solidificato, alla formazione di porosità e a difetti geometrici che minerebbero le qualità meccaniche del pezzo. Questa tesi presenta anche un algoritmo per la clusterizzazione delle aree del piatto di stampa caratterizzate da una storia termica simile utilizzando i dati funzionali termici. Infine. la tesi propone un'implementazione dell'algoritmo in Python.
\\[0.5cm]
\textbf{Parole chiave:} manifattura additiva, powder bed fusion, controllo qualità, hot-spot, machine learning, analisi di dati funzionali % Keywords (italian)

% 1814 caratteri 254 parole
\cleardoublepage


%-------------------------------------------------------------------------
%	THESIS MAIN TEXT
%-------------------------------------------------------------------------

\addtocontents{toc}{\vspace{2em}} % Add a gap in the Contents, for aesthetics
\mainmatter % Begin numeric (1,2,3...) page numbering

% --------------------------------------------------------------------------
% NUMBERED CHAPTERS % Regular chapters following
% --------------------------------------------------------------------------

% Introduction >>>

\chapter{Introduction}
\label{ch:Introduction}%
%\backgroundsetup{contents={DRAFT}}
% Introduction to additive manufacturing >>>
Additive manufacturing (AM) is described in ISO/ASTM 52900 \cite{organization_isoastm_2015} as the "process of joining materials to make parts from 3D model data, usually layer upon layer, as opposed to subtractive manufacturing and formative manufacturing methodologies".
In this first chapter, we will discuss and describe new AM processes or, popularly, 3D printing processes, which will most likely characterize the manufacturing scenario in upcoming years. Although the adjective "new" may not be quite appropriate, after all. Already from the 1950s - 1960s, before being called AM processes, these processes were called "Rapid Prototyping" (RP). From the name, it is easy to understand why this technology was firstly developed: there was the idea of a lean approach to new product development process, that would allow to contain costs and reduce time-to-market. It was only in 1984 and 1986 that the first 3D printing technologies were patented and the possibility of manufacturing a finished object by adding material layer by layer became a reality. The early 2000s were characterized by a push towards the use of these technologies also for large-scale production, trying to unlink RP from the research and development process alone. Indeed, in these years the first cost studies associated with 3D printing usage for mass production were conducted, among which it's worth mentioning \citeauthor{hopkinson_analysis_2003} (2003). 
\begin{figure}[H]
    \centering
    \includegraphics[width=0.55 \textwidth]{Images/costs.png}
    \caption[Traditional processes vs AM costs.]{Cost comparison for the manufacturing of a small plastic lever using different processes \cite{hopkinson_analysis_2003}.}
    \label{fig:costs}
\end{figure}
The early 2000s were characterized by a push towards the use of these technologies also for large-scale production, trying to unlink RP from the research and development process alone. Indeed, in these years the first cost studies associated with 3D printing usage for mass production were conducted, among which it's worth mentioning \citeauthor{hopkinson_analysis_2003} (2003). According to the author, for small production batches, AM technologies were more cost-effective and achieved a quality sufficient to ensure the products were saleable (Fig. \ref{fig:costs}). Furthermore, the primary cost factor in AM technologies was the initial machinery expenses. The same author believes these costs are set to decrease with mass adoption and technological advancements, leading to a reduction in overall production costs. The real turning point for widely spread 3D printing occurred in the early 2010s, when patents for some polymers used in fused deposition modeling systems (FDM) expired, and there was a real explosion of FDM 3D printers producers. We will discuss about FDM shortly. Nowadays, additive manufacturing is mainly used for the production of small batches of finished products, sometimes even single-unit batches of highly customized products, or in the prototyping phase of a new product. Moreover, AM processes allow for the creation of complex geometries and patterns that are not achievable with traditional manufacturing processes. As we will see in \emph{Chapter ~\ref{ch:Metal_AM}}, a typical example are lattice structures. 
I want to emphasize from the very beginning that these manufacturing processes do not intend to replace the traditional ones, but rather allow us to expand the limits of what we can build. It's important to understand that additive manufacturing allowed us to obtain components with mechanical functionalities and complexity that were unimaginable with traditional manufacturing processes.

%%%%%
%%%%%

% Additive Manufacturing Process Categories >>>
\section{Additive Manufacturing Process Categories} 
\label{sec:AMproc}
According to ISO/ASTM 52900 \cite{organization_isoastm_2015} (2015) and \citeauthor{gibson_additive_2015} (2015) there are seven different types of AM technologies: vat photo-polymerization processes, extrusion-based systems, material jetting, binder jetting, sheet lamination processes, powder bed fusion processes and direct energy deposition processes. Now we will briefly discuss about the different technologies and in \emph{Chapter~ \ref{ch:Metal_AM}} we will focus on AM of metals and powder bed fusion 3d printers.
\begin{itemize}
    \item \textbf{Vat photo-polymerization processes} (VP): VP processes use radiation-curable resins or photopolymers that can be solidified using controlled light sources, usually ultraviolet radiation. Typically, lasers are directed using scanning galvanometers, or multiple directional light sources such as the two-photon system or controlled LEDs. The polymers used in this printing technology are usually acrylate-based resins or epoxy resins, but also industrial ceramic materials, such as alumina, zirconia and silicone nitride. In recent years, new materials have been developed for the production of investment casting patterns \cite{3d_systems_investment_2023}, but also high technical performances material such as Cyanate Ester, which is used with Carbon Digital Light Synthesis\textsuperscript{TM} to produce finished pieces capable of withstanding high pressures (up to \SIrange[range-phrase = --]{3800}{4000}{\mega\pascal}) \cite{carbon_3d_carbon_2023}. The major advantages of this technology are the high level of accuracy of the finished products (up to \SI{0.05}{\milli\metre}), relatively high speed of manufacturing, and the variety of print sizes (ranging from very small printers to larger ones, up to \qtyproduct{1000 x 800 x 500}{\milli\metre}. The disadvantages are that the finished prints require significant post-processing, printers and manufacturing process are still quite expensive, and material choices are limited to photo-polymers only.
    \item \textbf{Fused filament fabrication} (FFF): FFF systems selectively extrude material, usually in the form of a filament, through a heated nozzle. They are also called fused deposition modeling (FDM). The FDM 3D printers we discussed earlier fall into this category and are characterized by a nozzle that is free to move in the horizontal plane and a system that moves up and down along the z-axis (usually a threaded tube with a diameter up to \SI{3}{\centi\metre}) allowing layer-wise material deposition. Typically, materials used are waxes, polyamide, acrylonitrile-butadiene-styrene, polyphenyl sulfone, polycarbonate, ceramics as well as biocompatible or biodegradable materials. The main advantages of FFF systems include their widespread popularity and economic accessibility, as well as the variety of materials that can be used, such as ABS, which provides excellent structural properties, or PLA (polylactic acid), which can be extruded at relatively low temperatures and it is compostable. The drawbacks of this technology include lower print speeds and a lower achievable minimum resolution if compared to other processes. So, usually we need to do some post-processing operations to achieve the desired finishing.
    \item \textbf{Material jetting systems} (MJ): MJ is a technique that allows 3D printing of a piece using tiny droplets of liquid material that are selectively deposited onto a surface. In a certain way, material jetting can be seen as the natural evolution of normal 2D inkjet printers but in three dimensions. Just like inkjet printers, in MJ there can be different "cartridges" of materials and colors. In fact, one of the main advantages of MJ is the possibility of printing high-resolution multi-colored or multi-material objects. Initially, the materials used with this technology were waxes, but in recent years the focus has been more on the deposition of acrylate photopolymers, which droplets of materials are solidified with an UV lamp directly attached to print nozzle. The main advantages of MJ are the almost absence of material waste, as the deposition-on-demand process is extremely precise (droplets of \numrange{25}{120} \unit{\micro\meter} at a rate of \numrange{0}{2000} \unit{drops / \second}), and the possibility of using different materials and colors simultaneously. The disadvantages of this solution are mostly the limited number of usable materials, since only polymers and waxes can be used and only materials that are not excessively viscous. Furthermore, the process and equipment required are relatively inexpensive and they can reach a remarkable printing speed.
    \item \textbf{Binder jetting systems}: Binder jetting, is very similar to MJ. In binder jetting, a liquid binding agent called "binder" is selectively deposited in small droplets to solidify powdered materials and allow different layers of powder to stick together. Due to the production method, the final objects may not always be suitable for withstanding large structural stresses and often require a long post-processing phase, which could increase variable costs. Binder jetting can be used with both thermo-plastic polymers, metals, and ceramic materials. The major advantages of this technology compared to MJ are the ability to use a much wider range of materials and the higher speed of printing. Moreover, it is possible to mix different powders to obtain final objects with different mechanical properties. The disadvantages, as mentioned earlier, are the fact that printed parts may not always be suitable to be used as structural parts and the costly post-processing operations necessary to improve the mechanical characteristics of the finished pieces.
    \item \textbf{Laminated object manufacturing} (LOM): In sheet lamination or laminated object manufacturing (LOM), a 3D printing technique, layers of material are bonded together to form a final object. So, each layer of the final object is made up of a sheet of laminate. There are two different approaches: "bond-then-form" and "form-then-bond". In the first case, the laminate is positioned and bonded to the substrate, and then cut following the model contour. In the second case, the laminate material is first cut and then placed on the substrate and bonded to the underlying layer. LOM techniques allow the usage of wood, thermo-plastic polymers, industrial ceramics, as well as paper, polyvinyl chloride, and composite materials. The major advantages of LOM are the ability to print very large parts quickly and the much lower fixed equipment costs compared to other types of AM processes.
    \item \textbf{Powder bed fusion} (PBF):  In these additive manufacturing techniques, a concentrated energy source is used to selectively melt layers of powder. Usually, the energy source is a laser (Selective Laser Sintering or SLS) or an electron beam (Electron Beam Melting or EBM) and the materials used can be polyamides, nylons, elastomers, and also metals such as stainless steel, titanium, aluminum, cobalt, and copper. It's important to note that this AM technique is widely used for direct manufacturing and allows for the production of finished objects with excellent mechanical properties and surface finishes.
    \item \textbf{Direct energy deposition} (DED): In DED, a concentrated heat energy source is used to selectively melt deposited material as soon as it is deposited. Depending on the energy source used, we can further divide the processes into laser metal deposition, electron beam metal deposition, and plasma metal deposition. Moreover, depending on the shape of the material used, we can distinguish between powder-fed and wire-fed systems. Finally, depending on the direction in which the material is fed, we can distinguish between off-axis feeding, in which the material is deposited by a side-mounted incident nozzle, and coaxial feeding systems, in which the powder or wire material is deposited co-axially with the energy beam. The major advantage of DED is that it can be used to add parts to existing components or repair broken components with a minimum need for support structures. There are several powders and printing area sizes available. DED drawbacks are that machining and post-processing are often required, and lower accuracy compared to PBF, which makes DED unsuitable to produce complex shapes.
    \end{itemize}
% <<< End of Additive Manufacturing Process Categories

%%%%%
%%%%%


\section{Aim and Structure of the Work}
\label{sec:aimwork}
Additive manufacturing processes for metallic parts, especially the so-called powder bed fusion processes, have been revolutionizing various industries in recent years, including aerospace, automotive, and energy industry. These processes are also used to create components that play a vital role during their use, often critical for safety. Consider, for instance, the importance of the brake system shown in the Fig. \ref{fig:bugatti}. Imagine the repercussions if the braking system failed on a Chiron that is speeding at \SI{350}{\kilo\metre /\hour}, or if a aerospace support, like the one in Fig. \ref{fig:supporto}, may break, causing a \SI{8500}{\metre} fall of a heavy piece. 
\begin{figure}
    \centering
    \subfloat[\label{fig:bugatti}]{
        \includegraphics[scale=0.4]{Images/bugatti.png}
    }
    \quad
    \subfloat[\label{fig:supporto}]{
        \includegraphics[scale=0.4]{Images/supporto.png}
    }
    \caption[Functional AM part printed in metal.]{On the left, a Bugatti brake caliper, one of the world’s largest functional parts produced in titanium alloy Ti6Al4V by AM. On the right, a support used in aerospace industry. From \citeauthor{du_plessis_beautiful_2019}.}
    \label{fig:funcpart}
\end{figure}
For these reasons, quality control, defect detection, and ensuring the production process adheres to standards are crucial now more than ever. Temperature, particularly in additive manufacturing for metallic parts, significantly influences the mechanical properties of the finished product as explained in \citeauthor{williams_situ_2019}. Indeed, unlike classic FDM 3D printing, the temperatures reached are much higher and the cooling process is more complex. Anyone who has somehow experienced FDM 3D printing would notice that the filament cools in a matter of seconds. On the other hand, with metallic PBF printing, one must wait considerably longer before final piece extraction from the production chamber. Moreover, there are many additional factors that could affect the temperature reached by the material during the PBF printing process and the consequent cooling process.
The aim of this thesis is to present a state-of-the-art review regarding the identification of so-called hot spots in PBF processes using machine learning algorithm. These are areas characterized by abnormal temperatures during the printing process, and the overheating leads to defects in the finished product.

This thesis is structured as follows. In \emph{Chapter~\ref{ch:Metal_AM}}, I will provide a comprehensive description of PBF processes for 3D printing of metal parts, the necessary technologies, materials, and some potential applications. In \emph{Chapter~\ref{ch:defects}}, I will outline the most common defects in PBF printing processes for metal parts, with a particular emphasis on the aforementioned hot spots, their causes and different monitoring systems used. \emph{Chapter~\ref{ch:state_ot_the_art}} will be a literature review to present the state-of-the-art in defect detection hot spots using machine learning algorithms. The research methodology will be meticulously detailed in Appendix~\ref{ap:research}. In \emph{Chapter~\ref{ch:baggingvoronoi}}, I will present the bagging voronoi classification algorithm as an innovative approach that might lead to a deeper understanding of the hot spots phenomenon in the printing processes. \emph{Chapter~\ref{ch:conclusions}} will conclude the thesis and propose some research streams that could be explored further. Finally, in Appendix~\ref{ap:Python}, I will present a simple Python \cite{python_software_foundation_python_2023} implementation of the algorithm proposed in \emph{Chapter~\ref{ch:baggingvoronoi}.}
\backgroundsetup{contents={}}
% <<< End of Introduction

%%%%%
%%%%%

% PBF Processes >>>
\chapter{Powder Bed Fusion Processes}
\label{ch:Metal_AM}
%\backgroundsetup{contents={DRAFT}}
\setlength{\tabcolsep}{10pt}
In this Section, I will provide a comprehensive overview of powder bed fusion (PBF) processes in metal additive manufacturing. In Section \ref{sec:pbf_proc}, we will discuss how 3d printers for metal PBF work, in Section \ref{sec:metalpowders} we will discuss about the powders used in those processes, while in Section \ref{sec:matterint}, we will delve into the physical process that allows the selective fusion of metal powders into the finished product. Lastly, in Section \ref{sec:examplesPBF} we will discuss some applications for metal PBF processes.

% Powder Bed Fusion Processes >>>
\section{Powder Bed Fusion Processes}\label{sec:pbf_proc}
As we have already discussed in Section \ref{sec:AMproc}, powder bed fusion processes can be further divided into laser powder bed fusion processes, also called selective laser sintering (SLS), or electron beam powder bed fusion processes also called electron beam melting (EBM). Despite the fact that the two processes are based on two different technologies and some differences in printers components, the two processes are based on the same fundamental idea: selectively melting metallic powders using external energy sources. All PBF 3D printers consist of two basic components: the powder bed, which is a container that can move along z-axis, and a highly concentrated energy source. 

% SLS >>>
\subsection{Selective Laser Sintering}
\label{subsec:LPBF}
\begin{figure}
    \centering
    \includegraphics[scale=1.2]{Images/PBF.jpg}
    \caption[Laser PBF in AM.]{Laser powder bed fusion process in additive manufacturing. Adapted from \cite{ozel_focus_2020}.}
    \label{fig:PBF}
\end{figure}
Now we will examine the SLS printing process from start to finish. Firstly, the plate mounting is calibrated and a vacuum atmosphere is created inside the printer chamber or a regular flow of inert gas is introduced to obtain controlled and uniform melting of the metal powder, in order to minimize oxidation and degradation of the powdered material. In industrial context, directing an inert gas is preferred to maintaining a vacuum atmosphere. This mainly for two reason: maintaining a vacuum atmosphere throughout the entire printing process require a lot of efforts and lowering the pressure in the process chamber can lower the boiling point of the metals, leading to the formation of an excessive amount of smoke. The mixture of vaporized metal and smoke can affect the effectiveness of the laser. Indeed, laser is an optical device and, as we all know, Micro particles suspended in the air could deflect the laser beam. Once the inert gas flow is established, the roller or the blade responsible for distributing the powder performs a rapid powder re-coating. Then an electric resistor located under the powder bed preheats the powder at a temperature of about \SI{500}{\degreeCelsius}. This operation helps to reduce the gap between the temperature of the chamber and the metal's melting temperature and also helps in maintaining an uniform temperatures within the build tray. This is necessary in order to prevent the warping of the part during the build due to non-uniform thermal expansion and contraction, which will lead to cracks and fractures (see Section \ref{sec:defects}). After these preliminary operations, the printing process can begin. A scanning system composed of a laser diode and a galvanometric mirrors system selectively melts the metal, alternating with the powder distribution blade, layer after layer. LASER is an acronym that stands for "light amplification by stimulated emission of radiation," and it is a monochromatic beam of photons characterized by low divergence and a focal spot size of \numrange[range-phrase = --]{30}{80}\unit{\micro\metre}. Laser equipment has the capability to generate powers in the range of thousands of Watt and can focus to beam spot sizes of fractions of a millimeter. These small spot sizes have the potential to create minuscule molten pools that can melt at extremely high travel speeds, reaching up to several meters per second. Nowadays, almost all lasers used in AM rely on active optical fiber sources. A schematic representation can be seen in Fig. \ref{fig:detailedfiber}. 
\begin{figure}
    \centering
    \includegraphics[scale=0.45]{Images/galvanometro.png}
    \caption[Galvanometric system.]{Dual-axis galvanometric system.}
    \label{fig:galvano}
\end{figure}
The energy transferred from the laser to the powder bed depends on the laser power (\numrange[range-phrase = --]{100}{1000}\unit{\watt}), the light-absorbing capacity of the material, and the scanning speed, which can be controlled (and it's limited) by the angular velocity of the galvanometer. The galvanometer system, Fig. \ref{fig:galvano}, system consists primarily of two mirrors rotating along their axis, enabling the laser beam to be directed to a desired point. The photon beam passes through a series of f-theta lenses that facilitate rapid laser movement and precise focusing of the laser spot. After completing the printing process, it is necessary to allow the object to cool down. Although supports are not always necessary as the powder gives support to the object itself, they may need to be inserted to ensure uniform heat transmission during both the printing and cooling processes. Once the object has cooled down, any excess powder can be removed. The excess powder can be recycled and reused after undergoing some preliminary processing and conformity controls \cite{strondl_characterization_2015}. Finally, if required, post-processing operations can be carried out, such as removing any supports, surface finishing, and annealing. The latter is required if the cooling phase inside the chamber has created an undesired micro structure in the printed piece.
% <<<SLS

%%%%%
%%%%%
% EBM >>>
\subsection{Electron Beam Melting}
\label{subsec:ebm}
In EBM printers, the process is fundamentally the same as in laser printers, with metal powder selectively fused in a powder bed, with a recoater blade to provide additional metal powder layer after layer. However, the energy source is not a laser (i.e., a beam of photons) but rather an beam of electrons. Secondly, the preheating stage is not accomplished through the use of an electric resistor, but instead through the use of the electron beam itself. The preheating operation is one of the biggest difference between EBM and SLS. To produce an \emph{electron beam}, an electric current is passed through a metallic filament, typically tungsten or tantalum. The resulting electron beam is then directed through a Wehnelt cylinder. When the cylinder is negatively or positively charged, blocks or allows the passage of electrons respectively.
\begin{figure}
    \centering
    \includegraphics[scale=0.5]{Images/EBM.png}
    \caption[Electron gun schema.]{Electron gun schema.}
    \label{fig:electrongun}
\end{figure}
A diagram of the electron gun beam's structure can be seen in Fig. \ref{fig:electrongun}. Finally, the electron beam passes through an anode charged with a voltage of up to \SI{60}{\kilo\volt}, which accelerates the electrons and allows the beam to be directed into the column through magnetic lenses. Unlike SLS printers which rely on optical lenses, the lenses in EBM printers are magnetic lenses. A more comprehensive description of how magnetic lenses function will be provided in Section \ref{sssec:magneticlens}. EBM printers require a vacuum atmosphere inside the printing chamber to correctly work and a minute amount of helium must be continuously injected into the chamber to prevent the accumulation of electric charges due to any residual electrons on the powder bed. The phenomenon will be discussed in Section \ref{subsec:ebminter}. To put the "minuscule amount" of helium required into perspective, consider that the pressure in the vacuum atmosphere inside the chamber is \SI{0.05}{\pascal}, and the pressure with helium flow is about \SI{0,2}{\pascal} (recall atmospheric pressure is \SI{1,01325e5}{\pascal}). The second major difference, as mentioned earlier, is that the preheating phase does not occur through an electric resistor but rather through an initial scan of the electron beam that happens at two distinct moments. During the first phase, the entire powder bed is scanned, and during the second phase, only sub-regions of the powder bed that will actually be printed are scanned. These preheating phases are useful in avoiding the so-called "smoke effect," an effect caused mainly by the electrostatic repulsion in the powder. Furthermore, the preheating phase is necessary for obtaining final objects with better performance and reduces the need for support structures during printing. This allows for the use of the entire volume of the powder bed container when printing, and even printing multiple objects one on top of the other, since there would be no printing supports interfering with them. This dual preheating phase causes the excess powder to solidify around printed pieces. Another structural difference of an EBM printer can bee seen in Fig. ~\ref{fig:ebm_printer}. Due to the vacuum inside the print chamber and the consequent lowering of the metal's vaporization temperature, a heat shield is required around the chamber to prevent the vaporized metal from condensing and creating a film on the entire printer structure. Essentially, the heat shield acts as a physical barrier that allows for easy removal and disposal of the metal film created during the printing process.
\begin{figure}
    \centering
    \includegraphics[scale=0.3]{Images/A-schematic-of-electron-beam-melting-EBM.png}
    \caption[EBM 3d-printer.]{Electron beam melting 3d-printer \cite{azam_-depth_2018}.}
    \label{fig:ebm_printer}
\end{figure}
\begin{table}
\small
    \centering 
    \begin{tabular}{|l l l|}
    \hline
    \rowcolor{bluepoli!40} % comment this line to remove the color
     & \textbf{EBM} & \textbf{SLS} \T\B \\
    \hline \hline
    \textbf{Energy source} & Electron beam & Laser  \T\B \\ 
    \textbf{Chamber atmosphere} & Vacuum(Almost) & Inert gas  \T\B\\
    \textbf{Scanning} & Magnetic lenses & Galvanometers \T\B \\
    \textbf{Energy absorption} & Conductivity-limited & Absorptivity-limited \T\B\\
    \textbf{Scan speeds} & Very fast & Limited by galvanometer inertia \T\B\\
    \textbf{Energy costs} & Moderate & High \T\B\\
    \textbf{Feature resolution} & \numrange{100}{200}\unit{\micro\metre} & \numrange{75}{100}\unit{\micro\metre}  \T\B\\
    \textbf{Materials} & Conductors & Polymers, metals and ceramic \T\B\\
    \textbf{Layer Thickness} & \SI{50}{\micro\metre} & \SIrange{10}{50}{\micro\metre}  \T\B\\
    \textbf{Min Wall Thickness} & \SI{0.6}{\milli\metre} & \SI{0.2}{\milli\metre} \T\B\\
    \textbf{Accuracy} & $\pm$\SI{0.4}{\milli\metre} & $\pm$\SI{0.3}{\milli\metre} \T\B\\
    \textbf{Build rate} & \numrange[range-phrase=--]{55}{80}\unit{\centi\metre^3 / \hour} & \numrange[range-phrase=--]{60}{100}\unit{\centi\metre^3 / \hour} \T\B\\
    \textbf{Powder Particle Size} & \numrange[range-phrase = --]{40}{105}\unit{\micro\meter} & \numrange[range-phrase = --]{15}{45}\unit{\micro\meter} \T\B\\
    \hline
    \end{tabular}
    \\[10pt]
    \caption{Comparison between SLS and EBM processes. Adapted from \cite{gallina_electron_2017}.}
    \label{table:slsvsebm}
\end{table}
There are also differences in the materials used for printing. In EBM printing, usually metal powder size is slightly larger than the powder used in laser printing, with a particle size of approximately \numrange[range-phrase = --]{40}{105}\unit{\micro\meter} against the \numrange[range-phrase = --]{15}{45}\unit{\micro\meter} in SLS processes. Furthermore, metal powder must be made of metals that exhibit good electrical conductivity. Indeed, in EBM also highly reflective metals can be used, because unlike photons, electrons are not repelled by mirrored surfaces since they penetrate the matters thanks to the material's own electrical conductivity.

To sum up what has been discussed so far, Table ~\ref{table:slsvsebm} highlights the main differences between the two processes described in previous sections.
% <<< End of EBM
% <<< End of Powder Bed Fusion Processes

%%%%%
%%%%%

% Metal Powders >>>
\vfill
\section{Metal Powders for AM} 
\label{sec:metalpowders}
As we have seen in the section \ref{sec:pbf_proc}, PBF processes require metal powders. Various metallic materials can be transformed into powders suitable for use in additive manufacturing processes, and can be applied in different sectors according to their mechanical properties. Table \ref{table:materialAMmetal} presents some examples of metallic materials currently available and the respective industrial sectors in which they are used. In the vast majority of cases, these powders are produced using atomization processes that exploit various physical methods to generate micro-particles, characterized by a more or less spherical shape and of which chemical purity depends on the method used. The inefficiency and cost of atomization processes used for manufacturing metal powders is the reason why, in the case of specific high-quality powders, the powder price can be up to 10 times the price of raw metal. \citeauthor{deng_origin_2020} (2020) have demonstrated that powders made of irregular micro-particles and with low chemical purity can lead to structural defects in the final parts, resulting in lower performance. Therefore, over the years, increasingly complex and expensive processes have been developed to obtain higher-quality powders.
\begin{table}[H]
\centering 
\small
    \begin{tabular}{|l l l|}
    \hline
    \rowcolor{bluepoli!40}
    \textbf{Material} & \textbf{Examples} & \textbf{Applications}\\
    \hline \hline
    \textbf{Stainless steel} & 316L, 174-PH, MS1, M300 & Food, biomedical, consumer \T\B\\
    \textbf{Ni-alloys} & In625, In718, In939 & Energy, motorsport\T\B\\
    \textbf{Al-alloys} & AlSi12, AlSi10Mg & Lightweight, aerospace, aviation\T\B\\
    \textbf{CoCr-alloys} & CoCrMo & Dental, biomedical\T\B\\
    \textbf{Ti-alloys} & Ti6Al4V, CP Ti & Biomedical, lightweight, aerospace\T\B\\
    \textbf{Tool steel} & Maraging 18Ni300 & Tooling, aerospace, automotive\T\B\\
    \textbf{Cu-alloys} & Bronze & Energy, heat exchanger\T\B\\
    \textbf{Precious} & Au, Pt, Ag & Jewellery, design\T\B\\
    \hline
    \end{tabular}
    \\[10pt]
    \caption{Material availability for metal AM.}
    \label{table:materialAMmetal}
\end{table}
\paragraph{Water atomization.}The process of water atomization of metals, Fig. \ref{fig:wateratom}, involves the production of small droplets of molten metal by exposing a stream of molten metal to a high-pressure jet of water. The water atomization process is typically carried out in a chamber where the molten metal is injected into a nozzle that directs the stream of liquid metal towards a high-pressure jet of water. As the molten metal comes into contact with the water, it rapidly cools and solidifies, forming small droplets that are collected at the bottom of the chamber. The size and shape of the metal droplets produced during the water atomization process can be controlled by adjusting the temperature of the molten metal, the pressure of the water jet, and the distance between the nozzle and the water jet. By controlling these parameters, it is possible to produce metal powders with a range of particle sizes between \numrange[range-phrase = --]{1}{500} \unit{\micro\metre}. The output-to-input ratio for this process is approximately 95\%, which means that from \SI{1}{\kilo\gram} of raw metal, it is possible to obtain up to \SI{950}{\gram} of powder on average. Compared to other powder production methods, the water atomization process offers several advantages including high production rates, the wide range of particle sizes obtainable, and the fact that metal ingots can be used as process input, which shortens the manufacturing chain of this production process. On the other hand, the process is characterized by a low chemical purity, irregular shape as can be seen in Fig. \ref{fig:waterpow}, and needs extensive post-processing operations to obtain a decent-quality product. Moreover, it cannot be used with reactive materials such as titanium and aluminum.
\paragraph{Gas atomization.} The process of gas atomization, Fig. \ref{fig:gasatom}, can be performed using a variety of gases, depending on processed metals. The most commonly used gases are inert gases such as nitrogen, argon, and helium, even if the latter is barely used in industrial applications due to its high cost. Inert gases are preferred because they do not react with the molten metal and do not introduce impurities into the metal powders. Metal ingots are molten and the flow is rapidly solidified by exposing it to a high-pressure stream of gas. Main advantages of this process are the high production rate, the wide range of obtainable particles, its applicability to reactive materials, and its ability to ensure good chemical purity. The main drawbacks, are mostly related to the porosity of the resulting powder and the formation of small satellite particles.
\begin{figure}
    \centering
    \subfloat[\label{fig:wateratom}]{
        \includegraphics[scale=0.6]{Images/wateratom.png}
    }
    \qquad
    \subfloat[\label{fig:gasatom}]{
        \includegraphics[scale=0.6]{Images/gasatom.png}
    }
    \qquad
     \subfloat[\label{fig:plasmaatom}]{
        \includegraphics[scale=0.4]{Images/plasma.png}
    }
    \qquad
    \subfloat[\label{fig:repatom}]{
        \includegraphics[scale=0.7]{Images/repatom.png}
    }
    \caption[Atomization processes]{Water atomization process (a), gas atomization process (b), plasma atomization process (c), rotating electrode process (d) \cite{material_technology_innovations_co_rotating_nodate, material_technology_innovations_co_water_nodate, material_technology_innovations_co_gas_nodate, inovar_communications_ltd_metal_2020}.}
    \label{fig:atom}
\end{figure}
\paragraph{Plasma atomization.} The process of plasma atomization, Fig. \ref{fig:wateratom}, involves the usage of high-temperature plasma torches to melt metal wire. The high-energy plasma arc is created by passing an electric current through a non-consumable tungsten electrode and an inhert gas jet to direct the welding arc into a focused area. As the molten metal droplets are expelled from the plasma jet, they rapidly solidify, sometimes also thanks to the water-cooled chamber, and break up into small, spherical powders, which are then collected in a container. The plasma atomization process offers several advantages over other powder production methods, including the ability to produce extremely spherical shape. Moreover, it can be used with reactive material. Plasma atomization can also produce powders with very high purity levels, as the high temperature of the plasma jet helps to eliminate impurities in the molten metal. However, it is more expensive than other powder production methods due to the need for specialized equipment and the high energy consumption required to generate the plasma arc. The process can also be challenging to control, as variations in parameters can affect the size, shape, and purity of the resulting powders. Powders produced using this method has a low size range of about \numrange[range-phrase = --]{1}{200}\unit{\micro\metre}.
\paragraph{Rotating electrode process.} In rotating electrode process, Fig. \ref{fig:repatom}, a consumable metal electrode is rotated at high speeds while it is melted by an electric arc, in a chamber with a continuous inert gas flow. As the molten metal is exposed to the high-velocity inert gas stream, it rapidly solidifies and breaks up into small droplets, which are then collected at the bottom of the atomization chamber. Due to its high cost, this production method is typically used when it is necessary to obtain a metal powder with a perfectly spherical shape and without any impurities, ad example gold. If we take a look at Fig. \ref{fig:reppow}, we can easily grasp the potential of this technology in obtaining perfect spherical micro-particles.
\begin{figure}
    \centering
    \subfloat[\label{fig:waterpow}]{
        \includegraphics[scale=0.22]{Images/waterpow.png}
    }
    \qquad
    \subfloat[\label{fig:gaspow}]{
        \includegraphics[scale=0.21]{Images/gaspowd.png}
    }
    \qquad
     \subfloat[\label{fig:plasmapow}]{
        \includegraphics[scale=0.22]{Images/plasmapowder.png}
    }
    \qquad
    \subfloat[\label{fig:reppow}]{
        \includegraphics[scale=0.21]{Images/erppowder.png}
    }
    
    \caption[Powders from atomization processes]{Micro-particles obtained by water atomization (a), gas atomization (b), plasma atomization (c) and rotating electode process (d) \cite{slotwinski_characterization_2014}.}
    \label{fig:powders}
\end{figure}
% <<< End of Materials

%%%%%
%%%%%

% Energy matter interactions >>>
\section{Energy - Matter Interactions}
\label{sec:matterint}
This section should be regarded as "optional", and it's up to the reader's discretion whether to engage with it or not. In this section, we will delve deeply into the physical processes underlying the technologies discussed in Section \ref{sec:pbf_proc}. As I have already said, in SLS there is the interaction between a photon and metal powder, while in EBM there is the interaction between electrons and metal powder. I will discuss first laser-metal interactions and then electrons-matter interaction.
\subsection{Laser - Matter Interaction}
\label{subsec:sintering}
The type and color of the laser are two extremely important features in SLS. Indeed, the energy transmitted to the material depends on energy absorption characteristics of the material itself that vary according to wavelength of the laser source. In order for the fusion process to be successful, metallic powder bed particles must receive enough energy to melt on the atomic level.
\begin{figure}
    \centering
    \subfloat[\label{fig:laserintera}]{
        \includegraphics[scale=0.4]{Images/laser interactions.png}
    }
    \qquad
    \subfloat[\label{fig:laserintensity}]{
        \includegraphics[scale=0.35]{Images/laserintensity.png}
    }
    
    \caption[Laser interactions and laser intensity.]{Schematic representation of physical interaction of absorption, reflection, transmission, refraction, and scattering between laser beam and material (a) \cite{katayama_fundamentals_2020}, and laser intensity schema (b).}
\end{figure}
For example, we can see that in Fig. \ref{fig:ondette} that absorptivity of copper is about 50\% and 58\% for green and blue lasers, respectively. In other words, green and blue lasers are more advantageous in processing copper sheets in terms of higher initial absorption than other laser typologies. As we will explore later, the electron beam transfers energy to the material through charge movement, while the laser beam operates as an optical device and hence it is more sensitive to the behaviors illustrated in Fig. \ref{fig:laserintera}. Indeed, the metal specimen is highly reflective and opaque. This is also the reason why the electron beam penetrates deeper into the metallic powder compared to the laser beam. 
\begin{figure}
    \centering
    \includegraphics[scale=0.4]{Images/ondette.png}
    \caption[Material absorptivity as function of wavelength.]{Absorptivity of incident light on several materials at room temperature as function of wavelength \cite{katayama_fundamentals_2020}.}
    \label{fig:ondette}
\end{figure}
The intensity of the laser emission varies as a function of the radius as 
\begin{equation}
    \label{eq:intensitylaser}
    I(r)=I_0\cdot e^{-2 \frac{r^2}{r_0^2}}
\end{equation}
and a schematic representation of density is shown in Fig. \ref{fig:laserintensity}.
The energy density of the laser beam is described by the equation
\begin{equation}
    \label{eq:energydensity}
    E_d = \frac{P}{h\cdot v \cdot d}
\end{equation}
where $E_d$ is the energy density (\unit{\joule/\milli\metre^3}), $P$ is the laser power (\unit{\watt}), $v$ is the laser scan speed (\unit{\milli\metre / \second}), $h$ is the hatch distance (\unit{\milli\metre}), $d$ is the layer thickness of the powder(\unit{\milli\metre}). For a more detailed description of heat transfer, please refer to next section. The intensity of the laser beam's energy plays a crucial role in SLS because if it is not correctly calibrated, the final piece might exhibit porosities that would compromise its mechanical features. Another critical parameter to consider during the SLS process is the duration of the laser pulse. As depicted in the Fig. \ref{fig:pulsipulsiocomepulsi} a prolonged laser pulse results in a broad heat zone, causing a shock wave that leads to an increased debris formation during the interaction that can interfere with new layers.
\begin{figure}
    \centering
    \includegraphics[scale=0.35]{Images/laserpulsing.png}
    \caption[Material absorptivity as function of wavelength.]{Schematic of laser interaction with materials under different pulse duration: (a) long pulse duration and (b) short pulse duration and respective holes fabricated on a steel foil by (c-d) \cite{lin_femtosecond_2021}.}
    \label{fig:pulsipulsiocomepulsi}
\end{figure}
For this reason, in recent years, there has been a shift towards using ultra-short pulse lasers. A schematic representation of one of this laser can be seen in Fig. \ref{fig:duripoco}. The laser in the figure is called "femtosecond laser" and is employed in manufacturing contexts such as precision processing or correction of semiconductors and liquid crystals, processing of transparent materials like glasses and sapphires, production of waveguide tubes for optical communication, drilling of engine components for cars and aircraft, among other applications \cite{katayama_fundamentals_2020}.
\begin{figure}
    \centering
    \includegraphics[scale=0.4]{Images/duripoco.png}
    \caption[Ultrashort pulse laser.]{Emission mechanism of ultrashort pulse laser \cite{katayama_fundamentals_2020}.}
    \label{fig:duripoco}
\end{figure}
\subsubsection{Heat transfer}
\label{sssec:heattransfer}
The heat and mass related events in SLS are influenced by both heat and mass movement. Laser heating is extremely rapid due to the fast scanning laser velocities. In a closed system, with respect to the first law of thermodynamics, the energy balance equation is formulated as follows \cite{bouabbou_understanding_2022}:
\begin{equation}
    \label{eq:tutteQ}
    Q_L = Q_C + Q_{CV} + Q_R,
\end{equation}
where $Q_L$, $Q_C$, $Q_{CV}$ and $Q_R$ respectively are the laser heat flux quantity, conduction, convection and radiation heat quantities. It is important to recall that only $Q_C$ will contribute to the melting process. Thus, some of the heat applied on the powder bed will be lost in a form of heat convection and radiation. To describe the heat conduction, we can use Fourier's law:
\begin{equation}
    \label{eq:nonsiecapitouncazzo}
    \frac{\delta}{\delta x}\left(k \frac{\delta T}{\delta x}\right)+\frac{\delta}{\delta y}\left(k \frac{\delta T}{\delta y}\right)+\frac{\delta}{\delta z}\left(k \frac{\delta T}{\delta z}\right)+\dot{q} =\rho C_p \frac{\delta T}{\delta t}
\end{equation}
where $T$ (\unit{\kelvin}) is the temperature, $k$ (\unit{\watt.\metre^{-1}.\kelvin^{-1}}) is the thermal conductivity, $\dot{q}$ (\unit{\watt.\metre^{-3}}) is the rate of which the heat is applied to the system, $r$ (\unit{\kilo\gram.\metre^{-3}}) is the material density, $C_p$ (\unit{\joule. \kilo\gram^{-1}.\kelvin^{-1}}) is the specific heat and $t$ (\unit{\second}) is the interaction time between laser beam and powder particles.

\subsubsection{Optical Fiber Laser}
\label{sssec:fiberlaser}
One of the most common types of lasers used in SLS is the fiber optic laser \cite{milewski_additive_2017}. A diode is a semiconductor component which basically acts like an unidirectional gate for electric current. It permits current to pass effortlessly in one direction but significantly hinders its flow in the opposite one. Diodes have a polarity, characterized by its anode (positive terminal) and cathode (negative terminal). Typically, diodes conduct current only when a positive voltage is supplied to the anode. In a laser diode (LD), an electric current is channeled directly through a dual hetero-junction structured semiconductor from the external side. When electrons recombine with positive holes at the active layer located between the n-type and p-type semiconductor zones, as shown in Fig. \ref{fig:np}, there is the emission of light beam.
\begin{figure}
    \centering
    \subfloat[\label{fig:np}]{
        \includegraphics[scale=0.35]{Images/np.png}}
    \qquad
    \subfloat[\label{fig:detailedfiber}]{
        \includegraphics[scale=0.35]{Images/detailedfiber.png}}
    \caption[Laser diode and fiber laser.]{Emission mechanism of laser diode (a) and a detailed fiber laser schema (b) \cite{katayama_fundamentals_2020}.}
    
\end{figure}
In fiber lasers, the beam is emitted by an LD pumping and a fiber of high purity \ce{SiO2} quartz glass doped with \ce{Yb^3+} rare earth element. The diameter of the fiber is about \numrange[range-phrase=--]{10}{20}\unit{\micro\metre}. In this laser technology, optical pump diodes are coupled to an active laser fiber that has a unique reflective coating and Bragg gratings, that reflect the laser light back-and-forth, along the length of the fiber, to create a coherent beam of light at the output of the laser \cite{milewski_additive_2017}. In Fig. \ref{fig:detailedfiber} we can see full reflection mirror and partial reflection mirror installed next to the fiber. This means also that the adjustment of the laser beam is not needed, which means easier handling of the system. It is also possible to combine multiple laser modules using a beam combiner for more powerful laser. We can use additional optical fibers to deliver and contain the light energy, providing a robust, flexible, and fully enclosed beam path for beam delivery \cite{milewski_additive_2017}. Fiber laser has many \textit{advantages} such as high beam quality, the fact they are small and lightweight, high intensity (which allows higher scan speed) and high efficiency (cost decreasing), long-distance delivery, and require basically no maintenance. 
\subsection{Electrons-Matter Interaction}
\label{subsec:ebminter}
According to \citeauthor{schultz_h_electron_1994} (1994), while the physical characterization of electron beams has been recognized since the 1960s, the microscopic processes are such that they remain without a comprehensive quantitative explanation also nowadays. However, a comprehensive explanation of electrons properties and their interaction with matter was given by \citeauthor{krumeich_properties_2015} (2015).
\subsubsection{What is an Electron?}
\label{sssec:electron}
An electron is a fundamental subatomic particle that carries a negative electric charge and it's one of the primary constituents of atoms. As explained in Section \ref{subsec:ebm} the electron beam is made of electrons generated through the thermionic effect with a metallic filament, typically tungsten or tantalum. At the cathode, an accelerating voltage $V$ is applied, causing electrons acceleration till velocity $v$. As explained in \citeauthor{krumeich_properties_2015} (2015), we can write the equation for kinetic energy associated to an electron as 
\begin{equation}
    \label{eq:energyelectron}
    E  = e\cdot V = \frac{1}{2}mv^2
\end{equation}
Given the fact that electron mass is \SI{9.109e-31}{\kilo\gram}, electron charge is \SI{-1.602e-19}{\coulomb} and that a common voltage for EBM application is \SI{60}{\kilo\volt}, from \ref{eq:energyelectron} we find that electron velocity is
\begin{equation}
\label{eq:velocityelectron}
v=\frac{\sqrt{2\cdot e\cdot V}}{m}
\end{equation}
which is approximately \num{1.453e8} \unit{m/s} (about half the speed of light). Recalling the dualism wave-particle of the electron, we can use De Broglie equation to compute the wavelength of the electron beam \cite{krumeich_properties_2015}:
\begin{equation}
\label{eq:wavelength}
    \lambda = \frac{h}{\sqrt{2\cdot m \cdot e \cdot V}}
\end{equation}
where $h=$\SI{6.62607015e-34}{\joule.\second} is Planck constant. In case of a voltage higher than \SI{300}{\kilo\volt}, it's necessary to include a term that accounts for the relativistic nature of mass in \ref{eq:wavelength}, as the electron's speed would approach the speed of light.
\subsubsection{Magnetic Lenses in EBM}
\label{sssec:magneticlens}
In EBM, electron beam is controlled by electromagnetic lenses. As the beam passes through these lenses, the astigmatic lenses adjust the shape of the electron beam spot, the focus coil changes the spot size, and the deflection coil moves the beam along the x and y axes.
\begin{figure}
    \centering
    \includegraphics[scale=0.9]{Images/Magnetic_lens.jpg}
    \caption[Magnetic lens.]{Forces schema when an electron beam pass through magnetic lens \cite{wikipedia_magnetic_2023}.}
    \label{fig:magneticlens}
\end{figure}
The force $\mathbf{F}$ which an electron of charge $-e$ experiences when traveling with a velocity $\mathbf{v}$ due to a magnetic field $\mathbf{B}$ is given by Lorentz's law:
\begin{equation}
    \mathbf{F} = -e\cdot \left( \mathbf{v} \times \mathbf{B}\right)
\end{equation}
We can decompose magnetic field into a radial component $\mathbf{B}_R$ and into an axial component $\mathbf{B}_L$. The initial direction of the electron is parallel to the axis, thus it is subjected to the radial component $\mathbf{B}_R$ only, which is responsible for an orthogonal force to $\mathbf{v}$. Hence, electrons start rotating. When they start rotating, they are subject to the axial field of the magnetic field too, which result in a reducing radius spiraling trajectory that converge the beam towards the target.
\subsubsection{Electrons Interaction}
\label{sssec:electroninteractions}
As explained in \citeauthor{korner_additive_2016} (2016), we can distinguish two interaction types between electrons and matters, defined on the basis of transferred energy quantity. 
\paragraph{Elastic interactions} The energy of the electron remains constant: when electrons delve into the surface, they engage as negatively charged entities with the metal's negative field and the positive charge from the nucleus's protons. Electrons can be rerouted on a different path without losing kinetic energy (known as elastic scattering). After several interactions, these deflected electrons might sometimes be ejected from the metal. This phenomenon is called referred to as backscattering. In Fig. \ref{fig:electronsscattering} there are a representation of all the possible interactions between an electron and an atom nucleus. This event occurs because of the Coulomb interaction, which happens when a negatively charged entity (electron) comes close to a positively charged one (nucleus). Coulomb force is described by the equation:
\begin{equation}
    \label{eq:coulomb}
    F=\frac{Q_1\cdot Q_2}{4\pi \varepsilon_0 r^2}
\end{equation}
where $r$ is the distance between the charges $Q_1$ and $Q_2$ and $\varepsilon_0=$ \SI{8.854e-12}{\farad / \metre} is the dielectric constant. 
\begin{figure}
    \centering
    \includegraphics[scale=0.4]{Images/electrons.png}
    \caption[Scattering and backscattering of an electron.]{Scattering and backscattering effect of an incident electron inside the electron cloud of an atom \cite{krumeich_properties_2015}.}
    \label{fig:electronsscattering}
\end{figure}
The backscatter phenomenon is determined by material characteristics and by the angle at which the beam hits the metal surface, see Fig. \ref{fig:backscattering}. It represents an energy deduction from the melting procedure. Indeed, if an electron is bounced outside the material, it cannot transfer energy to the material itself. So, for process efficiency, those interactions should be avoided. Backscattering effect is directly proportional to the atomic number $Z$ of the material of the powder, specifically as $Z^2$.
\begin{figure}
    \centering
    \includegraphics[scale=0.4]{Images/backscattering.png}
    \caption[Backscattering of an electron at different angles.]{Montecarlo simulation for an aluminum specimen hit by an electron beam at different angles, \ang{0}, \ang{45}, \ang{60}, \ang{75} respectively \cite{goldstein_scanning_2018}. In red, backscattered electrons.}
    \label{fig:backscattering}
\end{figure}
Having a fine and as spherical as possible powder is important to reduce backscattering effect caused by incidence angle. In EBM processes, the more spherical the particles are, the more we can assume that the electron beam will strike the particle surface perpendicularly. This can be clearly seen in Fig. \ref{fig:rotondette}.
\paragraph{Inelastic Interactions:} In inelastic interactions, the energy of the electrons diminishes. The energy imparted to the sample results in the heating of the metal and its correspondingly melting. These interactions causes the release of certain emissions, detailed below \cite{krumeich_properties_2015, goldstein_scanning_2018}:
\begin{itemize}
    \item \textbf{Inner-shell ionisation:} The energy is passed to an inner electron, which is adequate to eject it from the atom. This results in an unstable atom configuration, leading to the release of X-rays or of an Auger electron. \emph{X-rays} are produced when an inner electron is released into the vacuum, and an electron from the valence shell descends to occupy the void left by the electron from a deeper level. On the other hand, an \emph{Auger electron} is produced when the core electron, stimulated by the electron beam's energy, is expelled into the vacuum. The left space is occupied by another electron from a higher level, which then radiates energy to closer electron until it is released into the chamber. That's the reason why in EBM 3d printer is crucial to have a shield to protect operators from X-rays.
    \item \textbf{Braking radiation:} The electron that enters the atom slows down due to its interaction with the positively charged nucleus. The subsequent reduction in energy is released as X-rays. X-rays of lower energy are completely absorbed by the material, while X-rays of higher energy are released in the printing chamber.
    \item \textbf{Secondary electron:} Electrons within the valence band require minimal energy to surpass the potential barrier, and various electron-matter interaction mechanisms can provide this energy. This can cause the release of this electrons in the printing chamber. The flow of inhert gas in the chamber is used to remove these electrons from the powder bed, since they can interfere with printing process.
    \item \textbf{Cathodoluminescence:} when an electron from the valence band rises to the conduction band, a vacancy is created in the valence band. This vacancy is then occupied by another electron descending from the conduction band, which subsequently releases a photon. This phenomenon is responsible for the characteristic light produced when the electron beam hits the metal powder in EBM. 
    \item \textbf{Plasmon:} an electron moving through the collection of electrons in the valence band can cause a disturbance, leading to a collective vibration of the free electrons. This interaction is prevalent in metals but can also take place in any substance with free electrons. This phenomenon is responsible for the creation of plasma oscillation.
    \item \textbf{Phonos:} Phonons are understood as the collective oscillations of atoms within a crystal lattice, resulting from inelastic engagements with the electron beam. When an atom initiates this vibrational movement, it conveys this activity throughout the lattice, impacting a significant volume. Phonos are the only responsible for the melting of metal powders. Other products are undesirable, or even accountable for side effects during the printing process.
\end{itemize}
\begin{figure}
    \centering
    \subfloat[\label{fig:rotondette}]{
        \includegraphics[scale=0.7]{Images/EBMparticles.png}
    }
    \qquad
    \subfloat[\label{fig:unsaccodiroba}]{
        \includegraphics[scale=0.4]{Images/Screenshot 2023-08-21 at 11.42.27.png}
    }
    \caption[Laser interactions and laser intensity.]{Electron beam interaction with spherical powder particles (a) and Product from the interaction between an electron and a metal specimen (b) \cite{tushar_ramkrishna_mahale_electron_2009, krumeich_properties_2015}.}
\end{figure}
In Fig. \ref{fig:unsaccodiroba}, there is a schematic representation of the products resulting from the collision between an electron and a metal specimen.
% <<< End of Energy Matter Interactions

%%%%%
%%%%%

% Applications of PBF Metal Additive Manufacturing >>>
\section{Applications of PBF Metal Additive Manufacturing}
\label{sec:examplesPBF}
One of the primary advantages of PBF processes is the capability to produce finished objects unattainable with traditional manufacturing systems. In this section, we'll explore some examples of how this technology can completely disrupts production processes of metal manufacturing. The nozzle in Fig \ref{fig:nozzle} is one of the more famous successful and documented examples of additive manufacturing for metal. It was designed by General Electrics in 2015. With AM, GE was able to produce the new nozzle in a single piece, rather than 20 separated pieces welded together. The weight was reduced by 25\%, it was five times more durable and 30\% more cost-efficient \cite{amy_kover_transformation_2018}. GE Aviation has set the goal of 32,000 nozzles per year when in full production \cite{milewski_additive_2017}, opening this manufacturing process to large scale distribution. In 2015, NASA designed the first 3D-printed rocket nozzle made of copper. Within the combustion chamber, the propellant burns at temperatures exceeding \SI{3000}{\kelvin}. To prevent it from melting, hydrogen, at temperatures just above \SI{100}{\kelvin} flows through over 200 intricately designed cooling channels. Cooling inlets can be seen along the chamber's top rim. Thanks to this new rocket nozzle, costs were reduced by 50\%, and manufacturing times were reduced by ten times \cite{tracy_mcmahan_nasa_2015}, opening the path to a more affordable and lean space industry. An application of PBF processes in the automotive sector is illustrated in Fig. \ref{fig:piston} by Porsche. The company managed to print pistons by PBF for 911 GT2 RS engine. These new pistons have achieved a temperature on the piston's o-ring that is \SI{20}{\degreeCelsius} lower than usual. This lead to an additional \SI{23}{\kilo\watt} power available. However, before we see large-scale production, we'll have to wait until at least 2030 \cite{roberto_baldwin_porsches_2020} according to Porsche's press release. AM is also being employed in the medical field, for example, in dentistry. The traditional ways of fabricating items like crowns and dental implants are being transformed by custom-fitted designs. Advanced high-precision 3D printers such as EOS M 100 DMLS are used to craft devices from Cobalt Chrome SP2 alloy, a medical-grade approved material to be used in medical field \cite{milewski_additive_2017}. Products with a small batch size, high precision, and significant value, like these dental devices, are increasingly being manufactured using these new technologies. The advantages of AM include the swift creation of tailor-made items for immediate application, such as implants, or indirect uses like drill guides and fixtures, all based on the patient's own medical characteristics for a perfect fit. Moreover, the surface finish or porous structures offer huge advantages for bone growth. In Fig. \ref{fig:skul} there is an example of a titanium skull implant.
\begin{figure}
    \centering
    
    \subfloat[\label{fig:nozzle}]{
        \includegraphics[scale=0.20]{Images/nozzle.png}
    }
    \qquad
     \subfloat[\label{fig:nozzlerocket}]{
        \includegraphics[scale=0.24]{Images/nozzlerocket.png}
    }
    \\
    \subfloat[\label{fig:piston}]{
        \includegraphics[scale=0.21]{Images/piston.png}
    }
    \qquad
    \subfloat[\label{fig:skul}]{
        \includegraphics[scale=0.20]{Images/skull.png}
    }
    
    \caption[Examples of AM in metals.]{Examples of AM metal part: GE nozzle used in airplane turbines to mix fuel and air (a), a rocket nozzle in copper by NASA (b), Porsche engine pistons (c) and a titanium skull implant (d).}
    \label{fig:amexamples}
\end{figure}
\subsection{Lattice Structure and Cellular Material} \label{subsec:lattice}
A specific application of PBF processes is the printing of so-called lattice structures. According to ISO/ASTM 52900 \cite{international_standard_organization_isoastm_2015}, \emph{lattice structure} are "three dimensional geometric arrangement composed of connective links between vertices (points) creating a functional structure". Lattice structures are three-dimensional structures made up of different connected single elements called "cells" that can have different shapes, designed to meet the desired mechanical features of the final object. These objects are distinguished by empty spaces within the structure. These structures can only be obtained thanks to layer-wise AM technologies. In recent years, biomimicry technique has also spread in additive manufacturing. Biomimicry is a set of design techniques that allow us to take inspiration from nature to find solutions to engineering problems or to design functional components that can be used in engineering applications \cite{pathak_biomimicry_2019, du_plessis_beautiful_2019}. Nature, over the past 3.8 billions years, has been able to find the most efficient way to develop functional solutions to evolutionary problems characterized by immovable constraints, whether they are determined by the organism itself or by external factors. Moreover, during the evolutionary process of species, nature has been able to create nano, micro, and macro-structures that provide unique structural properties, adapting shape to function and using only what is truly necessary to achieve the evolutionary goal. Indeed, one of life's principles explained in \citeauthor{baumeister_biomimicry_2011} (2011) states \textit{"Life integrates and optimizes these strategies to create conditions conducive to life"}. Therefore, it not only concerns the possibility of creating new multi-functional structures with completely new characteristics, but also of doing so by optimizing materials and reducing waste. Engineers have always been fascinated by cellular materials, which is evident from the fact that the first reference to the idea that structure could influence the functional characteristics and behavior of a material was made by Robert Hooke in 1665 \cite{l_gibson_cellular_2010}. Only in recent years, thanks to the computational power of CAD softwares, it has been possible to experiment with the mechanical characteristics of 3D printed lattice structures. Cellular materials are used mainly for their mechanical characteristics.
\begin{figure}
    \centering
    \subfloat[\label{fig:beehive}]{
        \includegraphics[scale=0.15]{Images/honey-bees-337695_1280.jpg}
    }
    \quad
    \subfloat[\label{fig:beelattice}]{
        \includegraphics[scale=0.186]{Images/lattice1b.jpg}
    }
    \caption[Bio-inspiration design.]{An example of bio-inspiration design: from a beehive (a) to a lattice structure (b).}
    \label{fig:bioinsp}
\end{figure}
Just think that the aluminum cube in Fig. \ref{fig:beelattice} was printed using only \SI{3.9}{g} of material and is able to support a weight of \SI{408}{Kg}, which means it can withstand a stress 100,000 times its weight \cite{noauthor_3d_2014}. If we want to provide a more rigorous framework for classifying the applications of lattice structures in engineering, we can refer to the one proposed by \citeauthor{mcnulty_framework_2017}. According to this framework, cellular structures in nature exist primarily for three reasons:
\begin{itemize}
    \item \textbf{3D space-filling structures}, maybe this is the reason that most closely links the use of cellular structures in nature with AM, namely the need to confer strength to the structure while limiting its weight. Examples of this are the beehive, whose structure provides it with High specific stiffness under self-weight;
    \item \textbf{Surface structures}, enable surfaces to gain specific functional characteristics. For example, veining on the underside of the Amazon water lily leaf allows the leaf to gain excellent mechanical resistance, or the pomelo skin. with its open-cell structure, enables the fruit to resist impacts;
    \item \textbf{Cylindrical structures}: there are also examples of the use of cellular materials to confer particular characteristics to cylindrical structures (often hollow structure) that would otherwise be too fragile. For example hedgehog quills have a particular internal structure that confers ovalization and buckling resistance to the quill.
\end{itemize}
Lattice structures can be used for all these purposes, highlighting the wonderful connection between nature and engineering. Lattice structures are versatile and can be effectively applied in various engineering fields to address multiple challenges. Specifically, they are commonly used in four main areas according to \citeauthor{bhate_classification_2019} (2019): 
\begin{itemize}
    \item \textbf{Structural engineering} for vibration control, strain isolation, and weight reduction purposes like in Fig. \ref{fig:f1freno};
    \item \textbf{Thermal engineering} for applications such as heat exchangers in Fig. \ref{fig:heatexchanger}, flame arresters, or heat shields;
    \item \textbf{Fluid dynamics engineering} which employs lattice structures as catalyst carriers or packaging;
    \item \textbf{Biomedical engineering}, where lattice structure can be leveraged for bone integration in prosthetic and cell growth
\end{itemize} 
Due to their unique properties, lattice structures are an attractive choice for engineers and designers seeking innovative solutions to complex problems across a wide range of industries, especially in the automotive, aerospace, energy and commodities, and medical industries.
\begin{figure}
    \centering
    \subfloat[\label{fig:heatexchanger}]{
        \includegraphics[scale=0.25]{Images/heatexchanger.png}
    }
    \qquad
    \subfloat[\label{fig:f1freno}]{
        \includegraphics[scale=0.3]{Images/f1freno.png}
    }
    \caption[Lattice structure applications.]{Some lattice structure application: a complex heat exchanger (a) and brake pedal of F1 racing car (b) \cite{milewski_additive_2017, du_plessis_beautiful_2019}}.
\end{figure}
% <<< End of Applications of PBF Metal Additive Manufacturing
\backgroundsetup{contents={}}
% <<< end of PBF processes

%%%%%
%%%%%

% Defect in PBF Processes >>>
\chapter{Defects in PBF Processes}
\label{ch:defects}
%\backgroundsetup{contents={DRAFT}}
Like all industrial processes, Powder Bed Fusion (PBF) relies on parameters control to ensure proper functioning. When a process shows the expected behavior, it's called "in-control process". On the other hand, when the process deviates from the desired behavior, its output becomes unpredictable, and it's referred to as "out-of-control" (OOC). When a process is out of control, anomalies and defects can arise in the output. In Section \ref{sec:defects}, we will describe the most common defects in PBF processes and their causes. In Section \ref{sec:comelotrovo} we will discuss different approaches for defect detection, while in Section \ref{sec:sensoriniiniini} I will report available sensors used in anomaly detection.  Sections of this chapter are mainly based on \citeauthor{grasso_-process_2017} (2017), \citeauthor{grasso_process_2017} (2017), \citeauthor{mostafaei_defects_2022} (2022) and \citeauthor{wu_additively_2023} (2023). Finally, in Section \ref{sec:hotspot}, we will focus on defects caused by anomalies in the temperature profiles of the building bed, the so-called hot spots.

% Categories of Defect in PBF Processes >>>
\section{Defects Categories in PBF Processes}
\label{sec:defects}
Defects in PBF processes can be divided into five main classes: porosity defects, residual stresses and cracking, geometric defects and dimensional accuracy, balling, and surface defects.
\paragraph{Porosity.} Porosity is a significant parameter for many metal AM applications, as it strongly impacts fatigue resistance and crack propagation in the component \cite{edwards_electron_2013}. 
\begin{figure}
    \centering
    \subfloat[\label{fig:poris}]{
        \includegraphics[scale=0.4]{Images/pori e crack.png}
    }
    \qquad
    \subfloat[\label{fig:acircular}]{
        \includegraphics[scale=0.31]{Images/acircular.png}
    }
    \qquad
    \subfloat[\label{fig:delamination}]{
        \includegraphics[scale=0.4]{Images/delamination.png}
    }
    \qquad
    \subfloat[\label{fig:balling}]{
        \includegraphics[scale=0.4]{Images/balling.png}
    }
    \caption[Defect examples in PBF.]{Pores, micro-void and cracks in a specimen of FeCrCoMnNi (a) \cite{mostafaei_defects_2022}, an acircular pores on a specimen of Ti–6Al–4V (b) \cite{tammas-williams_xct_2015}, an example of delamination in SLS (c) \cite{sames_metallurgy_2016} and the balling effect in stainless steal powder (d) \cite{li_balling_2012}.}
\end{figure} 
Porosity is characterized by empty spaces within the mass of the fused material. These voids can appear within a layer, between adjacent layers, or on the surface of the printed piece. Voids most frequently occur within the layer and can vary in size, form, and distribution. The primary causes for these pores include incomplete fusion in the powder bed (lack of fusion or LOF), keyhole effects, and encapsulated gas. We can also differentiate between round and non-round pores. Some examples of round pores can be seen in Fig. \ref{fig:poris}. The elongated voids observed between layers are labeled as 'acicular pores'; their stretched form and larger size recognize them. For example, the acircular pore in Fig. \ref{fig:acircular} has a measure of \SI{20}{\micro\metre}. Pores can either be dispersed throughout the material or primarily situated between the inner patterned area and the outer boundary, known as under-skin pores. Pores can also appear on the exterior surface, typically termed 'surface porosity'.
\paragraph{Residual stresses, cracking, and delamination.} In material science, residual stresses are internal tensions that remain in the finished piece once the printing phase is complete. In PBF, these stresses can arise from two different causes: the thermal gradient mechanism and the cool-down phase of molten top layers \cite{mercelis_residual_2006}. Cracking phenomena occur as a consequence of stress relief when the tensile stress exceeds the ultimate tensile strength of the solid material. When the tension is suddenly released, a secondary crack can be generated from the main crack. Fig. \ref{fig:poris} shows an example of this kind of defect. On the other hand, delamination happens when cracks originate and propagate between adjacent layers (inter-layer cracking). Delamination occurs when the residual stresses exceed the binding linkage between the top layer and the previous one. Most of the time, delamination occurs from the partial disconnection of the part from the baseplate \cite{sames_metallurgy_2016}, as in Fig. \ref{fig:delamination}.
\begin{figure}
    \centering
    \subfloat[\label{fig:lackofsupport}]{
        \includegraphics[scale=0.45]{Images/lackofsupport.png}
    }
    \qquad
    \subfloat[\label{fig:oxide}]{
        \includegraphics[scale=0.34]{Images/oxide.png}
    }
    \caption[Examples of defect in PBF.]{Example of geometrical errors in SLS caused by a wrong support (a) \cite{grasso_-process_2017}, examples of two local oxide contaminations (b) \cite{casati_microstructure_2016}.}
\end{figure} 
\paragraph{Geometric defects and dimensional accuracy.} Dimensional anomalies in PBF can be categorized into (i) contraction and expansion effects, (ii) warping and curling, (iii) dross formation on bottom-facing surfaces, (iv) super-elevated edges, and (v) additional on-plane geometric distortions. Contraction is among the more common defects, though the opposite effect (i.e., components that are consistently bigger than expected) might arise in specific scenarios. Warping occurs from heat dispersion processes and thermal tensions. Similarly, curling is a specific bending effect caused by inconsistent thermal growth and component shrinkage, generally tied to disparate contraction rates between the top and underside of overhanging segments. A combination of shrinking and warping effects results in curved outlines for bottom-facing surfaces intended to be flat. Super-elevated edges represent another type of OOC geometrical distortion, featuring elevated ridges of solidified material. Elevated edges not only compromise the final component's integrity but might also cause defect propagation since they can interfere with the recoating blade. Other distortions affect critical features such as thin walls, overhang areas, and acute corners. In correspondence with these features, the melt pool is surrounded mainly by loose powder, which has a lower conductivity than the solid material. The diminished heat flux yields local overheating phenomena that may deteriorate the geometric accuracy. This overheated area is called a hot spot (HS), and we will discuss it in Section \ref{sec:hotspot}.
\paragraph{Balling.} Melt ball formation, a.k.a. balling, occurs when the molten material solidifies into spheres instead of solid layers. This phenomenon is caused by surface tension, which prevents the molten material from spreading over the previous layer. The result is a rough and bead-shaped surface that produces an irregular layer deposition with detrimental effects on the density and quality of the final product. \citeauthor{li_balling_2012}(2012) pointed out three main problems when balling happens: (i) increased surface roughness, (ii) a large number of pores between the discontinuous metallic balls, and (iii) protruding spheres may interfere with the recoating blade. The latter happens only in the presence of a very severe balling. Fig. \ref{fig:balling} shows the balling effect on a stainless steel powder.
\paragraph{Surface defects.} In PBF methodologies, as in most AM processes, the texture of the surface is influenced by the layer-by-layer manufacturing process and by the existence of surface impurities, inconsistencies, and voids.
\begin{figure}
    \centering
    \includegraphics[scale=0.4]{Images/surface.png}
    \caption[Surface roughness in PBF]{Good (left) and poor (right) surface textures measured via confocal microscopy on parts produced either in the presence of good and improper gas flow conditions in L-PBF \cite{ladewig_influence_2016}.}
    \label{fig:surface}
\end{figure}
Processing conditions determine the surface's texture, the perimeter scanning technique, and the dimensions of powder particles. The orientation of the surface in relation to building direction influences its texture. In particular, downward-facing and upward-facing surfaces are known to have considerably different texture properties. Even though many parts produced via PBF are refined by post-processing operations (like surface polish and thermal procedures), the texture of the surface is of functional importance since it impacts the part's fatigue resilience. Furthermore, for medical lattice structures, the texture of the surface can help the implant's assimilation during bone healing. Fig. \ref{fig:surface} displays instances of optimal (left) and poor (right) surface textures as observed through confocal microscopy on components made under optimal and improper gas flow situations in L-PBF \cite{ladewig_influence_2016}. Surface imperfections are further increased by the previously mentioned balling phenomenon, leading to less smooth surfaces.
\paragraph{Microstructural inhomogeneities and impurities.} PBF techniques employ localized intense heat inputs over brief durations of beam-material interaction, which influence the part's microstructural composition \cite{thijs_study_2010}. Microstructural discrepancies or non-stable microstructures can impair the mechanical and operational capabilities of the component. Microstructural disparities encompass (i) contaminants, (ii) grain dimension attributes, and (iii) crystallographic orientations \cite{dr_bree_m_sharratt_non-destructive_2015}. Material contaminants include inclusions, foreign material adulterations, and oxide layer formations. Fig. \ref{fig:oxide} illustrates a specimen with pronounced defects from oxide formations and conglobation. These imperfections are responsible for the specimen's premature failure and reduced strength \cite{casati_microstructure_2016}. Some cracks originate from a microstructure defect in Fig. \ref{fig:microstructuredefect}.
\begin{figure}
    \centering
    \includegraphics[width=0.8\textwidth]{Images/microstructure.png}
    \caption[Microstructure details.]{Typical microstructure details for AM part delamination and different types of AM cracking: solidification cracking (a); liquation cracking (b); strain aging cracking (c). \cite{mostafaei_defects_2022}.}
    \label{fig:microstructuredefect}
\end{figure}
\paragraph{Defect Causes} The causes of the defects just described can be divided into four main categories:
\begin{itemize}
    \item \textbf{Defects induced by feedstock material.} Various defects could impact the metal powder utilized in Powder Bed Fusion (PBF), such as irregular structures of the particles, absorbed gas, particles exhibiting satellites or fractures, and so forth. Problems related to dimensions and shapes influence the powder's apparent density and flow characteristics, thereby directly influencing the final part's quality and performance. Indeed, as we said in Sec \ref{sec:metalpowders}, we can set atomization parameters to get a specific particle size range, but we cannot control it very precisely, given the very nature of the process. In Fig. \ref{fig:particlesdistributions}, we can see different particle distributions for different processes. We can reduce these defects by measuring the particle size before printing begins. We can use more straightforward methods, such as sieving analysis, or more complex procedures, such as laser diffraction or microscopy. Alternatively, as the particle distribution influences the flowability of the powder, we can measure this using the Carney funnel. Moreover, metallic powders might be contaminated by moisture, organic substances, adsorbed gases, and the presence of oxides and nitride films on particle surfaces. These contaminations can increase across successive cycles of powder reuse, and they could lead to a degradation of the mechanical properties of the material.
\begin{figure}
    \centering
    \includegraphics[width=0.8\textwidth]{Images/particlesdistributions.png}
    \caption[Particles distribution.]{Particles size distributions for different powders used in different PBF processes}
    \label{fig:particlesdistributions}
\end{figure}
    \item \textbf{Equipment-induced defects.} This category of defects results from misuse or degradation of certain critical printer parts. A classic example is the optical apparatus of L-PBF printers or the electron gun of EB-PBF printers. Some defects are generated not only by excessive usage of these components but also by their incorrect calibration. The control of the atmosphere in the printing chamber can also give rise to some defects. The flow rate and the laminarity of the inert gas in L-PBF printers directly impact the final quality of the printed object, including its porosity and geometrical accuracy. In addition, the oxygen inside the chamber can also interfere with the laser beam. However, the powder recoating system is one of the most critical components for this defect category. Indeed, the linear motion of the recoating system causes defects that may have been generated in other printing areas to propagate all over the build area. This aspect is critical when carrying out batch production: one defective part could lead to all parts' defectiveness. In this sense, some automatic repair systems can be used in these industrial contexts. One example of such a system is Penelope. Penelope is a prototype realized by Politecnico di Milano that, through a high-resolution chamber, a high-speed chamber, and an MWIR chamber, can detect super-elevated edges that could be a source for defect propagation. Once the defects have been identified using computer vision techniques, a grinding cart passes over the defective layer, selectively removing the critical areas and preventing the recoating blade from spreading the defect \cite{colosimo_penelope_2020}.
    \item \textbf{Defects due to improper design or job preparation.} Even though software tools for 3D printing have become better at helping designers make intelligent decisions and create optimized designs using generative techniques, many of the choices about how to design and print things still come from the experience of the person operating the machines. Also, objects made with 3D printing can be complex, so sometimes people have to make trade-offs or try different things to get them right, which can lead to mistakes and parts that need to look or work better. One big reason for errors is the wrong support design. These supports are essential to stop finished pieces from warping because of thermal stress. Another possible defect source is improper part orientation within the build volume. 
    \item \textbf{Process setting–induced defects.} Printing parameters or scanning strategies often generate defects in PBF processes. Indeed, both process parameters and scanning strategies depend on the part's material and design. Since almost all PBF systems only allow to set printing parameters at the beginning of the process and cannot be changed from one part to the next, a trial and error approach may be necessary for particularly complex designs to obtain the proper parameter configuration. In general, it is almost impossible to achieve first-time-right production. We have extensively discussed how process parameters can affect defects in the previous paragraphs and will discuss this further in the next section.
\end{itemize}
% <<< End of Categories of Defect in PBF Processes

%%%%%
%%%%%

% >>> Defects Monitoring Methods
\section{Defects Monitoring Methods}
\label{sec:comelotrovo}
In recent years, thanks to advancements in both computational power available at lower cost and rising technology of sensors, quality control has evolved into an increasingly complex and sophisticated process. In AM (Additive Manufacturing), with the layer-wise production logic, we can collect vast amounts of data with remarkable granularity. This section outlines the various approaches to quality control, providing a terminology framework as established in \citeauthor{richard_leach_integrated_2020} (2020) and \citeauthor{grasso_-situ_2021} (2021).
\begin{figure}
    \centering
    \includegraphics[width=0.6\textwidth]{Images/dovevai.png}
    \caption[Measurement and monitoring techniques.]{Graphical representation of different terms associated with measurement and monitoring techniques \cite{richard_leach_integrated_2020}.}
    \label{fig:dovevai}
\end{figure}
\emph{In-process} methods involve measuring data during or between consecutive production phases within the same manufacturing chain. These measurements are synchronized with the different stages of the manufacturing process, thus allowing monitoring it synergistically. When in-process measurements are executed immediately before, after, or amidst manufacturing points, they're called \emph{"in-line"} measurements. These are measured on distinct measurement systems along the standard production line, where manufacturing is not occurring. Consequently, they fall under 'off-machine' measurements. \emph{off-machine} measurements are taken outside the machine where the manufacturing process occurs. On the other hand, if measurements taken during the process use sensors mounted on the production machine, they're labeled \emph{'on-machine'} measurements. Those that principally document data directly from the manufacturing site are called "in-situ" measurements. In AM processes, \emph{in-situ} predominantly denotes sensing and surveillance techniques aiming to capture details about process stability and product quality during manufacturing. Data from this last method are particularly significant as they are collected close to the final product, thus carrying much information. Conversely, when the metrics aren't "in-process", they're labeled as \emph{"off-line"} or \emph{"ex-situ"}. They're categorized as off-machine measurements since they're usually executed outside the core manufacturing setup, perhaps at a separate factory measurement station or a lab. Within the domain of in-situ measurements, there is a slight difference between 'in-situ measurement' and 'in-situ monitoring'. The former implies the capability of in-situ sensors to gauge aspects, aiming to understand the process or assess the product's quality attributes during production. The latter alludes to the real-time identification of discrepancies, abnormalities, and unregulated process conditions inducing product flaws. They often require establishing an alert system or classification algorithm. Furthermore, considering "in-situ" measurements, depending on the measured process signature, we can distinguish five different monitoring levels \cite{grasso_-situ_2021, grasso_process_2017}. These levels are referred to as "observable signatures" only, i.e. measurable quantity, while the "derived signatures", i.e. quantities estimated via process modeling, are included in the process modeling literature. Let's delve into measurement levels for observable signatures, summarized in Fig. \ref{fig:levels}.
\begin{figure}
    \centering
    \includegraphics[width=0.8\textwidth]{Images/Level Measurement.png}
    \caption[In-situ measurement levels.]{Four in-situ measurement levels applicable to PBF processes \cite{colosimo_-machine_2020}.}
    \label{fig:levels}
\end{figure}

\begin{itemize}
    \item \textbf{Level 0:} in this level, we use all the signals that can be measured using embedded sensors to detect anomalies and unstable process states. These sensors are usually employed by the printer control unit to control the machine state and the machine chamber environment under control. Indeed, the main advantage of Level 0 observation is the possibility of monitoring the process without needing external or additional sensors, especially in EBM printers, where many signals are often already available. However, these signals are often distant from the actual process dynamics in terms of “time distance”. Some examples of these signals are chamber pressure, ambient temperature, and inert gas flow;
    \item \textbf{Level 1:} this level involves measurements gathered once (or more than once) per layer of the entire build area. We are interested in both the homogeneity of the powder bed and the geometrical and dimensional features of the printed slice, and in general, all the anomalies that can be captured by using a general static 2-D view of the printing plane;
    \item \textbf{Level 2:} this level includes quantities that represent the process quality and stability while the beam scans the tracks within the print area. Usually, these quantities can be measured with a temporal resolution higher than the layer-wise resolution in the previous level, since we can rely on sensors with high sampling rates (some cameras can use up to 10,000 fps). It is also higher than the sampling frequency of embedded sensors;
    \item \textbf{Level 3:} this level process signatures are the ones with the highest granularity possible in PBF systems, i.e. the melt pool. The melt pool is known to be a primary feature of interest in any process that involves a beam-material interaction aimed at achieving a local fusion of the material since it highly impacts the quality of the part and stability of the process. Some examples of process signatures are the melt-pool size, its shape, and the temperature intensity/profile of the melt pool itself;
    \item \textbf{Level 4:} this level regards the capability of gathering information about phenomena occurring under the currently processed layer. These measurements can be obtained with ad-hoc prototype machine configurations that enable transverse x-ray imaging and ultrasound and acoustic emissions caused by the release of elastic energy and plastic deformation of solidified layers. The research stream from which this measurement level originated is relatively recent. Indeed, it is not included in Fig. \ref{fig:levels}.
\end{itemize}
% <<< End of Defect Monitoring Methods

%%%%%
%%%%%

% >>> Sensors for anomalies detection
\section{Sensors for Anomalies Detection}
\label{sec:sensoriniiniini}
Each level described in Section \ref{sec:comelotrovo} involves different sensor typologies. This section will describe the most widely used sensors in in-situ process monitoring and measurement in AM. 

In \emph{level ~0}, we can use embedded sensors integrated within the printer, the same sensors the printing control unit uses to adjust printing parameters. These process signatures are particularly suitable in EB-PBF applications since many embedded systems are already mounted on EBM printers. Embedded sensor signals in EB-PBF are also known as ‘log signals’. Publications proposed a more advanced in-process use of signals from such sensors only in EB-PBF applications as a possible source of information for in-situ anomaly detection as in \citeauthor{grasso_data_2018} (2018) and \citeauthor{steed_falcon_2017} (2017). The same publications pointed out that many of these EB-PBF log signals are correlated with process errors and variations in process conditions. \citeauthor{steed_falcon_2017} (2017) proposed a software tool for the visualization and analysis of significant multivariate time series log signals called Falcon. 
\begin{figure}[H]
    \centering
    \subfloat[\label{fig:offaxiallaser}]{
        \includegraphics[scale=0.35]{Images/off axial laser.png}
    }
    \qquad
    \subfloat[\label{fig:fig:offaxialebm}]{
        \includegraphics[scale=0.35]{Images/off axial EBM.png}
    }
    \caption[Off-axis sensing.]{Example of off-axis sensing architectures in SLS (a) and EBM (b) \cite{colosimo_-machine_2020}.}
\end{figure}
In \emph{level ~1}, i.e. measurement and characterization of layer properties, we can distinguish two further types of measurements. The first is conducted before laser scanning, which allows us to identify areas where the powder is not uniform, potential defects produced by the recoating system, or super-elevated edges. On the other hand, we can perform a measurement after the laser scanning, through which we can detect any contamination of the powder of the printed slice or any geometric or dimensional deviations from the nominal shape. Several types of sensors can be employed for anomaly detection at this level.
\begin{figure}
    \centering
    \subfloat[\label{fig:quadratoir}]{
        \includegraphics[scale=0.45]{Images/quadratoir.png}
    }
    \qquad
    \subfloat[\label{fig:cilindroir}]{
        \includegraphics[scale=0.45]{Images/cilindroir.png}
    }
    \caption[Optical tomography examples.]{Examples of optical tomography images for cubic samples with variation of energy density (a) and a cylinder produced under shielding gas flow variation (b)\cite{bamberg_-process_2016}.}
\end{figure}
\paragraph{Off-Axis Sensing.} In both SLS and EBM, off-axis sensing is suitable for gathering information. Conventional cameras operating in the visible range are usually employed in L-PBF (Fig. \ref{fig:offaxiallaser}), while the high temperatures in EBM recoating and the difficulty of installing additional sensors on EB-PBF machines (Fig. \ref{fig:fig:offaxialebm}), makes their usage more challenging. The off-axis optical camera within the visible range, however, has a significant drawback: numerous publications such as \citealt{kleszczynski_error_2012} (2012) and \citeauthor{foster_bk_optical_2015} (2015) have shown how appropriate illumination condition significantly impacts the camera's acquired result. Both the surface pattern of the printed slice and the powder bed, as well as the contrast between the background (powder) and foreground (printed slice) areas, are influenced by the intensity, nature, and relative angle of the illumination source. Hence, non-uniform illumination conditions may mask some anomalies and involve some post-acquisition operations that might be challenging. The relative angle between the camera and the illumination source is also relevant in this measurement, as shown in  \citeauthor{caltanissetta_characterization_2018} (2018). Moreover, we must correct the distorted perspective given by the camera's positioning through specific geometric adjustment operations before analysis. In addition to traditional cameras, those operating in the NIR/IR range can monitor level 1 process signatures. They are particularly suitable in EBM processes due to the high temperatures encountered during printing. Fig. \ref{fig:quadratoir} and Fig. \ref{fig:cilindroir} are two examples of acquired images. One of the most employed systems for this kind of measurement is "LayerQam", developed by Arcam (GE Additive) for integration in its EB-PBF machines. The system consists of an NIR camera that acquires an image of the layer after the melting phase and uses local pixel intensity variations to detect anomalies. Furthermore, bright spots are used as a proxy for possible volumetric flaws and material discontinuity defects.
\paragraph{Fringe Projection.} The off-axis imaging techniques mentioned in the preceding paragraph can generate a 2D reconstruction of the powder bed and the printed section. We can achieve a height map of the powder bed with fringe projection, a 3D representation of the powder. We need to merge the visualization layer with topographic examination to achieve this map. The standard configuration for such measurements consists of a camera (single-view configuration) and a projector. The graphic representation of the system can be seen in Fig. \ref{fig:fringesetup}, while in Fig. \ref{fig:fringeprojection}, there is an example of a height map.
\begin{figure}
    \centering
    \subfloat[\label{fig:fringesetup}]{
        \includegraphics[scale=0.35]{Images/fringesetup.png}
    }
    \qquad
    \subfloat[\label{fig:fringeprojection}]{
        \includegraphics[scale=0.35]{Images/fringeprojection.png}
    }
    \caption[Fringe projection.]{Setup for fringe projection measurement system (a) and reconstructed height map of the power bed (b) \cite{zhang_situ_2016}.}
\end{figure}
Other authors proposed and tested multi-view configurations with two or more cameras, suitable for higher resolution and accuracy, like in \citeauthor{kalms_new_2019} (2019).
\paragraph{Blade Mounted Sensor.} Another possible approach for level 1 process monitoring is to employ an optical sensor directly mounted on the recoating blade of PBF printers. This way, we can avoid the geometric distortion typical of off-axis systems. Moreover, these sensors can obtain a 2D representation of the print bed and reconstruct the height maps. In the former case, we can get a resolution up to \SI{4800}{DPI} over a length of \SI{210}{\milli\metre} with a spatial resolution of \SI{5.3}{\micro\metre / pixel} \cite{tan_phuc_high-resolution_2019}, while in the latter with a little lower resolution of \SI{20}{\micro\metre / pixel} \cite{barrett_micron-level_2018}.
\paragraph{Co-axial sensing.} L-PBF can leverage the back-scattered light from the melt pool and its adjacent areas to gauge specific metrics indicative of melt-pool dynamics. Scattered light passes through the optical path of the laser, passing the same lenses of the processing laser, and a semi-reflective mirror redirects the laser beam to the scanner while simultaneously allowing the light emitted from the melt pool to be detected by sensors. This emitted light is then divided and directed to each sensor. Various co-axial detection designs might encompass multiple pyrometers, bi-wavelength or multi-wavelength, and assorted imaging devices in the visible or infrared spectrum. Fig. \ref{fig:coaxiallaser} shows a typical co-axial sensing system. With this method, we can reach high sample rates of \SI{50}{\kilo\hertz}, but the observation area is typically limited to \qtyproduct{0.5 x 0.5}{\milli\metre}, making this method unsuitable to capture time evolution. Alterations in the melt pool's dimensions and luminosity/temperature can induce changes in the pyrometer's readings, facilitating the observation of melt pool consistency and any unusual signal shifts. This architecture was first developed by the Katholieke Universiteit Leuven and licensed by Concept Laser \cite{colosimo_-machine_2020, kruth_jean-pierre_feedback_2007, berumen_quality_2010}. This measurement approach has achieved outstanding results, with a lateral resolution of \SI{30}{\micro\metre} and a vertical resolution of \SI{7}{\micro\metre} \cite{fleming_tracking_2020}.
\begin{figure}
    \centering
    \subfloat[\label{fig:coaxiallaser}]{
        \includegraphics[scale =0.35]{Images/coaxial laser.png}
    }
    \qquad
    \subfloat[\label{fig:coaxialheatmap}]{
        \includegraphics[scale=0.5]{Images/coaxialmap.png}
    }
    \caption[Co-axial sensing and results.]{A graphical representation of a system for co-axial sensing in L-PBF \cite{colosimo_-machine_2020} (a) and a corresponding in-situ topography reconstruction (b) \cite{neef_low_2014}.}
\end{figure}
\paragraph{Electronic imaging.} In EBM, the process gives rise to by-products like secondary and backscattered electrons and x-rays. The electrons generated from the beam and metal powder interaction could be harnessed and used to produce an electronic representation of the layer. Thus, one of the undesired effects outlined in Section \ref{sssec:electroninteractions} can be leveraged to extract valuable insights about the current state of the process. Using only backscattered electrons, \citeauthor{wong_pilot_2019} (2019) were able to reach a spatial resolution of \SI{60}{\micro\metre / pixel} We can see an example in Fig. \ref{fig:electronic1}. Moreover, we can use an electron beam to scan the layers and metal surfaces, using the heat shield as an electron collector during the melting phase as done by \citeauthor{arnold_operando_2020} (2020), a real-time measurement. This approach lead to a resolution of \numrange[range-phrase=--]{50}{100}\unit{\micro\metre / pixel}. An image from the study can be seen in Fig. \ref{fig:electronic2}.
\begin{figure}
    \centering
    \subfloat[\label{fig:electronic1}]{
        \includegraphics[scale=0.44]{Images/electronic1.png}
    }
    \qquad
    \subfloat[\label{fig:electronic2}]{
        \includegraphics[scale=0.44]{Images/electronic2.png}
    }
    \caption[Co-axial sensing result.]{Examples of electronic images in EBM: images with different magnification factors (a) \cite{wong_pilot_2019} and images of squared printed areas generated via real-time back-scattered signal acquisition from different materials and with increasing hatch spacing \cite{arnold_operando_2020}.}
\end{figure}
For \emph{level ~2} signatures, we use sensors for in-process measurement of fast transient phenomena and high-speed emissions during laser or electron beam scanning. The main goal of this level is to understand underlying physical phenomena that lead to the formation of the defect. Most methods used at this level involve off-axis mounted sensors, mainly cameras in the visible range or thermal cameras with high temporal resolution. Indeed, this is needed to capture fast and transient phenomena. There are different research streams.
\paragraph{Measurement of Process Heatmap and Temperature Profiles.} We will address this paragraph with greater detail as these methods are used to detect the defects discussed in Section \ref{sec:hotspot}.
\begin{figure}
    \centering
    \includegraphics[width=0.8\textwidth]{Images/voxel.png}
    \caption[Example of voxel thermic reconstruction.]{(a) SWIR camera observation of SLM process, (b) time series data from location in center of part cross-section (location 1) and border scan region (location 2), (c) definition of thermal features in time series data with melt pool image, and (d) thermal feature maps showing generation of volumetric data and slicing \cite{lough_-situ_2019}.}
    \label{fig:voxel}
\end{figure}
The layerwise manufacturing paradigm allows the "seeing" of the thermal history of the process in a new dual-paradigm. In this new paradigm, time and space overlap over the z-axis. Indeed, the production of an object using additive manufacturing grows in the same direction over time and space. \citeauthor{williams_situ_2019} (2019) demonstrated that almost all major quality characteristics of the final part and its mechanical performance depend on the thermal history. Local and global variations of heating and cooling patterns may affect material solidification at volumetric, microstructural, and geometrical levels. We can identify two primary goals linked to this monitoring system. The first focuses on building a heatmap of the layer using data collected at slower rates (up to 50 fps), and the other targets recording rapid thermal fluctuations, with time resolutions ranging from \numrange{300}{10000} \unit{fps}. In terms of spatial resolution, we can use measurement setups with a limited field of view, which enable us to reach a resolution of \numrange[range-phrase=--]{8}{100}\unit{\micro\metre / pixel}, or with a field of view covering the entire build area, with a resolution typically above \SI{100}{\micro\metre / pixel}. In both L-PBF and EBM, we can use thermal cameras in either short, medium or long-wave IR. To minimize measurement errors caused by the laser's light emissions, we should use a short wave IR camera, with a bandwidth from \numrange{1.35}{1.6}\unit{\micro\metre} \cite{heigel_situ_2020}. \citeauthor{lough_-situ_2019} (2019) used this type of camera, and they produced a thermal map of a specimen printed with L-PBF. Then, using extracted features from the map, they achieved a voxel-based representation of the part associated with its local and global quality attributes. There is the specimen reconstruction in Fig. \ref{fig:voxel}. Other researchers have employed high temporal resolution IR video imaging in L-PBF to improve the ability to map temperature profiles spatially and temporally, capturing quick and transient events. Although they exhibit a high sensitivity to emissivity values when determining absolute temperatures, thermal cameras in the medium or long-wave IR spectrum can be adjusted over a broader temperature span than short-wave IR cameras. They maintain high sensitivity even at elevated temperatures, which makes them particularly suitable for in-situ thermal video imaging in EB-PBF applications, especially when we need a larger field of view to construct a thermal map. In EB-PBF, the methodologies for in-situ video imaging have to be tailored to cater to the unique aspects of the process. Both standard and thermal imaging devices must be shielded against X-ray emissions and metal deposition. We can use x-ray shielding glass and a rotating Kapton film to protect sensors, ensuring that metallic vapors don't settle on the glass. The Kapton film allows for approximately 79\% IR transmission, while a \SI{10}{\milli\metre} x-ray shielding window lets through around 1.08\% of IR, as noted by \citeauthor{ralph_b_dinwiddie_thermographic_2013} (2013). However, PBF processes experience rapid state changes, transitioning from powder to a liquid state and then to a solidified form. The fast transitioning, combined with ongoing alterations in surface attributes and vaporized material emissions, could affect the accurate estimation of temperature measurements. Luckily, most of the time, we are interested more in the thermal signature fluctuation over time rather than pinpointing an absolute temperature. Data analysis and monitoring algorithms can be employed directly on original signals in these cases. On the other hand, if we need a precise temperature value, we can use more advanced thermal cameras, even if they are bigger, costly, and generally call for alterations in machinery hardware and viewports for assimilation into working setups. Standard cameras are more cost-effective and straightforward to incorporate, and with the right equipment, they can offer high temporal precision. Even though they don't grant precise temperature assessments, variations in pixel brightness in the visible spectrum can be used as a proxy to measure thermal shifts to spot abnormalities and imperfections. A rising stream of research consists of identifying localized intense heat events, referred to as 'hot spots' (HS), using rapid video imaging in the visible domain. A hot spot is characterized by localized excessive layer heating due to unregulated thermal interactions with the nearby matter.
\paragraph{Process By-Products Measurement.} Due to the different nature of the interactions in the SLS and EBM processes, as already described in \ref{sec:matterint}, we need to distinguish between the measurement systems of the by-products of these two processes. Since we have already discussed the measurement of the by-products of EBM in level 1, in this paragraph, we will focus on L-PBF processes. It has been widely demonstrated that the partial vaporization of the metal, also called plume, and the spatter ejected together with it are correlated with detrimental effects on part quality. Indeed, intense emissions of these by-products could lead to the deflection and partial absorption of the laser beam, resulting in a change in the laser spot geometry or energy density. Fig. \ref{fig:plume} visually represents the phenomenon just described. We can use visual and IR video imaging techniques to obtain information about spatter and plume. In Fig. \ref{fig:colosimoplume}, we can see an example of an in-control plume generation phenomenon acquired using an IR camera. Recent research has introduced a high-speed stereo vision system to identify and trace individual spatters in the 3D space above the layer. Through this approach, it's possible to monitor the trajectory of spatters, allowing for the calculation of both their speed and their spatial-temporal behaviors. Tracking the spatter along its course can enhance the understanding of process by-products, offering further clarity regarding their formations and the effect of processing conditions. We can use a high-speed, high-energy X-ray video imaging system to follow the spatter's origination mechanism.
\begin{figure}
    \centering
    \subfloat[\label{fig:plume}]{
    \includegraphics[scale = 0.4]{Images/plume.png}}
    \qquad
    \subfloat[\label{fig:colosimoplume}]{
        \includegraphics[scale=0.4]{Images/plumecolosimo.png}
    }
    \caption[Spatters and plume in SLS]{Schematic representation of spatters and plume emission in SLS (a) \cite{grasso_-situ_2021}, and plume emissions captured with long wave IR video imaging (b) \cite{grasso_statistical_2019}.}
\end{figure}
\begin{figure}
    \centering
    \label{fig:barretspatter}
    \includegraphics[scale=0.3]{Images/splatter.png}
    \caption[3D spatter reconstruction.]{Reconstructed of 3D spatter localisation via high-speed stereo vision (b) \cite{barrett_statistical_2019}.}
\end{figure}
\paragraph{Measurement of air-borne acoustic emission.} In L-PBF processes, a research stream involves air-borne acoustic emission sensors. The measurement system captures air density variations during the laser scanning by placing the sensor near the melted area. To measure acoustic vibration, we can use a Bragg grating opto-acoustic sensor installed into the build chamber at about \SI{200}{\milli\metre} from the process zone, with a sampling frequency of \SI{1}{\mega\hertz}. In \citeauthor{wasmer_situ_2019} (2019), the sensor was placed so that the fiber's longitudinal axis was perpendicular to the acoustic wave to increase its sensitivity. The result of the study can be seen in Fig. \ref{fig:acustic}. We can appreciate how different levels of porosities correspond to varying levels of audio waves. We can also use more common sensors like in \citeauthor{ye_defect_2018} (2018). In the study, authors installed a microphone into the build chamber at an angle of \ang{30} above the build area, with a frequency response in the range \numrange[range-phrase=--]{0}{100}\unit{\kilo\hertz}. The resulting measurements and their information content in the temporal and frequency domain can be viewed as signatures of the laser-material interaction.
The melt pool carries much information since it is the phenomenon with the highest granularity possible in PBF. Monitoring the properties of the melt pool is a process involved in \emph{level 3}. These systems have primarily been explored in L-PBF since co-axial sensors can be easily employed. Indeed, co-axial sensing is prevalently employed at this level. Using off-axis video imaging, we have already discussed some level 2 methods in EBM to extract details at both the track and melt pool levels. Recently, there's been emerging research focused on analyzing melt pool data using machine learning methods. We can identify two primary measurement systems: methods that integrate over space and resolve spatial details.
\begin{figure}
    \centering
    \includegraphics[width = 0.8\textwidth]{Images/acustic.png}
    \caption[Example of AE of different porosities specimens.] {From left to right, typical light microscope cross-sectional images with different porosities, their corresponding acoustic emission signals with a \SI{160}{\milli\second} time span, and their corresponding wavelet spectrogram of the regions produced with different scanned velocity \SI{300}{\milli\metre / \second} (a), \SI{500}{\milli\metre / \second} (b) and \SI{800}{\milli\metre /\second} (c) \cite{wasmer_situ_2019}}
    \label{fig:acustic}
\end{figure}
\paragraph{Spatially integrated methods.} Real-time co-axial monitoring of melt pool attributes has been employed since 2010. At this level, process signatures are the size of the melt pool, size, shape, and the intensity and spectrum of emitted radiation. We can use spatially integrated pyrometry with single or multiple photodiodes to capture the power of radiation from the melt pool with impressive temporal accuracy, typically at a \SI{100}{\kilo\hertz} sampling rate. Significant parameters influencing acquired data quality include the spectral range of measurement and the sensor's field of view. Most sensors work within a range of \numrange{400}{1100} \unit{\nano\metre}, even if it's hard to find sensors for wavelengths beyond \SI{1000}{\nano\metre}. Today, co-axial photodiodes are equipped in the majority of industrial L-PBF setups. In a certain way, they are slowly becoming embedded sensors of level 0. \citeauthor{montazeri_-process_2018} (2018) showed that the chemical composition of the material could be determined via melt pool radiation measurement, finding a way to detect potential powder contaminations by analyzing spectral data. \citeauthor{lough_-situ_2020} (2020) achieved similar results using a co-axial optical emission spectroscope to capture the spectral data and infer the chemical composition of vaporized by-products. There are also sensors with larger fields of view, up to the entire build area, which enables measuring radiation emitted by surrounding hot spots. However, these sensors have a much lower resolution, not even comparable to co-axial sensors.
\paragraph{Spatially resolved methods.} Co-axial video imaging techniques can obtain More detailed insights into the melt pool's attributes and consistency over time. Commonly, a high-speed camera (\numrange{1000}{50000} \unit{fps}) paired with a narrow-band NIR filter is utilized to maximize the dynamic range and record the primary melt pool emissions. In some cases, it is also possible to use off-axis video imaging methods, using high-magnification optics combined with a limited field of view. In \citeauthor{lane_measurements_2020} (2020), the authors employed an off-axis high-speed camera paired with a mirror, facilitating an up-close examination of the melt pool without interfering with the laser beam. They achieved a spatial resolution \SI{3}{\micro\metre / pixel}. In the same paper, the authors also used a thermal camera for melt pool thermography and achieved a lower resolution of \SI{36}{\micro \metre / pixel}.
\begin{figure}
    \centering
    \subfloat[\label{fig:colorini}]{
        \includegraphics[scale=0.12]{Images/colorini.png}
    }
    \qquad
    \subfloat[\label{fig:spectrograrphy}]{
        \includegraphics[scale=0.3]{Images/spectrography.png}
    }
    \qquad
    \subfloat[\label{fig:poolthermal}]{
        \includegraphics[scale=0.2]{Images/meltpool.png}
    }
    \caption[Level 3 measurement methods.]{Example thermographic video frames aquired using an high speed camera (a) \cite{zhirnov_accurate_2020}, average optical emission spectra of 304L stainless steel (b) \cite{lough_-situ_2020}, and (c) radiance temperature thermal video frames from a thermal camera (left) and high speed camera (right) \cite{lane_measurements_2020}.}
\end{figure}
Another method involves dual-wavelength video imaging, i.e. the acquisition of video sequences at two distinct wavelengths (\SI{700}{\nano\metre} and \SI{950}{\nano\metre}) to enable a temperature estimate via two-colour thermography. This method gets around the problems of guessing the melt pool's brightness by comparing the light levels from the two different wavelengths, thinking they shine consistently at both \cite{williams_situ_2019}.
All the layers discussed so far have one common element: they are used for monitoring the last printed layer, either before, during, or after the melting phase. For \emph{level ~4}, the subject of the monitoring process is under the last layer. Indeed, as the layer is being printed, the material characteristics underneath are also modified due to partial remelting of adjacent layers and heat exchanges within the build volume. Since we cannot use optical devices to monitor what we cannot see, level 4 monitoring relies on other sensors. We can use high-speed, high-energy X-ray imaging systems to gain information about the energy penetration depth and pore formation, or we can use X-ray diffraction measurement to detect any strain and stress formation. Fig. \ref{fig:xray4} shows the result of an x-ray video imaging in L-PBF. Moreover, a new research stream about in-situ x-ray micro-tomography was born in recent years. In \citeauthor{lhuissier_situ_2020} (2020), during the in-situ microtomography scan, 1500 projections were acquired resulting in a scan time of \SI{45}{s}, gathering a spatial resolution of \SI{3.64}{\micro\metre / pixel}. Another technique regards the in-situ measurements of acoustic emissions (AEs), specifically of structure-borne. This technique also has direct application potential in industry. Structure-borne AEs are suitable for detecting sudden releases of elastic energy that propagate within the material, allowing us to identify crack formations, detachments of overhang areas from supports, or delamination phenomena. We can also use multiple sensors placed at different locations to triangulate the position of the epicenter of the energy release. A pioneer in the use of AEs in PBF processes is Rieder. In \citeauthor{rieder_-_2016} (2016), he proposed an ultrasonic monitoring device in L-PBF mounted on the underside of the baseplate focusing on the bottom plate interface echo and the back wall echo patterns as proxies of discontinuities in the specimen.
\begin{figure}
    \centering
    \includegraphics[width=0.65\textwidth]{Images/xray4.png}
    \caption[X-ray video frame.]{Example of an in-situ x-ray video frame with relative defects identified \cite{paulson_correlations_2020}.}
    \label{fig:xray4}
\end{figure}
% <<< End of Sensors for anomalies detection

%%%%%
%%%%%

% Temperature Anomalies and Hot Spots>>>
\section{Temperature Anomalies and Hot Spots}
\label{sec:hotspot}
As discussed in \citeauthor{williams_situ_2019} (2019), almost all major quality characteristics of the final part and its mechanical performance depend on the thermal history. A specific defect characterized by anomalous temperature profiles is known as hot spot (HS). The primary cause of these defects is the laser beam being repeatedly focused on thermally insulated regions. These are areas predominantly surrounded by powder and lacking any support structure to ensure the proper cooling down phase, as explained in Section \ref{sec:pbf_proc}, but also overhanging walls and sharp edges. This would result in a very localized overheating of the powder bed, leading to the following defects \cite{bugatti_towards_2022}:
\begin{itemize}
    \item \textbf{High surface roughness:} overheated areas can lead to melting unwanted areas of powder and partially melted powder particles attach to the surface. This will result into small melted powder on th surface of the final object, thus increasing its roughness;
    \item \textbf{Change in microstructure of the material:} normal melting zone are characterized by a high cooling rate that leads to finer grain formation, while overheating regions tend to develop a coarser microstructure due to the slower cooling transient;
    \item \textbf{Porosity formation:} if the region is already hot, new laser scans may lead to material vaporization, hence unstable keyhole formation which is often correlated with porosity defects.
\end{itemize} An example of these defects can be seen in Fig. \ref{fig:cane}.

\begin{figure}
    \centering
    \subfloat[\label{fig:dio}]{
        \includegraphics[scale=0.45]{Images/modehs.png}
    }
    \qquad
    \subfloat[\label{fig:cane}]{
        \includegraphics[scale=0.4]{Images/printedhs.png}
    }
    \caption[Examples of hot spot defects.]{Examples of a complex shape CAD model with the position of observed hot spot (a) and relative observed defects in printed part (b) \cite{grasso_-process_2017}.}
\end{figure}
% <<< End of Temperature Anomalies and Hot Spots
\backgroundsetup{contents={}}
% <<< end of Hot Spot Defect

%%%%%
%%%%%

% Hot Spot State of the Art >>>
\chapter{Hot Spot Detection: a State of the Art Research}
\label{ch:state_ot_the_art}
\backgroundsetup{contents={DRAFT}}
This chapter will focus on machine learning algorithms for hot spot detection. Since I want this chapter to be entirely understandable to as many readers as possible, I will provide a brief introductory overview to elucidate the fundamental aspects in Section \ref{sec:introperritardati}. For a detailed understanding of the research strategy employed to write this state-of-art review, please refer to Appendix \ref{ap:research}. 

\section{Warm Up}
\label{sec:introperritardati}
In this section, I will explain some fundamental concepts that are essential for fully understanding of next section and also \emph{Chapter ~\ref{ch:baggingvoronoi}}. If the reader is already familiar with these concepts, they can comfortably proceed to Section \ref{sec:hotspotstateart}. 

% Machine learning>>>
\subsection{Machine Learning}
\label{subsec:ml}
In recent years, the volume of data and the speed at which data are available have increased exponentially. Thus, there is a need for new tools we can leverage in order to find valuable and actionable insights from this huge amount of data. The set of activities involved in the analysis of these large sets of data, usually with the purpose of extracting useful knowledge to support decision-making, has been referred to as machine learning \cite{vercellis_business_2009}. \emph{Machine learning} (ML) is an approach that involves the development of algorithms that enable computers to learn from and find patterns in a large amount of data. Rather than being explicitly programmed to perform a task, a machine learning system uses inference to make predictions or decisions based on input data, trying to minimize somehow the output error or a cost function. Formally, \textit{“A computer program is said to learn from experience $E$ with respect to some class of task $T$ and a performance measure $P$, if its performance at tasks in $T$, as measured by $P$, improves because of experience $E$”} \cite{zaki_data_2020, pierluca_lanzi_data_2021}. Suppose we have the experience $E$, a set of observations of a specific phenomenon, encoded as a dataset
\begin{equation}
    \label{eq:dataset}
    \mathbf{E}=\{\mathbf{x}_1, \dots, \mathbf{x}_N\} 
\end{equation}
where $\mathbf{x}_i\in \mathbb{R}^p$ is a $p$ dimensional vector, containing the realization of the $p$ features of the dataset. We have three possible scenarios:
\begin{itemize}
    \item \textbf{Supervised learning:} given the records of the dataset and respective desired outputs $t_1, t_2, \dots, t_N$, also known as labels, the algorithm learns how to produce the correct output given a new set of input;
    \item \textbf{Unsupervised learning:} given $\mathbf{E}$, the algorithm is used to exploit regularities, usually used as input for supervised learning algorithm, for explaining the observations or for classification;
    \item \textbf{Reinforcement learning:} the algorithm producing actions $a_1, a_2, \dots, a_H$ which affect the environment, and receiving rewards $r_1, r_2, \dots, r_H$, learn how to act in order to maximize rewards in the long term.
\end{itemize}
To further highlight the difference between supervised and unsupervised learning, it's important to note that unsupervised learning is also defined as "learning without a teacher" in \citeauthor{tibshirani_elements_2008} (2008). So, the main difference between supervised and unsupervised learning is the input data: in the former case, data is labeled a priori, while in the latter case is not. 
% <<< end of Machine learning

%%%%%
%%%%%

% Clustering algorithms >>>
\subsection{Clustering Algorithms}
\label{sec:clustering}
\emph{Clustering} is an unsupervised learning technique. All clustering algorithms aim to organize a collection of entities $\mathbf{E}$ into subsets or "clusters" wherein elements within a cluster are as similar as possible. In contrast, data in different groups are as dissimilar as possible. So, clustering aims to minimize the within-clusters variance while maximizing the between-clusters variance. But how do we define if two observations in $\mathbf{E}$ are similar? We have to rely on the concept of distance between two observations. There are several different definitions of distances. In general, given a space $\mathbb{R}^p$ and a set of points on this space, a distance measure $d\left(\mathbf{x}_1,\mathbf{x}_2\right)$ is a function mapping two points $\mathbf{x}_1\in\mathbb{R}^p$ and $\mathbf{x}_2\in\mathbb{R}^p$ to a real number, and satisfies these axioms:
\begin{enumerate}
    \item $d\left(\mathbf{x}_1,\mathbf{x}_2\right) \geq 0$
    \item $d\left(\mathbf{x}_1,\mathbf{x}_2\right) = d\left(\mathbf{x}_2,\mathbf{x}_1\right)$
    \item $d\left(\mathbf{x}_1,\mathbf{x}_2\right)=0 \Leftrightarrow \mathbf{x}_1 = \mathbf{x}_2$
    \item $d\left(\mathbf{x}_1,\mathbf{x}_2\right) \leq d\left(\mathbf{x}_1,\mathbf{x}_3\right) + d\left(\mathbf{x}_1,\mathbf{x}_3\right)$
\end{enumerate}
where $\mathbf{x}_3 \in \mathbb{R}^p$ is a third point in the space. Recalling that $\mathbf{x}_1=\left(x_{11}, x_{12}, \dots, x_{1p} \right)$ and $\mathbf{x}_2=\left(x_{21}, x_{22}, \dots, x_{2p} \right)$, most used distance measure:
\begin{itemize}
    \item Euclidean distance
\begin{equation}
    \label{eq:euclidean}
    d\left(\mathbf{x}_1,\mathbf{x}_2\right) = \sqrt{\sum_{i=1}^p\left(x_{1i}-x_{2i} \right)^2}
\end{equation}
\item Manhattan distance
\begin{equation}
    \label{eq:mandistance}
    d\left(\mathbf{x}_1,\mathbf{x}_2\right) = \sum_{i=1}^p\left|x_{1i}-x_{2i} \right|
\end{equation}
\item Mahalanobis distance
\begin{equation}
    \label{eq:mahadistance}
    d(\mathbf{x}_1, \mathbf{x}_2) = \sqrt{(\mathbf{x}_1 - \mathbf{x}_2)^\mathrm{T} \Sigma^{-1} (\mathbf{x}_1 - \mathbf{x}_2)}
\end{equation}
where $\Sigma^{-1}$ it's the inverted covariance matrix of the dataset.
\item Jaccard distance, used in the case of cathegorical features, once they have been encoded using one-hot encoding
\begin{equation}
    \label{eq:jaccard}
    d\left(\mathbf{x}_1, \mathbf{x}_2\right) = 1 - J\left(\mathbf{x}_1, \mathbf{x}_2\right)
\end{equation}
where 
\begin{equation}
J\left(\mathbf{x}_1, \mathbf{x}_2\right) = \frac{\left|\mathbf{x}_1 \cap \mathbf{x}_2\right|}{\left|\mathbf{x}_1 \cup \mathbf{x}_2\right|}    
\end{equation}
\end{itemize}
There are several different clustering algorithms that can be divided into two main groups, as in \citeauthor{james_introduction_2021} (2021): K-Means clustering and hierarchical clustering. In \emph{K-Means clustering}, we try to divide observation into a specified number of clusters. On the other hand, in \emph{hierarchical clustering}, we do not specify the number of clusters in advance, but we rely on the dendrogram, a tree-like visual representation of all clustering steps. Then, we can "cut" the dendrogram in order to select the most appropriate number of clusters for us. In Fig. \ref{fig:clustiteration}, there is an example of simulated data in the plane clustered into 3 clusters using the K-Means algorithm (that will be discussed later).
% <<< End of Clustering algorithm

%%%%%
%%%%%

% K-means algorithms >>>
\subsubsection{K-Means Clustering Algorithm}
\label{sss:kmeans}
We will delve deeper into it in a dedicated section since we will use K-Means in Section \ref{sec:bvc} and in Appendix \ref{ap:Python}. Let $C_1,\dots,C_k$ denote the sets containing all the indices of the observation in respective the respective cluster. K-Means clustering is a straightforward clustering algorithm we can use for complete partitioning of a data collection of observations $E=\{\mathbf{x}_1, \dots, \mathbf{x}_n \}$, described by $p$ features, into $K$ distinct, non-overlapping clusters. This means that:
\begin{enumerate}
    \item $C_1 \cup C_2 \cup \dots \cup C_K=\{1,\dots,n\}$. So, each observation belongs to at least one of the $K$ clusters.
    \item $C_k\cap C_{k'}=0$ for all $k\neq k'$. In other words, no observation belongs to more than one cluster.
\end{enumerate}
Thus, each observation will be assigned to one cluster and one cluster only. Moreover, each cluster $C_i$ has a representative point that summarizes the cluster. That's the centroid $\bm{\mu}_i$
\begin{equation}
    \label{eq:centroid}
    \boldsymbol{\mu}_i=\frac{1}{\left|C_i\right|} \sum_{x_j \in C_i} \mathbf{x}_j
\end{equation}
K-means begins by randomly placing $\bm{\mu}_i \in \mathbb{R}^p$ for $i=1,\dots,K$ points as initial cluster centroids. This is usually achieved by generating a random value within the bounds of each dimension. 
\begin{figure}
    \centering
    \includegraphics[width=0.6\textwidth]{Images/clustiteration.png}
    \caption[K-Means clustering iterations.]{The progress of the K-means algorithm with $K=3$ \cite{james_introduction_2021}.}
    \label{fig:clustiteration}
\end{figure}
The K-means algorithm progresses through iterative cycles involving two primary actions: (i) assigning data points to clusters and (ii) updating centroids. Given the $K$ centroids, the assignment phase allocates each data point to its nearest centroid, resulting in distinct clusters, where each cluster $C_i$ consists of points closer to its centroid $\bm{\mu}_i$ than to any other centroid. Then, in the update phase, each cluster's centroid is recalculated based on its (updated) member data points. This iterative process of assignment and updating continues until the centroids stabilize. In practical terms, K-means can be considered stable when there's no shift in centroid positions between consecutive iterations. We can adjust the tolerance parameter $\epsilon \geq 0$ for early termination of the algorithm or let the algorithm converge to the local optimum until observations no longer change cluster. In algorithm \ref{alg:kmeans}, there is the pseudo-code of the K-Means clustering algorithm.
\begin{algorithm}
\small
    \caption{K-Means Clustering}
    \label{alg:kmeans}
    \begin{algorithmic}[1]
    \STATE {$t=0$}
    \STATE {Randomly initialize $K$ centroids: $\bm{\mu}_1^t, \bm{\mu}_2^t, \dots, \bm{\mu_k^t} \in \mathbb{R}^p$}
    \REPEAT
    \STATE {$t \leftarrow t + 1$}
    \STATE {$C_i \leftarrow \empty$ for all $i=1,\dots,K$}
    \STATE {//Cluster assignment step}
    \FORALL{$\mathbf{x}_j\in D$}
    \STATE {$i^* \leftarrow \arg \min _i\left\{\left\|\mathbf{x}_j-\boldsymbol{\mu}_i^{t-1}\right\|^2\right\}$}
    \STATE {$C_{i^*} \leftarrow C_{i^*} \cup\left\{\mathbf{x}_j\right\} / / \text { Assign } \mathbf{x}_j \text { to closest centroid }$}
    \ENDFOR
    \STATE {//Centroid update Step}
    \FORALL{$i=1,\dots,K$}
    \STATE{$\boldsymbol{\mu}_i^t \leftarrow \frac{1}{\left|C_i\right|} \sum_{\mathbf{x}_j \in C_i} \mathbf{x}_j$}
    \ENDFOR
    \UNTIL{$\sum_{i=1}^k\left\|\boldsymbol{\mu}_i^t-\boldsymbol{\mu}_i^{t-1}\right\|^2 \leq \epsilon$}
    \end{algorithmic}
\end{algorithm} 
Since the initial guess of the centroids can lead to different labeling of the same observation, the K-means algorithm is typically run several times. In bagging clustering, we run the algorithm multiple times, and the final label is defined by majority voting across the different intermediate clusterizations.
% <<< End of K-means clustering algorithmclassification


% Principal component analysis >>>
\subsection{Principal Component Analysis}
\label{subsec:PCA}
Another unsupervised machine learning technique is the principal component analysis (PCA). This method leverages the information patterns within the dataset to reduce the problem's dimensionality. Principal components (PC) of a set of data in $\mathbb{R}^p$ provide a sequence of best linear approximation to that data of all ranks $q\le p$ \cite{james_introduction_2021, tibshirani_elements_2008}. \emph{Principal component analysis} refers to the process by which principal components are computed. Suppose we have a dataset that represents experience $\mathbf{E}$ of a certain phenomenon, of unlabeled data and each realization $\mathbf{x}_j \in \mathbb{R}^p$ for all $j=1,\dots,n$. Thus, $\mathbf{E} \in \mathbb{R}^{n\times p}$. The first principal component of a set of features $\mathbf{X}_1, \mathbf{X}_2, \dots, \mathbf{X}_p$ is the normalized linear combination of the features
\begin{equation}
    \label{eq:firstPC}
    Z_1=\phi_{11} \mathbf{X}_1+\phi_{21} \mathbf{X}_2+\cdots+\phi_{p1} \mathbf{X}_p
\end{equation}
By normalized, we mean that $\sum_{j=1}^p\phi_{j1}^2=1$. The elements $\phi_{11},\dots,\phi_{p1}$ are the loadings of the first principal component. But how can we find PCs? Suppose we want to find a linear model representing the $i$-th observation. We are trying to find a linear combination such as:
\begin{equation}
    \label{eq:pcaf}
    f(\lambda)=\bm{\mu}+\mathbf{V}_q\bm{\lambda}
\end{equation}
where $\bm{\mu}\in \mathbb{R}^p$ is a location vector, $\mathbf{V}_q$ is a $p\times q$ matrix with $q$ orthogonal unit vectors as columns, and $\bm{\lambda \in \mathbb{R}^q}$ is a vector of $q$ parameters. So, we can fit the model to the data by minimizing the sum of least squared errors:
\begin{equation}
    \label{eq:SSE}
    \min _{\bm{\mu},\left\{\lambda_i\right\}, \mathbf{V}_q} \sum_{j=1}^N\left\|\mathbf{x}_j-\bm{\mu}-\mathbf{V}_q \lambda_j\right\|^2
\end{equation}
We can partially optimize for $\bm{\mu}$ and the $\lambda_j$ to obtain
\begin{equation}
    \hat{\bm{\mu}} = \left(E\left(\mathbf{X}_1\right), E\left(\mathbf{X}_1\right), \dots, E\left(\mathbf{X}_p\right)\right)
\end{equation}
\begin{equation}
    \hat{\lambda}_j = \mathbf{V}_q^{T}(\mathbf{x}_j-\hat{\bm{\mu}})
\end{equation}
This leads to 
\begin{equation}
\label{eq:diocaro}
    \min _{\mathbf{V}_q} \sum_{i=1}^N\left\|\left(\mathbf{x}_j-\hat{\bm{\mu}}\right)-\mathbf{V}_q \mathbf{V}_q^T\left(\mathbf{x}_j-\hat{\bm{\mu}}\right)\right\|^2
\end{equation}
The $q \times q$ matrix $\mathbf{H}_q=\mathbf{V}_q\mathbf{V}_q^T$ is a projection matrix and maps each point $\mathbf{x}_j$ onto it's rank~-~$q$ reconstruction $\mathbf{H}_qx_j$, the orthogonal projection of $x_j$ onto the subspace spanned by the columns of $\mathbf{V}_q$. The solution can be expressed as follows. Stacked the centered observations into rows of an $N\times p$ matrix $\mathbf{E}$, we can construct the singular value decomposition
\begin{equation}
    \label{eq:SVD}
    \mathbf{E}=\mathbf{U}\bm{\Sigma}\mathbf{V}^T
\end{equation}
where $\mathbf{U}$ and $\mathbf{V}$ are unitary matrices and $\bm{\Sigma}$ is a $p \times p$ diagonal matrix, such that $\sigma_{11}\ge \sigma_{22} \ge \dots \ge \sigma_{pp}\ge 0$. This means that $\mathbf{u}_1$ and $\mathbf{v}_1$ correspond to $\sigma_{11}$ and are somehow more important with respect to $\mathbf{u}_2$ and $\mathbf{v}_2$ which are associated to $\sigma_{22}$ in explaining the information in $\mathbf{E}$ and so on and so forth. The columns of $\mathbf{U}$ are the basis vector of the new space, and the columns of $\mathbf{U}\mathbf{D}$ are the principal components. The $\sigma_{ii}$ are also called singular values and are the variance the $i$-th PC explains. Lastly, $\mathbf{v}_i=(\phi_{11} \dots \phi_{p1})$ re the loading vectors of such $i$-th PC.  For each rank $q$, the solution $\mathbf{V}_q$ to Eq. \ref{eq:diocaro}, consist of the first $q$ columns of $\mathbf{V}$. A good method to determine the number $q$ of principal components (PC) is to rely on the scree plot of explained variance. Indeed, the variance explained by each additional PC follows the law of diminishing returns. Therefore, we can decide how many PCs when we observe a significant knee in the graph.

\subsubsection{Handwritten Digits}
This example is  from \citeauthor{tibshirani_elements_2008} (2008).
PCA can be applied to images as well. For instance, take a dataset containing handwritten digits. Let's consider its application on the digit "3". Authors experimented on a set of 130 3s randomly sampled from a total of 658 total observations. Each image is a \numproduct{16 x 16} gray-scale image. We can see considerable variation in writing styles, character thickness, and orientation (Fig. \ref{fig:3s}). We can unfold images and treat each of them as vectors $\textbf{x}_j \in \mathbb{R}^{256}$, and we can use PCA on the dataset $\mathbf{E}\in\mathbb{R}^{256\times 130}$.
\begin{figure}
    \centering
    \subfloat[\label{fig:3s}]{
        \includegraphics[scale=0.35]{Images/1303.png}
    }
    \qquad
    \subfloat[\label{fig:principal3}]{
        \includegraphics[scale=0.5]{Images/pc3.png}
    }
    \caption[Example of PCA of images.]{Example of PCA applaied to images dataset: the random sample of handwritten 3s (a) and the scatterplot of the first two principal components (b) with a superimposed grid in which images match red circles on the scatterplot\cite{tibshirani_elements_2008}.}
\end{figure}
From Fig. \ref{fig:principal3}, we can see that the first PC mainly accounts for the length of the lower tail of the three, while the second PC accounts for character thickness. Then, we can reshape every PC in a matrix and interpret them as images. These images will be the base for the new vectorial space. Using just 50 PCs, we can explain 90\% of the variance of all the 3s, while with 12 PCs, we can account for 63\% of the total variance.
% <<< End of Principal component analysis

%%%%%
%%%%%


% Deep Learning, ANN and CNN >>>
\subsection{Artificial Neural Network and Deep Learnig}
\label{subsec:deepl}
In 1956, there was a summer research project on artificial intelligence at Dartmouth University. A team of professors and an expert from IBM Corporation began to think about the possibility of creating virtual neural networks. These networks would learn and enhance the purpose for which they were designed through a self-improvement process, and they could be based on the mathematical model of the human brain \cite{mccarthy_proposal_1955}. Indeed, the human brain model is especially well-suited for computational purposes. The human brain possesses a massive amount of computing units: the neurons. Approximately \num{e11} neurons in an adult human brain form about 7,000 synaptic connections with other neurons, amounting to nearly \num{e15} total synapses. Moreover, the computational model of the brain is distributed among simple nonlinear units, redundant and fault-tolerant, and capable of performing calculations in parallel \cite{matteo_matteucci_perceptrons_2021}. The mathematical virtual model of a single neuron is called a perceptron, the unitary computational unit of artificial neural networks (ANN). Just like in human neurons, the perceptron has dendrites that "collect" the charge from synapses, and once a certain threshold is exceeded, the accumulated charge is released and sent to other perceptrons. The portion of the schema in Fig. \ref{fig:perceptron} enclosed in the dashed line is called the activation function and is responsible for releasing the perceptron's accumulated charge. There are various activation functions, with some of the most common ones can be seen in Fig. \ref{fig:actfunc}.
\begin{figure}
    \centering
    \includegraphics[width=0.6\textwidth]{Images/neurone.png}
    \caption[Perceptron schema]{Perceptron schema \cite{matteo_matteucci_perceptrons_2021}.}
    \label{fig:perceptron}
\end{figure}
Perceptrons can be combined to form networks with a layered architecture. Indeed, every neural network will have an input layer, an output layer, and a varying number of intermediate layers, known as hidden layers. Deep learning is a subfield of artificial intelligence (AI) and machine learning that involves training artificial neural networks on vast amounts of data to perform classification or regression tasks. ANNs can learn from experience $\mathbf{E}$ and understand complex patterns in data. Deep feedforward networks, also called simply feedforward neural networks or multilayer perceptrons (MLPs), are used to approximate some function $f^*$. Let's take as an example a classifier. Given an unknown function $\Lambda_0:E \rightarrow \{1,\dots,L\}$ where $\{1,\dots,L\}$ are label classes, the $f^*$ is the best approximation of the function $\Lambda_0$, defining a mapping $\mathbf{y}=f^*\left(\mathbf{x}, \bm{\theta} \right)$. The ANN learns the value of parameters $\bm{\theta}$ from $\mathbf{E}$, resulting in the best function approximation. These models are called feedforward because information flows through forward all in the network. To understand how neural networks learn from experience $\mathbf{E}$ using the chain rule and the back-propagation mechanism, readers are referred to \citeauthor{goodfellow_deep_2016} (2016).
\begin{figure}
    \centering
    \includegraphics[width=0.8\textwidth]{Images/activationfunction.png}
    \caption[Activation functions.]{Some examples of activation functions with relatives analytical forms \cite{matteo_matteucci_perceptrons_2021}.}
    \label{fig:actfunc}
\end{figure}
Let's examine examples of how networks can address classic machine learning problems such as regression and classification. A best practice is to use tanh or sigmoid as hidden layers activation functions. Indeed, the main difference is the activation function of the output layer: 
\begin{itemize}
    \item \textbf{Regression:} since the output of a regression problem spans all $\mathbb{R}$, the activation function of the output layer will be a linear activation function;
    \item \textbf{Binary Classification:} in binary classification, the output domain is $\Omega=\{0, 1\}$. Thus, the necessary activation function for the output layer is the sigmoid activation function. The value of the function can be interpreted as class posterior probability;
    \item \textbf{Multi-class Classification:} when dealing with multiple classes ($K$), we need to use as many neurons as classes, and each output neuron will use a softmax unit. The unique feature of this activation function is that the sum of the outputs from the $K$ output neurons is equal to 1. Thus, the $k$-th output can be viewed as the probability that the input belongs to the $k$-th class.
\end{itemize}
Another typical application of neural networks is as classification algorithms for images in computer vision. However, images cannot be used directly as input since they carry too much information. We need some intermediate steps to extract useful information and reduce the problem's dimensionality by computing a feature array that summarizes the initial image's information \cite{giacomo_boracchi_convolutional_2021}. The feature array for the input layer of the classification algorithm will have a dimension of $d\ll r_1 \times c_1$ where $r_1$ and $c_1$ are the rows and the columns of the matrix representation of the image. Recall that an image can be represented as a three-dimensional matrix: the first two dimensions are pixel indexes, and the third identifies the pixel's color based on the color profile used. Take, for instance, the image in figure \ref{fig:cazzle}. The image has a width of 472 pixels a height of 376 pixels, and uses an RGB color profile. This means that three values are needed to identify the color of each pixel uniquely: the first for red, the second for green, and the third for blue. This means that we would require \numproduct{472 x 376 x 3} input neurons, one for each value of each color of each pixel of the image (more than 500,000 neurons). Feature extraction algorithms can be categorized as hand-crafted or data-driven features. With hand-crafted algorithms, we can leverage prior knowledge of the phenomenon, and features are interpretable, allowing us to exploit them even with limited data. However, this approach demands significantly more time and design/programming efforts. Moreover, it is very domain-specific and isn't so much "portable". On the other hand, data-driven feature extraction algorithms are less interpretable but offer a more robust mathematical foundation. Convolutional Neural Networks (CNNs) are neural networks that can be used for data-driven feature extraction. Given the matrix representation of the image discussed earlier, while in ANN, we talked about layers, in CNNs, we will have volumes. As the depth of such volumes increases, the height and width of the volume decreases. CNNs enable us to perform some fundamental operations essential for accurate feature extraction.
\paragraph{Convolution.} Convolutional layers "mix" all the input components. Convolution is an operation that allows us to combine all the input values from a specific region into a linear combination. Specifically, the convolution is described by the following equation:
\begin{equation}
    \label{eq:convolution}
    y_j=\sum_iw_{ij}x_i+b_j
\end{equation}
The linear combination operation is also referred to as filtering, and the parameters of this combination are called filters. The same filter is used as a sliding window through the whole spatial extent of the input images. In Fig. \ref{fig:convmatrice}, a typical convolution operation is graphically represented.
\begin{figure}
    \centering
    \subfloat[\label{fig:cazzle}]{
    \includegraphics[width=0.7\textwidth]{Images/featureextraction.png}
    }
    \qquad
    \subfloat[\label{fig:convmatrice}]{
    \includegraphics[width=0.5\textwidth]{Images/convolution.png}}
    \caption[Convolution operation.]{Feature extraction operation (a) adapted from \cite{giacomo_boracchi_convolutional_2021}, and convolution operation graphically represented (b)\cite{giacomo_boracchi_convolutional_2021}.}
\end{figure}
In practice, the convolution operation can be viewed as a two-dimensional window that slides across the entire image, producing an output that is a linear combination of the matrix values and their corresponding values.
\paragraph{Poolig layers.} Pooling layers reduce the spatial size of the input volume. Using the logic of a sliding window that traverses the image, a specific function is applied, often the maximum, to spatially resize the input.
\begin{figure}
    \centering
    \includegraphics[width=0.5\textwidth]{Images/pooling.png}
    \caption[Pooling layer.]{Example of pooling layer that performs a max operation \cite{giacomo_boracchi_convolutional_2021}.}
    \label{fig:poolingmax}
\end{figure}
The stride parameter specifies how many elements the window shifts at once. In the figure, the window will move by two pixels simultaneously on the y-axis.
\paragraph{Activation layers.} Activation layers are used to introduce non-linearities in the network; otherwise, CNN might be equivalent to a simple linear classifier. In practice, they are layers of neurons characterized by activation functions that introduce non-linearity into the levels, just as in ANNs. We have yet to discuss two commonly used functions in these layers: the RELU (rectified linear units) and LEAKY RELU activation. The former is described as
\begin{equation}
    \label{eq:RELU}
    T(x)=\left\{\begin{array}{rr}
    x, & \text { if } x \geq 0 \\
    0, & \text { if } x<0
\end{array}\right.
\end{equation}
Whereas LEAKY RELU is similar to RELU but incorporates a slight slope for negative values, preventing neurons from dying off and becoming unresponsive after several layers. For a deeper understanding of this topic, one can refer to the vanishing gradient problem as discussed by Goodfellow.
\begin{equation}
    \label{eq:LEAKYRELU}
    T(x)=\left\{\begin{array}{rr}
    x, & \text { if } x \geq 0 \\
    k\cdot x, & \text { if } x<0
\end{array}\right.
\end{equation}
where $k$ is usually in the order of $10^{-2}$.
The convolution output will be a dense feature layer, which implies that this layer can serve as input for an ANN, which will have as many output neurons as classes in the original domain. Fig. \ref{fig:typicalcnn} shows an example of how a general CNN works.
\begin{figure}
    \centering
    \includegraphics[width=0.8\textwidth]{Images/CNNtypical.png}
    \caption[CNN general structure.]{Graphical representation of a tipical CNN used for image classification \cite{giacomo_boracchi_convolutional_2021}.}
    \label{fig:typicalcnn}
\end{figure}
% End of Deep learning, ANN and CNN

%%%%%
%%%%%

% Hot spot >>>
\section{Hot-Spot Detection Using ML Algorithm: A-State-Of-The-Art Review}
\label{sec:hotspotstateart}
% <<< End of Hot spot  
\backgroundsetup{contents={}}
% <<< end of PBF Hot Spot State of the Art


% Proposed method Bagging Voronoi >>>
\chapter{Bagging Voronoi Clustering Algorithm}
\label{ch:baggingvoronoi}
\backgroundsetup{contents={DRAFT}}
This chapter presents a novel approach proposed by the author for clustering building plate areas characterized by similar thermal profiles, firstly described in \citeauthor{secchi_bagging_2013} (2013). In general, it can be applied to any functional data referenced by a lattice in $\mathbb{R}$, $\mathbb{R}^2$ or $\mathbb{R}^3$. The analysis of functional data allows us to gain a comprehensive and detailed understanding of the anomalies and the thermal history, which would enable us to discern if there are specific zones on the print plane behaving irregularly and to understand how hot spots might influence the mechanical properties of the printed part, given the fact that we have a functional knowledge of those phenomena. Section \ref{sec:fda} briefly introduces functional data analysis. Section \ref{sec:bvc} gives a detailed explanation of the Bagging Voronoi algorithm for geo-referenced functional data clustering, and Section \ref{sec:irradiance} presents a successful case study of the algorithm application.

% Functional Data Analysis >>>
\section{Functional Data Analysis}
\label{sec:fda}
This section discusses Functional Data Analysis (FDA) and Functional Principal Component Analysis (FPCA). It provides a brief overview to give readers the foundational knowledge needed to understand the algorithm described in Section \ref{sec:bvc}. This section is based mainly on \citeauthor{ramsay_functional_2006} (2006) and \citeauthor{ramsay_functional_2009} (2009).

Functional data are measurable observations of a given phenomenon that we can assume, with a fair degree of certainty, that a function can describe it. Indeed, although each observation is a single discrete value, these values reflect the smooth behavior of the generative function of those data. Fig. \ref{fig:bambine} shows the height data in \unit{\centi\metre} of 10 girls measured at 31 different temporal instants in the Berkeley Growth Study. We can see the smooth behavior typical of analytic functions, suggesting a "height" function, relating height to girls' age. Like all mathematical functions, the height function has interesting properties, including continuity and the possibility to compute function derivatives. Indeed, we can calculate the first derivative of the height function and interpret it as the growth velocity. Similarly, we can calculate the second derivative, the growth acceleration, as Fig. \ref{fig:accelerazione} shows.
\begin{figure}
    \centering
    \subfloat[\label{fig:bambine}]{
        \includegraphics[scale=0.28]{Images/crescita.png}
    }
    \qquad
    \subfloat[\label{fig:accelerazione}]{
        \includegraphics[scale=0.28]{Images/accelerazione.png}
    }
    \caption[Functional data]{The heights of 10 girls measured at 31 ages (a) and their estimated accelerations (b) \cite{ramsay_functional_2006}. The dotted line is the cross-sectional mean acceleration.}
    \label{fig:tuttoinsieme}
\end{figure}
Even if the height observations are taken nominally at the same age for each girl, it is not strictly needed. Indeed, the points at which the function is observed may change from one subject to another. It can be helpful in quality control that relies on sensors we discussed in Section \ref{sec:sensoriniiniini} to deal with missing values due, for example, to the obstruction of the off-axis sensor by the recoating blade. Furthermore, the signals transmitted by these sensors may be affected by disturbances that prevent the observed quantity from being correctly measured, causing discontinuities or missing values. It is also important to note that the FDA does not suffer from the problems of classical multivariate methods. It can also be used when the number of variables $p$ is larger than the available sample size $n$. In the case of Fig. \ref{fig:tuttoinsieme}, we have $p=31$ and $n=10$. In general, the goal of the FDA is to represent data as functions, to provide a representation that helps to analyze and highlight various characteristics of the collected data. This allows us to explain patterns of variation in the data and compare data containing different sets of observations of the same function or functions for a standard set of replicates. We can write each functional observation as the sum of a functional component and a random component that we will call noise. Formally, we can write every observed profile $y_j$ as 
\begin{equation}
    y_j = f(t_j) + \varepsilon_j \qquad \varepsilon_j\sim N(0, \sigma_j^2)
\end{equation}
We can approximate the unknown functional component as a linear combination of a functional basis. A basis function system is a set of known functions $\phi_k$ such that they are mathematically independent of each other and have the property that they can arbitrarily well approximate any function by taking a weighted sum (linear combination) of a sufficiently large number $K$ of these functions.
\begin{equation}
\label{eq:linear}
    \hat{f}(t)=\sum_{k=1}^K c_k \phi_k(t)
\end{equation}
Now, let $\mathbf{c}$ indicate the vector of length $K$ of the coefficients $c_k$ and $\bm{\phi}$ as the functional vector whose elements are the basis functions $\phi_k$, we can also express \ref{eq:linear} in matrix notation as
\begin{equation}
\label{eq:matrice}
    \hat{f}(t)=\mathbf{c}'\bm{\phi}=\bm{\phi}'\mathbf{c}
\end{equation}
When $K=n$, we have a limit case called exact representation or interpolation. It is a limit situation because we can choose the coefficients of the weighted sum such that $\hat{f}(t_j)=y_j$ for each $j$. Thus, $K$ controls the smoothness of the approximation $\hat{f}(t)$. We have to choose this parameter according to the used functional basis and the characteristics of the data. Another parameter to be selected is the type of functional basis itself. There are different functional bases, each with different characteristics, which ideally reflect the features of observed data. In a certain way, we can use the parameter $K$ to understand whether the used basis was more or less suitable: the smaller $K$ is, the better the base reflects the characteristics of the data. The type of basis selected is critical when one is interested in analyzing derivatives. Formally, when we are interested in
\begin{equation}
    \label{eq:derivative}
    D \hat{f}(t)=\sum_k^K \hat{c}_k D \phi_k(t)=\hat{\mathbf{c}}^{\prime} D \phi(t)
\end{equation}
where $D$ is the derivative operator defined as $Df(t)=\frac{d}{dt}f(t)$. This operator can also define further derivatives as powers of $D$.
\begin{equation}
    \label{eq:deroperator}
    D^2f(t)=\frac{d^2}{dt^2}f(t), \dots, D^kf(t)=\frac{d^k}{dt^k}f(t)
\end{equation}
The following are some of the most common functional bases used, highlighting their characteristics.
\paragraph{Polynomial basis.} The polynomial basis consists of monomials of different order.
\begin{equation}
    \label{eq:polinomial}
    \phi_k(t)=(t-\omega)^k \qquad k=0,\dots, K
\end{equation}
where $\omega$ is a shift parameter that is usually chosen to be in the center of the interval of approximation. The polynomial basis tends to fit the center of the data very well, but the tails tend to exhibit less ideal behavior. This is why they are never used for extrapolation or forecasting. However, the basis has a remarkable behavior: if $t_j$ are equally spaced, the basis can be fit with minimal computational effort. Moreover, derivatives are simple to compute and tend to be good estimators for actual derivatives, although they may present a wiggling effect when $K$ is excessively large. 
\paragraph{Exponential and power basis.} The exponential base consists of a series of exponential functions 
\begin{equation}
    \label{eq:exponential}
    \phi_k(t)=\exp (\lambda_kt) \qquad k=1,\dots, K
\end{equation}
where the rate parameters $\lambda_k$ are all distinct and often $\lambda_1=0$. This basis is particularly useful as linear differential equations with constant coefficients have expansions in terms of an exponential basis as solutions. While power basis
\begin{equation}
    \label{eq:power}
    \phi_k(t)=t^{\lambda_k} \qquad k=1,\dots, K
\end{equation}
is less useful, but can be helpful when $t$ is strictly positive.
\paragraph{Fourier basis.} One of the most famous and widely used bases is the Fourier basis, mainly because it can capture the periodic behavior of data and is very useful for stable functions when curvatures tend to have the same order in all the data. However, the Fourier basis is unsuitable for data known or suspected to reflect some discontinuities. Fourier basis is defined as:
\begin{equation}
    \label{eq:fourier}
    \phi_k(t)=\left\{ \begin{array}{ll}
         \phi_k(t)=1 & \text{for} \quad k = 0 \\
         \phi_{2k-1}(t)= \sin(k\omega t) & \text{for} \quad k = 1, \dots, K \\
         \phi_{2k}(t)= \cos(k\omega t) & \text{for} \quad k = 1, \dots, K
    \end{array}\right.
\end{equation}
The period of the basis is determined as $\mathcal{T}=\frac{2\pi}{\omega}$. The algorithm used to calculate the coefficients of this basis, the Fast Fourier transform, makes it possible to find the coefficients extremely efficiently when $n$ is a power of 2 and the observations are equally spaced.
\paragraph{B-splines basis.} The B-splines base is a splines base. \emph{Splines} are bases defined by a number $L$ of knots and an order $m$. Knots are values $\tau_l$ for $l=1,\dots, L-1$ that divide the interval a function has to be estimated into $L$ sub-intervals. On each sub-interval, the spline is a polynomial function of a specific order $m$, the degree of the polynomial itself. Polynomials defined on adjacent intervals join smoothly at knot values, such that the function at that point is constrained to the value of the knot itself. By increasing the number of knots, the fit of the estimated function to the observed data can be improved up to the extreme situation where the knot numbers have the same value as the observed data. Thus, the number of nodes will affect the goodness of fit to the observed data. In a certain way, the parameter $L$ regulates the flexibility of our fit. This is why having more nodes is a good practice for regions with complex variations. Thus, to summarise what has been said so far, a spline is defined by the order of the polynomials and the sequence of nodes. \emph{B-splines} is a particular spline basis developed by de Boor in 2001. 
\begin{figure}
\centering
\includegraphics[width=0.4\textwidth]{Images/bspline.png}
\caption[Bsplines basis.]{Thirteen basis functions defining an order four spline with nine interior knots, shown as vertical dashed lines \cite{ramsay_functional_2006}.}
\label{fig:bspline_13}
\end{figure}
\begin{figure}
\centering
\includegraphics[width=0.7\textwidth]{Images/bsplines.png}
\caption[Bsplines fitting.]{In the left panels, the solid line indicates the spline function of a particular order that fits the sine function shown as a dashed line. The corresponding fits to its derivative, a cosine function, are shown in the right panels. The vertical dotted lines are the interior breakpoints or knots defining the spline fits \cite{ramsay_functional_2006}.}
\label{fig:bsplines}
\end{figure}
From Fig. \ref{fig:bspline_13}, we can see that the three functions on the left and the last three on the right are different from the other. When transitioning from the left edge toward the center, the spans in which these functions remain positive expand from one to four. However, they consistently have the same smooth transition that's twice-differentiable. On the other hand, if we move from the center to the left or right, functions have different levels of smoothness: the far-left spline is discontinuous, the following one is only continuous, and the last is differentiable once. It makes perfect sense that we experience a loss in differentiability at the edges, given that we typically lack knowledge about the behavior of the function we are approximating outside the interval.
\paragraph{Wavelets basis.} We can construct a basis for all functions on $\left(- \infty, \infty \right)$ that are square integrable by choosing a suitable mother wavelet function $\psi$ and then considering all dilations and translations of the form
\begin{equation}
    \label{eq:wavelet}
    \psi_{j k}(t)=2^{j / 2} \psi\left(2^j t-k\right)
\end{equation}
for integers $j$ and $k$. The wavelet expansion of a function $f$ gives a multiresolution analysis in the sense that the coefficient of $\psi_{jk}$ yields information about $f$ near position $2^{-j}k$ on scale $2^{-j}$, i.e., at frequencies near $c2^j$ for some constant $c$. Wavelets are suitable for modeling sharp local features, even discontinuity, localized in space and frequency, which can be computed efficiently. Multidimensional wavelets provide coherent estimates of ECG traces, where the different projections of the heart dynamics are consistent.
\subsection{Model Fitting}
There are different ways to fit a model to observed data. When we have to estimate a model, we want the model to fit our data somehow. To quantify how well it fits, we can define the concept of sum of squared errors defined as
\begin{equation}
    \label{eq:sse}
    \operatorname{SSE}(\mathbf{y} \mid \mathbf{c})=\sum_{j=1}^n\left[y_j-\sum_k^K c_k \phi_k\left(t_j\right)\right]^2
\end{equation}
or equivalently, we can re-write it in matrix form as
\begin{equation}
    \label{eq:ssematrix}
    \operatorname{SSE}(\bm{y} \mid \mathbf{c})=(\mathbf{y}-\bm{\Phi} \mathbf{c})^{\prime}(\mathbf{y}-\bm{\Phi} \mathbf{c})
\end{equation}
where $\bm{\Phi}$ is the $n \times K$ matrix containing the values $\phi_k(t_j)$. By deriving equation \ref{eq:ssematrix} with respect to $c$, we obtain equation
\begin{equation}
 \label{eq:c}
    2 \bm{\Phi} \bm{\Phi}^{\prime} \mathbf{c}-2 \bm{\Phi}^{\prime} \mathbf{y}=0
\end{equation}
and solving for $\mathbf{c}$, this provides the estimate $\hat{\mathbf{c}}$ that minimizes the least squares solution
\begin{equation}
    \label{eq:cestimated}
    \hat{\mathbf{c}}=\left(\bm{\Phi}^{\prime} \bm{\Phi}\right)^{-1} \bm{\Phi}^{\prime} \mathbf{y}
\end{equation}
The vector $\hat{\mathbf{y}}$ of fitted values is
\begin{equation}
    \hat{\mathbf{y}}=\bm{\Phi} \hat{\hat{\mathbf{c}}}=\bm{\Phi}\left(\bm{\Phi}^{\prime} \bm{\Phi}\right)^{-1} \bm{\Phi}^{\prime} \mathbf{y}
\end{equation}
If we want to fit a model on data with a non-stationary and/or auto-correlated stochastic component, however, we must use weighted least squares fits. In this case, equation \ref{eq:ssematrix} becomes
\begin{equation}
    \label{eq:wls}
    \operatorname{SSE}(\mathbf{y} \mid \mathbf{c})=(\mathbf{y}-\bm{\Phi} \mathbf{c})^{\prime} \mathbf{W}(\mathbf{y}-\bm{\Phi} \mathbf{c})
\end{equation}
where the matrix $\mathbf{W}$ is a symmetric positive definite matrix that allows for unequal weighting of squares and products of residuals. The matrix $\mathbf{W}$ is defined as
\begin{equation}
    \mathbf{W}=\bm{\Sigma}_e^{-1}
\end{equation}
where $\bm{\Sigma}_e^{-1}$ is the variance-covariance matrix for the noise. As before, if we derive \ref{eq:wls} with respect to $\mathbf{c}$ and by solving the equation, we get the weighted least squares estimate $\hat{\mathbf{c}}$ of the coefficient vector $\mathbf{c}$ defined as
\begin{equation}
    \hat{\mathbf{c}}=\left(\bm{\Phi}^{\prime} \mathbf{W} \bm{\Phi}\right)^{-1} \bm{\Phi}^{\prime} \mathbf{W} \mathbf{y}
\end{equation}
Finally, in the case of a splines basis, we can smooth the model to the data by introducing a roughness penalty. We have already encountered the derivative operator in \ref{eq:deroperator}. The first two derivative operators are particularly relevant because  $Df(t)$ is the local slope of $f(t)$ and $D^2f(t)$ is its curvature. We can therefore define the roughness of the $m$-th derivative as
\begin{equation}
    \label{eq:roughness}
    \begin{aligned}
\operatorname{PEN}_m(x) & =\int\left[D^m x(s)\right]^2 d s \\
& =\int\left[D^m \mathbf{c}^{\prime} \boldsymbol{\phi}(s)\right]^2 d s \\
& =\int \mathbf{c}^{\prime} D^m \boldsymbol{\phi}(s) D^m \boldsymbol{\phi}^{\prime}(s) \mathbf{c} d s \\
& =\mathbf{c}^{\prime}\left[\int D^m \boldsymbol{\phi}(s) D^m \boldsymbol{\phi}^{\prime}(s) d s\right] \mathbf{c} \\
& =\mathbf{c}^{\prime} \mathbf{R} c
\end{aligned}
\end{equation}
We can then find the right trade-off between the fitting of the data and the smoothness of fitted curve by minimising the penalised sum of squared errors
\begin{equation}
    \label{eq:pensse}
    \operatorname{PENSSE}_\lambda(f)=[\mathbf{y}-f(\mathbf{t})]^T[\mathbf{y}-f(\mathbf{t})]+\lambda \operatorname{PEN}_2[f]
\end{equation}
where $\lambda$ is a smoothing parameter measuring the compromise between fit and smoothness. As $\lambda$ increases, roughness is increasingly penalized, and $\hat{f}(t)$ will become linear. As it decreases, the penalty is reduced, allowing $\hat{f}(t)$ to fit the data better. Fig. \ref{fig:lambda} shows different fit of yearly precipitations data of Vancouver obtained by using different values of $\lambda$.
\begin{figure}
    \centering
    \subfloat[\label{fig:-1}]{\includegraphics[width = 0.3\textwidth]{Images/-1.png}}
    \subfloat[\label{fig:7}]{\includegraphics[width = 0.3\textwidth]{Images/7.png}}
    \subfloat[\label{fig:15}]{\includegraphics[width = 0.3\textwidth]{Images/15.png}}
    \caption[PENSSE with different $\lambda$]{PENSSE on Vancouver precipitations data with $\lambda=-1$ (a), $\lambda=7$ (b) and $\lambda=15$ (c) \cite{ramsay_functional_2009}.}
    \label{fig:lambda}
\end{figure}
Finally, the estimated coefficient vector $\hat{\mathbf{c}}$ will be 
\begin{equation}
    \hat{\mathbf{c}}=\left(\boldsymbol{\Phi}^{\prime} \mathbf{W} \boldsymbol{\Phi}+\lambda \mathbf{R}\right)^{-1} \boldsymbol{\Phi}^{\prime} \mathbf{W} \mathbf{y}
\end{equation}
and thus the data-fitting vector $\mathbf{\hat{y}}$ is
\begin{equation}
    \label{eq:unbelcasino}
    \hat{\mathbf{y}}=\boldsymbol{\Phi}\left(\boldsymbol{\Phi}^{\prime} \mathbf{W} \boldsymbol{\Phi}+\lambda \mathbf{R}\right)^{-1} \boldsymbol{\Phi}^{\prime} \mathbf{W} \mathbf{y}=\mathbf{S}_{\phi, \lambda} \mathbf{y}
\end{equation}

\subsection{Parameters Tuning} 
As we have seen in previous paragraphs, increasing the number of basis functions $K$ can increase the goodness of fit. Indeed, the larger $K$, the better the fit to the data. However, by increasing $K$ disproportionately, our model would also capture the stochastic component of the noise. On the other hand, by decreasing $K$ too much, we may risk missing some data features. It is the classic problem of the trade-off between bias and variance. The bias is defined as
\begin{equation}
    \label{eq:bias}
    \operatorname{Bias}[\hat{x}(t)]=x(t)-\mathrm{E}[\hat{x}(t)]
\end{equation}
and for large values of $K$, this quantity is small. On the other hand, for larger $K$, the variance
\begin{equation}
    \label{eq:variance}
    \operatorname{Var}[\hat{f}(t)]=\mathrm{E}\left[\{\hat{f}(t)-\mathrm{E}[\hat{f}(t)]\}^2\right]
\end{equation}
will be larger. We can write SSE from \ref{eq:sse} as 
\begin{equation}
    \label{eq:ssebias}
    \operatorname{SSE}[\hat{f}(t)]=\mathrm{E}\left[\{\hat{f}(t)-f(t)\}^2\right]] = \operatorname{Bias}^2[\hat{f}(t)]+\operatorname{Var}[\hat{f}(t)]
\end{equation}
Therefore, to find the correct number of bases $K$, it is sufficient to minimize the SSE, as shown in Fig. \ref{fig:tradeoff}.
\begin{figure}
\centering
\includegraphics[width=0.5\textwidth]{Images/tradeoff.png}
\caption[Bias variance trade-off.]{The heavy solid line indicates mean squared error totaled across the ages of observation between three and sixteen. The dashed line shows the totaled sampling variance, and the dotted-dashed line shows the totaled squared bias \cite{ramsay_functional_2006}.}
\label{fig:tradeoff}
\end{figure}
When we have to decide which parameter to use to penalize \ref{eq:pensse}, we can carry out the same operation as in the previous paragraph by minimizing generalized cross-validation defined as
\begin{equation}
\label{eq:gcv}
    \operatorname{GCV}(\lambda)=\left(\frac{n}{n-d f(\lambda)}\right)\left(\frac{\operatorname{SSE}}{n-d f(\lambda)}\right)
\end{equation}
where $df(\lambda)$ are the degrees of freedom value for the spline smooth and are computed as
\begin{equation}
   df(\lambda)=\operatorname{trace} \mathbf{S}_{\phi, \lambda}
\end{equation}
and $\mathbf{S}_{\phi, \lambda}$ is defined in \ref{eq:unbelcasino}.

\subsection{Functional Principal Component Analysis}
Functional principal component analysis (FPCA), which is nothing other than PCA applied to functional data, is one of the most used technique when dealing with functional data. Indeed, it is a common practice to look for features that characterize typical functions profile. For example, we might be interested in the sinusoidal nature of periodic data \cite{ramsay_functional_2006}. When we talked about classical PCA, the linear combination consisted of a vector of weights and the observations of the dataset. In the case of FPCA we have to translate this approach into the domain of functions, thus, into the domain of the "continuum". Equation \ref{eq:pcaf} becomes
\begin{equation}
    f_i= \int \beta y_i = \int \beta(t) y_i(t) dt
\end{equation}
and weights $\lambda$ become functions $\beta(t)$. 
\begin{figure}
    \centering
    \subfloat[\label{fig:mediafpca}]{
        \includegraphics[scale=0.2]{Images/canadamedie.png}
    }
    \qquad
    \subfloat[\label{fig:canadapc}]{
        \includegraphics[scale=0.28]{Images/PCs.png}
    }
    \caption[Functional Principal Component Analysis.]{Mean temperature curve of the 35 different stations over the considered time period (a) and the first four PCs with relative proportion of explained variance (b) \cite{ramsay_functional_2006, ramsay_functional_2009}.}
    \label{fig:FPCA}
\end{figure}
As in PCA, each weight function has the task of capturing the greatest amount of variance and being orthogonal to all other weight functions. Thus, every weight function will have a decreasing associated amount of explained variance. For a more in-depth explanation of the calculation method, the reader can refer to \citeauthor{ramsay_functional_2006} (2006). Consider the example in Fig. \ref{fig:FPCA}. The dataset consists of daily temperature data recorded at 35 locations in Canada from 1960 to 1994. Fig. \ref{fig:mediafpca} shows the 35 annual mean temperature curves. We can see that the first PC, although always strictly greater than zero, has values in the winter months up to four times higher than in the summer months. This means that temperatures in the winter tend to be much more variable than those in the summer months. Thus, weather stations with high scores will have hotter winters and warm summers. Indeed, the two highest scores are assigned to Vancouver and Victoria on the Pacific Coast. On the other hand, the largest negative score goes to Resolute in the High Arctic. The second PC shows similar overall behavior to the first PC but is not strictly greater than zero. This means it corresponds to a measure of temperature uniformity throughout the year. For example, the highest score goes to Prince Roupert, a city on the Pacific coast, characterized by a slight temperature difference between the summer and winter. Fig \ref{fig:canadapc} shows also the relative percentage of explained variance. The third PC helps to explain only 1.6\% of the total variance, so it is much less significant. Each PC has its logical interpretation, although as the relative variance explained decreases, it becomes more difficult to find a satisfactory one.
% <<< End of Functional Data Analysis
%%%%%
%%%%%
% Bagging Voronoi Clustering Algorithm >>>
\section{Bagging Voronoi Classifier}
\label{sec:bvc}
\citeauthor{secchi_bagging_2013} (2013) presents the Bagging Voronoi Clustering Algorithm (BV). It is an algorithm for unsupervised classification of functional data that exploits spatial dependence by repeatedly generating random connectivity maps and by clustering, at each replicate, local representatives of neighboring functional data. The algorithm is entirely non-parametric. Thus, we do not need any a priori data distribution hypothesis, making it more robust. The algorithm is structured in five consecutive steps. The first step involves using a Voronoi tessellation to partition the domain lattice. Voronoi diagram divides the plane into regions close to each object of a given set. In our case, these objects are just the sites of the lattice. We call these objects nuclei. Each nucleus has a corresponding region, called a Voronoi cell, consisting of all points of the plane closer to that nucleus than to any other. Let $X$ be a metric space and $d(\cdot, \cdot)$ be a distance function, in our case, Euclidean distance. Let $\Phi_n$ be the set of $n$ selected nuclei and let each nucleus be a tuple of coordinates $\left(P_k\right)_{k\in \Phi_n}$. The formal definition of a Voronoi cell $R_k$ associated with the site $P_k$ is 
\begin{equation}
    \label{eq:voronoicell}
    R_k=\{x\in X \mid d(x, P_k)\leq d(x, P_j),\forall j\neq k\}
\end{equation}
The Voronoi tessellation is the collection of all the cells in the space.
Fig. \ref{fig:voronoi} shows an example.
\begin{figure}
    \centering
    \includegraphics[width=0.3\textwidth]{Images/A-set-of-atoms-the-associated-Voronoi-tessellation-solid-lines-and-the-Delaunay.png}
    \caption[Voronoi tessellation.]{Example of Voronoi tessellation with highlighted nuclei.}
    \label{fig:voronoi}
\end{figure}
\subsection{Algorithm Steps}
\label{subsec:algsteps}
This section describes in detail the algorithm steps. In Alg. \ref{alg:bvc}, there is the pseudo-code of the BV algorithm. Suppose a latent field of labels $\Lambda_0:\mathcal{S}_0 \rightarrow \{1,\dots,L\}$ which is defined on the lattice $\mathcal{S}_0$. $\Lambda_0(\mathbf{x})$ is the true unknown label associated to the site $\mathbf{x}\in\mathcal{S}_0$, where $\mathcal{S}_0 \subset \mathcal{S}$ and $\mathcal{S}$ is a measurable subset of $\mathbb{R}^d$. In addition, suppose that a functional datum is observed in each site $\mathbf{x}\in\mathcal{S_0}$. The steps needed to find an estimation of true label $\Lambda_0(\mathbf{x})$ the following steps are needed.
\begin{algorithm}
\scriptsize
    \caption{\footnotesize{Bagging Voronoi Classifiers}}
    \label{alg:bvc}
    \begin{algorithmic}
    \STATE \textbf{Bootstrap:}
    \STATE Initialize $B, n, p, K$. Choose a metric $d(\cdot,\cdot).$
    \FOR{$b:=1$ to $B$}
    \STATE{Randomly generate a set of nuclei $\Phi^b_n=\{\mathbf{Z}_1^b, \dots, \mathbf{Z}_n^b\}$ among the sites in $\mathcal{S}_0$}
    \FOR{$i=1$ to $n$}
    \STATE $\mathbf{Z}_i^b\sim\mathcal{U}(\mathcal{S}_0)$, where $\mathcal{U}$ is the uniform distribution on the lattice. Obtain a random Voronoi tessellation of $\mathcal{S}_0, \{V(\mathbf{Z}_i^b|\Phi_n^b\}^n_{i=1}$ by assigning each site $\mathbf{x}\in \mathcal{S}_0$ to the nearest nucleus $\mathbf{Z}_i^b$, according to the specified distance $d(\cdot,\cdot).$
    \ENDFOR
    \FOR {$i:=1$ to $n$}
    \STATE Compute the function $g_i^b$, acting as local representative, by summarizing information carried by the functional data associated to sites belonging to the $i$-th element of the tessellation $V(\mathbf{Z}_i^b|\Phi_n^b)$.
    \ENDFOR
    \STATE Perform dimensional reduction of the local representatives $\{g_1^b,\dots, g_n^b\}$ by projecting them on the space spanned by a proper $p$-dimensional scores vectors $\{\mathbf{g}_1^b,\dots, \mathbf{g}_n^b\}$, which are then clustered in $K$ groups according to suitable unsupervised method.
    \ENDFOR
    \STATE \textbf{Aggregation:} perform cluster matching
    \FOR {$k:=1$ to $K$}
    \FOR {$b:=1$ to $B$}
    \STATE indicate with $C_k^b$ the set of $\mathbf{x} \in \mathcal{S}_0$ whose label is equal to $k$, and match the cluster labels across bootstrap replicates, to ensure identifiability.
    \ENDFOR
    \ENDFOR
    \FOR {$\mathbf{x} \in \mathcal{S}_0$}
    \STATE Calculate the frequencies of assignment of the site to each of the $K$ clusters along iterations, \textit{i.e.}, $\pi_{\mathbf{x}}^k=\#\{b\in\{1,\dots,B\}:\mathbf{x}\in C_k^b\}/B, \forall k=1,\dots,K$
    \ENDFOR
    \end{algorithmic}
\end{algorithm} 
\paragraph{Step 1: Voronoi Tessellation.} Select a set of points $\mathbf{x}\in\mathcal{S}$ as nuclei for the Voronoi tessellation. Thus, let $\Phi_n=\{\mathbf{Z}_1, \dots, \mathbf{Z}_n\}$ be the set of $n$ points in $\mathcal{S}$ sampled from a proper distribution $F$ defined on S: this will be the set of nuclei of the Voronoi tessellation. For each $\mathbf{Z}_i\in\Phi_n$ define the polyhedron
\begin{equation}
    \label{eq:polyedron}
    V\left(\mathbf{Z}_i \mid \Phi_n\right)=\{\mathbf{x}\in\mathcal{S} : d\left(\mathbf{x},\mathbf{Z}_i\right) \leq d\left(\mathbf{x}, \mathbf{Z}_j\right), \quad \text{for all}\quad \mathbf{Z}_j \in \Phi_n, i\neq j\}
\end{equation}
to be the closest Voronoi cell with nucleus $Z_j$ for the Voronoi tessellation induced by $\Phi_n$.
\paragraph{Step 2: Functional Local Representatives.} Consider $T$, a bounded interval of $\mathbb{R}$ and a realization $f_{\mathbf{x}}:T\rightarrow \mathbb{R}$ of a functional random variable is observed in each site $\mathbf{x}$ of the lattice $S_0$. The Voronoi tessellation of the lattice $S_0$ computed in the previous step defines random neighborhoods. For each element $V_i$ of the Voronoi tessellation, we sum up the information contained in the sub-sample $\{f_{\mathbf{x}}\}_{\mathbf{x} \in V_i}$ by estimating a functional local representative through a method that exploits spatial dependence of neighboring data. For example, the authors suggest using a weighted mean with a Gaussian kernel. Still, we can use any functional local representative we prefer, depending on our needs and problem domain. For $i=1, \dots, n$ the functional local representative $g_i$ aforementioned is defined as
\begin{equation}
    \label{eq:gaussianmean}
    g_i(t)=\frac{\sum_{\mathbf{x}\in V_i}w^i_{\mathbf{x}}\cdot f_{\mathbf{x}}(t)}{\sum_{\mathbf{x}\in V_i}w^i_{\mathbf{x}}}
\end{equation}
where $w_{\mathbf{x}}^i$ is a Gaussian weight centered in $Z_i$ and decreasing with respect to $d\left(\mathbf{x}, \mathbf{Z}_i\right)$. Intuitively, we assume that spatial dependence between two sites decreases when the distance between them increases. The kernel covariance matrix is $\sigma^2\mathbb{I}_2$, where $\sigma=d_{max}/d_{min}$, being $d_{max}$ and $d_{min}$ the maximum and minimum distance between two nuclei of the tessellation, respectively. This choice connects $\sigma$ to the mean dimension of the tessellation via an estimator of the element's mean diameter. The choice of $n$, which sets the tessellation dimension and, thus, the number of local representatives to be computed, greatly influences the algorithm behavior. In general, we have to find an optimal $n$ to find a good compromise between variance and bias. We will see two possible approaches for choosing the optimal $n$ value in Section \ref{sec:irradiance}.
\paragraph{Step 3: Dimensional Reduction and Clustering.} The algorithm's third step aims to perform data dimensional reduction and cluster the dimensionally reduced data. First, we must find the base explosion that best approximates the observed profiles. This depends mainly on the domain of the phenomenon, or, more quickly, we can use a Gaussian smoother. This operation is something that can be done a priori. For dimensional reduction of functional data, we need to find the best projection local representatives $\{g_1, \dots, g_n \}$ onto the space generated by a proper functional basis of finite dimension $p$. To choose the correct dimension $p$, we can use the amount of variance explained by each functional principal component. For instance, by setting a threshold of the relative amount of total variance explained, keep the first $p$ principal components that account for a variance greater than or equal to this threshold. Alternatively, we could use the scree plot of the absolute variance explained by each principal component and retain the $p$ components before the knee. This latter approach offers more freedom of choice to the analyst. Thus, a deep understanding of the problem domain becomes essential. Using a bootstrapping method, typical of bagging clustering, we can repeat steps 1-3 $B$ times, where $B$ is arbitrarily chosen. 
\paragraph{Step 4: Cluster Matching.} The output of the $b$-th replicate of the bootstrap phase of the algorithm is a label assignment for each local representative estimated during the $b$-th run. All sites $\mathbf{x} \in V_i^b$ get the same label associated with the function local representative $g_i^b(t)$. To obtain a classification map of the lattice $\mathcal{S}_0$, we consider the frequency distribution of assignment of each site to each of the $K$ clusters along the $B$ replicates. To compute this frequency distribution, we need, in turn, to assume that cluster labels $\{C_1^b, \dots, C_K^b\}$ are coherent along the $B$ replicates. More specifically, we want cluster labels $\{C_1^b, \dots, C_K^b\}$ to be coherent with $\{C_1^m\, \dots, C_K^m\}$, for all $m<b$ and $b\geq2$. This is obtained through cluster matching. Cluster matching is a set of techniques that ensures that different clustering algorithm outputs are coherent. A simple approach is based on a contingency table. The algorithm looks for the label permutation that minimizes the total sum of the off-diagonal frequencies in the contingency table describing the joint distribution of sites along the classifications at two subsequent replicates. At the end of clustering matching, for each site $\mathbf{x}$ we will have a vector $\pi_{\mathbf{x}}^K=\left(P_1, \dots, P_K\right)$ where each $P_i, \quad i=1,\dots,k$ is the frequency the label $k$ was assigned to the site $\mathbf{x}$. The final label will be the label associated with maximum $P_k$. In other words, it is a majority voting process among the different replicates $b$.
\paragraph{Step 5: Output Evaluation} We can use spatial entropy to have an overview of the statistical validity of the performed clustering. \emph{Spatial entropy} is directly derived from the classical notion of entropy in statistics. Considering the frequency distribution vector $\boldsymbol{\pi}=(\pi_{\mathbf{x}}^1,\dots,\pi_{\mathbf{x}})^K$ of each site $\mathbf{x} \in \mathcal{S}_0$ to each of the $K$ clusters obtained after the bootstrapping step of the algorithm. The entropy associated to site $\mathbf{x}$ classification in the site $\mathbf{x}\in\mathcal{S}_0$ is defined as
\begin{equation}
\mathbf{\eta}_{\mathbf{x}}=-\sum_{k=1}^K\pi_{\mathbf{x}^k}\cdot \log(\pi_{\mathbf{x}^k})
\end{equation}
which assumes the minimum value 0 when there is an $r$ such that $\pi_{\mathbf{x}}^r=1$ while $\pi_{\mathbf{x}}^k=0$ for all $k\neq r$ and maximum value $\log(K)$ when $\pi_{\mathbf{x}}^k=\frac{1}{K}$ for $k\in\{1,\dots,K\}$. We can see spatial entropy as a measure of the disorder of the frequency of site assignment to cluster $k$. Moreover, we can compute the average normalized entropy to have a global evaluation index of the clustering performed. \emph{Average normalized} entropy is defined as
\begin{equation}
\label{eq:avgentropy}
    \eta^K=\frac{\sum_{\textbf{x}\in\mathcal{S}_0}\eta_{\mathbf{x}}^K}{\log(K)\cdot |\mathcal{S}_0|}
\end{equation}
In \ref{eq:avgentropy}, the index has been normalized with respect to maximum value $\log(K)$ to allow comparison over different choices of $K$. Another possible index for evaluating a classification in the ratio
\begin{equation}
    \label{eq:theta}
    \theta=\frac{tr(B)}{tr(B)+tr(W)}
\end{equation}
where $B$ and $W$ are the final between and within-cluster sum of the squares matrix, respectively. This index can be interpreted as the percentage of variance explained by splitting the observations into clusters.
% <<< End of Bagging Voronoi Clustering Algorithm
%%%%%
%%%%
% Succesful case of the algorithm >>>
\section{Clustering of Irradiance Data}
\label{sec:irradiance}
In this section, we will discuss an example of how the algorithm outlined in the previous paragraph has been successfully applied in the clustering of irradiance data performed in \citeauthor{secchi_bagging_2013} (2013). This section addresses any potential questions or uncertainties the reader may have concerning using Algorithm \ref{alg:bvc}. 

Insolation refers to the amount of solar radiation energy captured on a specific surface area over a defined time period. Typically, it is computed as average irradiance in kilowatt-hours per square meter per day (\unit{\kilo\watt.h/\metre^2.day}). The authors focused on direct insolation, which refers to the solar irradiance measured at a particular location on Earth with a surface facing directly toward the sun's rays, omitting diffuse insolation. Diffuse insolation refers to the solar energy scattered or reflected by the atmosphere or the cloud. The value of direct insolation is computed as solar constant minus the atmospheric losses. While the solar constant is influenced by factors like Earth's distance from the sun and solar cycles, the losses depend on factors such as the time of day (which affects the light's path through the atmosphere based on the sun's elevation angle), cloud coverage, atmospheric moisture, and other factors. More specifically, they decided to analyze the maximum deficit under the average incident solar radiation on a horizontal surface over a consecutive day. This quantity can be seen as a proxy for the extra energy capacity required to compensate for the lack of energy produced by solar panels installed in a certain area. Available data for the analysis consists of vectors in $\mathbb{R}^{12}$ indexed by each site of the lattice. Sites are located on a non-uniform lattice $\mathcal{S}_0=\bigcup_{\lambda \in Z_{\mathbf{1}};\theta \in Z_{\mathbf{2}}}A_{\lambda\theta}$ where $Z_{\mathbf{1}} = \{ -180, -179, \dots, 178, 179\}$ and $Z_{\mathbf{2}} = \{ -66, -65, \dots, 65\}$: each element $A_{\lambda\theta}$ is the portion of the earth surface which is included between the meridians at longitude $\lambda$ and $\lambda + 1$ in degrees, and between the parallels at latitude $\theta$ and $\theta + 1$ in degrees. This lattice includes 47,520 worldwide non-polar districts. In each site, the 12 measures correspond to the monthly maximum energy deficit values with respect to the monthly average. The maximum and average values are computed over the 22 years from July 1983 to June 2005. For each site, authors obtained a functional datum $Y_{\lambda,\theta}(t)$ by smoothing $\{Y_{\lambda\theta}^1, \dots, Y_{\lambda\theta}^12\}$ with a Gaussian kernel with bandwidth equal to 1.5. These functional profiles are the input for the BV algorithm. The first $p=3$ functional principal components were retained since the proportion of explained variance was more than 95\% with only these three components. For classifying the scores of the $n$ representatives, they employed the K-Means algorithm with the $L^2$ semi-metric induced by the principal components.
\begin{figure}
    \centering
    \includegraphics[width=0.6\textwidth]{Images/baggingscreeplot.png}
    \caption[Result on irradiance data.]{Algorithm performances with irradiance data. On the left is the average normalized entropy; on the right is the index $\theta$ \cite{secchi_bagging_2013}.}
    \label{fig:baggingscreeplot}
\end{figure}
Subsequently, they performed the clustering multiple times with a fixed $B=100$ to obtain the scree plot shown in Fig. \ref{fig:baggingscreeplot} and to select the optimal number of clusters $K$. The same operation was necessary to determine the optimal number of nuclei, $n$, for the initial Voronoi tessellation. Fig. \ref{fig:baggingscreeplot} shows the scree-plots. The optimal choice for these two parameters was made using average normalized spatial entropy, the index $\theta$, and the normalized spatial entropy map. Starting from the scree-plot of average normalized entropy, we can see that for $n=500$, we have two local minima, specifically for $K=3$ and $K=5$. The plot of parameter $\theta$ suggests that the number of clusters $K=6$ appears to be the most reasonable. Indeed, using $K=7$ would not be a good choice, as it would not enhance the obtained clusterization. Consequently, the classification should yield good results for $K=5$ and $K=6$. This was also confirmed by the maps of the normalized entropy, shown in Fig. \ref{fig:irrentropy}.
\begin{figure}
    \centering
    \includegraphics[width = 0.8\textwidth]{Images/irrentropy.png}
    \caption[Normalized spatial entropy for irradiance data.]{Normalized spatial entropy maps associated to the classification with K = 5 (left) and K = 6 (right). Colors from red to white correspond to values from 0 to 1; higher values identify areas where classification is more or less uncertain. \cite{secchi_bagging_2013}.}
    \label{fig:irrentropy}
\end{figure}
The final choice $K=5$ was made looking at maps in Fig. \ref{fig:irrentropy}, which tell us that for $K=5$, the classification is more robust, and from the fact that with $K=5$, the algorithm identified different homogeneous macro-areas which seem interpretable in terms of the observed phenomenon. In Fig. \ref{fig:resbagging}, there is the final classification for irradiance data. The two classifications are not in contrast but somehow support each other. As mentioned earlier, we can take advantage of the functional nature of the data to derive additional insights that would not be possible with traditional clustering algorithms. For instance, the red cluster exhibits a non-seasonal pattern that remains consistent throughout the year, and the green and yellow clusters have opposite seasonal patterns. From an engineering point of view, the red cluster is the one that shows an annual buffer capacity pattern, which is optimal in terms of electricity production by solar power because it needs minimal buffer capacity installation. %The advantage of using $K=6$ is that we can identify some sites with a very high seasonality, highlighted with the orange cluster.
To conclude, it is noteworthy that spatial entropy is systematically higher near the boundary between two clusters. This is because the functional data is spatially dependent, so data near the boundaries are more challenging to cluster.
\begin{figure}
    \centering
    \includegraphics[scale=0.3]{Images/resbagging.png}
    \caption[Classification of irradiance data.]{In the top panels, final classification maps are obtained via a majority vote on frequencies of assignment and by setting $K = 5$ (left) or $K = 6$ (right).In the bottom panels, a set of functional local representatives obtained with $n = 500$ in one of the iterations of the algorithm and clustered with $K = 5$ (left) and $K=6$ (right). From \citeauthor{secchi_bagging_2013} \citeyear{secchi_bagging_2013}.}
    \label{fig:resbagging}
\end{figure}

%\subsection{Analysis of Spatio-Temporal Mobile Phone Data in Milan}
% \label{subsec:mobilephonemilan}
% I have chosen to present this example to illustrate to the reader the flexibility of the previously discussed algorithm. It demonstrates how the method can adapt to different data dimensionality reduction techniques and how a Fourier basis can represent functional data exhibiting certain seasonality, as with lattice structures in Section. \ref{subsec:devcell}.
% In \citeauthor{secchi_analysis_2015} \citeyear{secchi_analysis_2015}, authors examined high-dimensional, geo-referenced mobile phone usage data from Milan's urban area to discern activity patterns within the city over time. The goal was to identify specific urban zones with similar temporal patterns, which could indicate specific activities or events in those areas. To achieve the goal, they decided to use the algorithm explained in Section \ref{sec:bvc} coupled with treelet analysis (I'll explain it in just a moment). 
% The data available for analysis are from Telecom Italia. In this dataset, the metropolitan city of Milan is divided into a uniform lattice $\mathcal{S}_0$ consisting of 97x109 sites. In each location, the average number of mobile phones simultaneously using the network for calling is provided every 15 minutes over 14 days.

% %%%%% Grammarly eseguito sul testo qua sopra ^^^


% This quantity is called Erlang and it is defined as
% \begin{equation}
%     \label{eq:erlang}
%     E_{\mathbf{x}j}=\frac{1}{15}\sum_{q=1}^Q \left| T_{\mathbf{x}j}^q \right|
% \end{equation}
% where $T_{\mathbf{x}j}^q$ indicates the time interval in which the $q$-th mobile phone is using the network for calling while moving within site $\mathbf{x}$ and during the $j$-th quarter of an hour. The Erlang data was recorded every quarter of an hour, from March 18th, 2009, 00:15, till March 31st, 2009, 23:45. Fig. \ref{fig:secchino} there is an example of an Erlang profile.
% \begin{figure}[H]
%     \centering
%     \includegraphics[scale=0.43]{Images/secchierlang.png}
%     \caption[Erlang profile.]{Erlang profile for a site in the lattice $\mathcal{S}_0$ \cite{secchi_analysis_2015}.}
%     \label{fig:secchino}
% \end{figure}
% However, for some sites of the lattice, Erlang profiles had some missing values, and in some instances, they were entirely missing. Therefore, after the data cleaning phase, the authors decided to focus only on a non-uniform time grid of $p = 1,308$ elements, each element of the time grid being relative to a quarter of an hour for which an Erlang measurement has been observed in at least one site of the lattice. In Algorithm \ref{alg:bvc} for dimensional reduction of functional data, FPCA was suggested. While FPCA stands as an effective technique for identifying optimal subspaces to represent data, it has a limitation in its global approach, making it less appropriate for multi-resolution analysis. This is primarily because its basis elements typically involve a linear combination of all foundational variables. So, authors decided to use a treelet basis, introduced in \citeauthor{lee_treeletsadaptive_2008} (2008). For this specific case, since the data exhibits an evident seasonality, authors decided to use a Fourier basis with a period equal to 1 week. Thus, the reconstructed functional form of the Erlang profile for site $\textbf{x} \in \mathcal{S}_0$ is a function $E_{\textbf{x}(t)}$ such that:
% \begin{equation}
%     E_{\textbf{x}}(t)=\frac{c_0^{\textbf{x}}}{2} + \sum_{h=1}^H \left[a_h^{\textbf{x}}\cos \left( h\omega t\right) + b_h^{\textbf{x}}\cos \left( h\omega t\right) \right]
% \end{equation}
% where $t \in \left[ 0; T\right]$, $\omega = 2\pi/T$ and $T=60 \times 24 \times 7$ is the period expressed in minutes. To choose the optimal basis dimension $H$, we can use the power spectrum associated to the site~-~wise smoothing of the Erlang data with a Fourier basis of large dimension $H = 200$. The power spectrum of the Fourier expansion of a signal represents the amplitude of the signal as a function of the frequency, and at the $h$-th frequency it is related to the amplitude of the $h$-th harmonic

% \begin{equation}
%     \label{eq:spectrum}
%     P_{\mathbf{x}}(h)=\sqrt{(a^{\mathbf{x}}_h)^2 + (a^{\mathbf{x}}_h)^2}
% \end{equation}

% Thus, the more the $h$-th harmonic is relevant in the explanation of features occurring in the data, the more $P_{\mathbf{x}}(h)$ will be large. From the spectrum in Fig. \ref{fig:fantasmino}, we can see that local maximum happening at multiple of 7 (the dotted lines), while for $H >100\omega$, power spectre becomes almost negligible.

% \begin{figure}[H]
%     \centering
%     \includegraphics[scale=0.6]{Images/fantasmino.png}
%     \caption[Spectrum of Erlang data.]{Average power spectrum $P(h)$ obtained via site-wise smoothing of the Erlang measures with a Fourier basis of dimension $H = 200$. Only the values of $P(h)$ for$ h = 1, \dots, 100$ are shown in the plot. Dotted vertical lines are drawn for multiples of 7 \cite{secchi_analysis_2015}.}
%     \label{fig:fantasmino}
% \end{figure}

% As did in Section , the metric $d(\cdot, \cdot)$ used is Euclidean distance, the total number of bootstrap $B=50$, and the dimension $n$ of the Voronoi tessellation ranging from 50 to 1250. As we already know, the optimal $n$ will be found by using the scree plot.


% \subsection{One Last Thing}
% \label{subsec:devcell}
% In previous section, I have chosen to provide an example of how functional data and clustering algorithms can be used to conduct a layer-wise temperature profile analysis. However, these concepts can be readily adapted to other scenarios, including, for instance, particular application in three dimensions. In this section, I want to go beyond the primary domain of the thesis. I want to suggest how the algorithm could also be leveraged for the classification of functional data in three dimensions. A possible application in three dimensions could be the clustering of functional data regarding the deviation of cell profiles in lattice structures. Recall that in 3 dimension, Voronoi regions are represented as polyhedra, like in Fig. \ref{fig:3dvoronoi}.

% \begin{figure}[H]
%     \centering
%     \includegraphics[scale=0.43]{Images/3D-Voronoi-tessellation-with-1000-grains-unit-cell-with-its-internal-grains-not.png}
%     \caption[3d Voronoi tessellation.]{An example of a Voronoi tesselation in 3d.}
%     \label{fig:3dvoronoi}
% \end{figure}

% In \citeauthor{colosimo_complex_2022} (2022), the authors aim to compare the nominal slices with the actual slices measured through x-ray computed tomography of a lattice structure composed of 64 dodecahedron unit cells of size \SI{10}{\milli\metre} within a specimen of $40 \times 40 \times 40$ \unit{\milli\metre}. Subsequntely, the 2D image of the as-built single cell and as-designed one were compared, slice by slice, in order to compute the deviation index. Given a pair of images that correspond to the $z-$th slice of the as-designed cell and the $z-$th slice of the as-built cell, let $i_{u,v,z}(\text{as.designed}$ and let $i_{u,v,z}(\text{as.build}$ be the intensities of $(u,v)^{th}$ pixel in the as-designed slice and in the as-built slice, respectively. In Fig. \ref{fig:slices} there are some examples of the procedure. A superimposition of these two images leads to four possible cases:
% \begin{itemize}
%     \item Region 1: pixels in which both $i_{u,v,z}(\text{as.designed})=0$ and $i_{u,v,z}(\text{as.build})=0$, i.e., pixels where material is present in both the images.
%     \item Region 2: pixels in which both $i_{u,v,z}(\text{as.designed})=1$ and $i_{u,v,z}(\text{as.build})=1$, i.e., pixels that correspond to the background in both the images.
%     \item Region 3: pixels in which both $i_{u,v,z}(\text{as.designed})=0$ and $i_{u,v,z}(\text{as.build})=1$, i.e., pixels where material is present in the as-designed slice but not in the as-built slice (i.e., less material has been produced than indicated in the CAD model)
%     \item Region 4: pixels in which both $i_{u,v,z}(\text{as.designed})=1$ and $i_{u,v,z}(\text{as.build})=0$, i.e. pixels where material is present in the as-built slice but not in the as-designed slice (i.e., more material has been produced than indicated in to the CAD model).
% \end{itemize}


% \begin{figure}
%     \centering
%     \includegraphics[scale=0.5]{Images/slicescolosimo.png}
%     \caption[Deviation index.]{Example of as-built (top row) and as-designed (central row) slice images at different z heights with the corresponding deviation (bottom row) \cite{colosimo_complex_2022}. }
%     \label{fig:slices}
% \end{figure}


% The overall number of pixels that belong to the union of region 3 or region 4 is the number of pixels for which $i_{u,v,z}(\text{as.build}) - i_{u,v,z}(\text{as.designed}) \neq 0$. So the \emph{deviation index} for the cell $(i,j,k)$ is defined as
% \begin{equation}
% \label{eq:deviationindex}
% \delta_{i, j, k}(z)= \sum_{u=1}^p \sum_{v=1}^p \mathcal{I}\left(i_{u, v, z}(\text { as.built })\right. -i_{u, v, z}(\text { as.designed })\neq 0)_{i, j, k},
% \end{equation}
% where $\mathcal{I}$ is the indicator function, that is 1 if the condition in brackets is true and 0 otherwise.
% After calculating the deviation index for each cell, the authors suggest using a cubic B-splines basis to fit a functional profile on the deviation index as a function of the slice position along cell z-axis. The position of the knots was chosen to eliminate the first and second discontinuity present in the deviation profiles, resulting from the typical trabecular structure of lattice designs. Fixed knots were then placed at the discontinuities of the as-designed geometry. Subsequently, intermediate knots were added iteratively, until a knee in the mean squared error (MSE) scree plot of the B-spline model residuals was identified.
% At the conclusion of step 5 in Fig. \ref{fig:colosimofunctional}, as outlined in the article, we will have 64 functional profiles, one for each cell. These profiles and the centroid of each cell in a lattice in $\mathbb{R}^3$ and will be the input for the classification algorithm described in Section \ref{sec:bvc}. In doing so, we could identify if there are critical areas within the lattice where cells exhibit similar anomalous behavior. Coupled with the ability to quantify the magnitude of the defect using the functional profiles, this could lead to a deeper understanding of the lattice structure's behavior and determine whether, despite its imperfections, it remains fit for its intended use. Indeed, localized geometric and dimensional deviations of the unit cells might adversely affect the structure's elastic modulus and compressive strength \cite{colosimo_complex_2022}.

% \begin{figure}
%     \centering
%     \includegraphics[scale=0.5]{Images/colosimosplines.png}
%     \caption[Functional control chart in complex geometries.]{Schema of the proposed methodology in \cite{colosimo_complex_2022}. The output of step 5 would be the input for the classification algorithm.}
%     \label{fig:colosimofunctional}
% \end{figure}
% % <<< End of Succesful case of the algorithm

% ^^^^^^^^^^^^^^^^^^^^^^^^^^^^^^^^^^^^^^^^^^
% Fatta revisione
%%%%%%%%%%%%%%%%%%%%%%%%%%%%%%%%%%%%%%%%%%%%


%>>> Simulation Study
%%%%%%%%%%%%%%
% REVIEW FATTA
%vvvvvvvvvvvvv
\section{Simulation Study}
\label{sec:simstudy}
In this section I want to take a little extra step on my thesis aim. As additional food for thought, this section briefly presents an example of the algorithm applied on simulated data. The purpose is to test in a controlled environment whether the algorithm could work and then move on to the real data collection phase. In addition, the experiment also wants to test the possibility of applying the algorithm described in the previous chapter in real-time context. The simulated thermal profiles were obtained using the approach proposed in \cite{colosimo_spatially_2018} (2018) and \citeauthor{yan_real-time_2021} (2021). We can describe the intensity of each pixel in a spatial location $s$ at a certain frame $t$ with a proper sigmoid function
\begin{equation}
    \label{eq:simulation}
    a_{s, t}=\frac{255}{1+\exp (0.2(t-H \tau))}, t=1, \ldots, \tau
\end{equation}
the parameter $\tau$ can be used to make the HS persistent $\tau$ frames, and the parameter $H$ is a constant that controls the shape of the cooling profile in the hot-spot regions. In general, $\tau \in \left(0,180\right)$ and $H=0.95$. For $\tau=0$, we have an in-control profile. Fig. \ref{fig:simulationprofile} shows examples of simulated thermal profiles according to different parameters.
\begin{figure}
    \centering
    \includegraphics[width=0.6\textwidth]{Images/simulation.png}
   \caption[Simulated profiles.]{Different thermal profiles obtained by varying $\tau$ and $H$ parameters.}
    \label{fig:simulationprofile}
\end{figure}
A \numproduct{50 x 50} pixels grid was used during the experiment. As suggested in previous papers, cross-shaped clusters of pixels of varying sizes were inserted to simulate different HS sizes: small ($n=9$ pixels), medium ($n=21$ pixels), and large ($n=45$ pixels). In each scenario, two medium size HSs were injected with different $\tau$. The persistence was randomly selected in the range $\tau\in (1,51)$. The HSs were inserted in random locations for. In addition, a random component $\varepsilon \sim \mathcal{N}(0, 1)$ was also added to each profile to simulate the noise present in real production settings.
For the experiment, paramenters $n$ and $B$ were varied, while the first $p$ PCs were kept contributing to obtain a variance of the explained percentages of 95\%. Since we are evaluating the general performance of the algorithm and not its power in defect detection, the number of clusters was superimposed to $k=4$. The experiment was performed on a Thinkpad T14 with an AMD Ryzen Pro 5 processor and 16 GB of RAM. Fig. \ref{fig:time} shows the computational time to classify a frame according to different value parameters.
\begin{figure}
    \centering
    \includegraphics[width = 0.5\textwidth]{Images/time.png}
    \caption[Computational time of BVC on simulated data.]{Computational time needed to process a single frame according to different algorithm parameters. The plot was cut to $n=70$ for clarity.}
    \label{fig:time}
\end{figure}
It seems obvious that it cannot be used for real-time applications. Presumably, implementing the algorithm in C or C++ would greatly improve performances mainly for two reasons: (i) the functions that generate random numbers in python are inefficient, and (ii) the low-level functions of C and C++ would allow the algorithm to be implemented by exploiting all processor cores by parallelizing bootstrap operations. Indeed, the algorithm by its very nature is particularly suited to be programmed using multi-threading logic. Fig. \ref{fig:frames} shows some frames of the simulated thermal profiles.
\begin{figure}
    \centering
    \includegraphics[width = 0.6\textwidth]{Images/frames.png}
    \caption[Frames of simulated data.]{Frames from simulated data of the second experiment.}
    \label{fig:frames}
\end{figure}
To test the ability to correctly cluster observations, the correct classification rates defined as
\begin{equation}
    \label{eq:classrate}
    \operatorname{CLASS} = \frac{\# correct classified}{\# Total}
\end{equation}

The problem with \ref{eq:classrate} is that it is strongly biased in this case: since the defects are very small relative to the domain, the large number of pixels in control will greatly affect the $\operatorname{CLASS}$.
Therefore, to overcome this problem, it was decided to use the classification rate only on OOC pixels.
Table \ref{tab:result} reports results for different parameter configurations. The standard deviation is shown in parentheses.

\begin{table}[]
\begin{tabular}{|llll|}
\hline
\rowcolor{bluepoli!40}\textbf{n} & \textbf{B} & \textbf{Class} & \textbf{Entropy} \\ \hline
30         & 10         & 0.083(0.06)    & 0.108(0.033)     \\
           & 40         & 0.083(0.06)    & 0.113(0.036)     \\
           & 100        & 0.042(0.11)    & 0.125(0.045)     \\
50         & 10         & 0.250(0.034)   & 0.077(0.026)     \\
           & 40         & 0.250(0.038)   & 0.081(0.033)     \\
           & 100        & 0.294(0.05)    & 0.094(0.028)     \\
70         & 10         & 0.354(0.042)   & 0.047(0.019)     \\
           & 40         & 0.332(0.054)   & 0.066(0.017)     \\
           & 100        & 0.362(0.06)    & 0.067(0.016)     \\
300        & 10         & 0.431(0.054)   & 0.027(0.005)     \\
           & 40         & 0.444(0.06)    & 0.029(0.005)     \\
           & 100        & 0.431(0.098)   & 0.033(0.007)     \\ \hline
\end{tabular}
\caption{$\operatorname{CLASS}$ and normalized entropy for different scenarios. Standard deviation is shown in parentheses.}
\label{tab:result}
\end{table}
From the table, we can see that as $n$ changes, the $\operatorname{CLASS}$ changes considerably. This is because small values of $n$ tend to "hide" smaller defects. On the other hand, we can see that increasing $B$ does not significantly impact the mean, but slightly reduces the standard deviation.Thus, if we want to get a good $\operatorname{CLASS}$ we need to increase $n$, in a sense $n$ is a measure of how accurate we want to be in capturing even small defects, whereas $B$ is a measure of how sure we want to be about clustering. The results obtained are not excessively good, probably because $n$ is too small relative to the domain. Indeed, for $n=600$ and $B=100$ we begin to get acceptable results. Fig. \ref{fig:risultati} shows the result of frame classification in Fig. \ref{fig:frames}.
\begin{figure}
    \centering
    \includegraphics[width = 0.7\textwidth]{Images/final_figure_1.png}
    \caption[Bagging Voronoi Classificator maps.]{Cluster map and normalized entropy map obtained with $n=600$ and $B=30$.}
    \label{fig:risultati}
\end{figure}
% <<< End of Simulation Study
\backgroundsetup{contents={}}
% <<< End of proposed method Bagging Voronoi


% Conclusion and future development
\chapter{Conclusions}
\label{ch:conclusions}
\backgroundsetup{contents={DRAFT}}
\begin{minipage}{0.7\textwidth}
\small
			\textit{We can imagine that this complicated array of moving things which constitutes “the world” is something like a great chess game being played by the gods, and we are observers of the game. We do not know what the rules of the game are; all we are allowed to do is to watch the playing. Of course, if we watch long enough, we may eventually catch on to a few of the rules.
            \\[1.5ex]
            -Richard Feynman}
\end{minipage}
\\[2.5ex]
\textcolor{red}{\textbf{To be completed.}}
%This chapter is the end of our journey into additive manufacturing for metallic components and functional data world. In \emph{Chapter ~\ref{ch:baggingvoronoi}} were just some demonstrations of how the classification algorithm could be employed with functional data to detect hot spots. Fundamentally, the strength of the algorithm lies in its ability to identify spatial dependencies and evolution in time or space, since in additive manufacturing time and space along z-axis are strongly linked. Such functional data afford us a deep understanding of the events unfolding in space/time, also allowing us to understand how the printing process evolve in time. Finally, the algorithm is highly versatile, allowing us to select the size of the lattice site and, consequently, the "magnification lens" through which we wish to approach the analysis. Indeed, we can use given its bootstrapping approach, it also fully leverages the capabilities of today's multi-core processors. But not just temperature profiles: functional data benefits from all the properties functions dos. Other examples of functional data could include electrons emitted during fusion in EBM processes over time, or the concentration of inert gas flowing through the printing chamber. The purpose of this thesis, in addition to providing an overview of AM processes for metals, was to reflect on how a relatively young branch of mathematics (late 90s) can contribute to the understanding of physical phenomena in AM.
\backgroundsetup{contents={}}
% <<< End of conclusion

%-------------------------------------------------------------------------
%	BIBLIOGRAPHY
%-------------------------------------------------------------------------

\addtocontents{toc}{\vspace{2em}} % Add a gap in the Contents, for aesthetics
\bibliography{references}

%-------------------------------------------------------------------------
%	APPENDICES
%-------------------------------------------------------------------------

\cleardoublepage
\addtocontents{toc}{\vspace{2em}} % Add a gap in the Contents, for aesthetics
\appendix

% Appendix A: Research Methodology >>>
\chapter{Appendix A}
\label{ap:research}
%\backgroundsetup{contents={DRAFT}}
In this appendix, I will explain in detail the process I undertook during the systematic review of literature during my thesis project, highlighting the search queries I used, the outcomes of those queries, and the sources from which I retrieved the papers.
\backgroundsetup{contents={}}
% <<< End of Appendix A

% Appendix B: Python Code >>>
\chapter{Appendix B}
\label{ap:Python}%
\backgroundsetup{contents={DRAFT}}
In this appendix, I want to suggest a potential implementation of the Bagging Voronoi clustering algorithm discussed in \ref{ch:baggingvoronoi}. I decided to implement the algorithm in Python due to its straightforward nature and the ability to execute code concurrently, leveraging on the parallel computing capacities that today's computers offer. If there is the need to implement the algorithm for real-time in-situ monitoring, it would be essential to rewrite the algorithm in C, as it offers significantly greater computational speed compared to Python. However, there isn't any existing library for functional data analysis already implemented in C, but in C++, which has comparable performances to C. The provided code represents an initial implementation of the function: it operates correctly, but it doesn't fully exploit the parallelization potential inherent in the bootstrapping approach. I chose to construct a class structure to allow for easy integration into a Python package. The development of a package for this analysis in Python could be a subsequent step. Moreover, concurrent Python programming can be seamlessly integrated into cloud environments (such as Google Cloud Platform, Amazon AWS, or Microsoft Azure), leveraging also containerization. With proper economic consideration and a well-thought-out business plan, this could potentially be marketed as a Software as a Service (SaaS) to companies that use PBF to print parts using the technologies detailed in this thesis. To perform clustering I used k-means clustering, to compute the distance i used the Euclidean distance and as a clustering matching method the permutation of contincengy matrix explained in Section \ref{sec:bvc}.
%\lstinputlisting[language=Python]{analysis-to-latex.py}
\backgroundsetup{contents={}}
% <<< End of Appendix B

\blankpage
\cleardoublepage
\end{document}
